% Sensibilisation sur la permanence du renouvellement moléculaire des cellules — SVT 1ère D — Cours signé OnBuch+
% Programme officiel MINESEC. Compiler avec : tectonic ND02-renouvellement-moleculaire-cellules.tex
\newcommand{\DOCMATIERE}{SVT}
\newcommand{\DOCNIVEAU}{1ère D}
\newcommand{\DOCMODULE}{Module 1 : Le Monde Vivant}
\newcommand{\DOCLECON}{ND02}
\newcommand{\DOCTITRE}{Sensibilisation sur la permanence du renouvellement moléculaire des cellules}
% OnBuch+ — préambule « cours signés » ENRICHI (compilé avec tectonic / XeLaTeX)
% Système visuel « Pop » d'OnBuch+ : fond crème (jamais blanc), bordures encre
% épaisses, ombres « dures » décalées, titres Archivo Black en capitales.
% Version MATHÉMATIQUES : graphes (pgfplots/tikz), boîtes pédagogiques
% (définition, propriété, méthode, attention, le savais-tu, exercice résolu,
% exercice, TP, bilan), page de garde et signature OnBuch+ ancrée en bas.
%
% Macros attendues dans main.tex AVANT \input{preamble} :
%   \DOCMATIERE  \DOCNIVEAU  \DOCMODULE  \DOCLECON (numéro)  \DOCTITRE
\documentclass[11pt]{article}

\usepackage{fontspec}
\usepackage[french]{babel}
\usepackage[a4paper,margin=2cm,top=3.4cm,bottom=2.8cm,headheight=1.4cm,footskip=1.3cm]{geometry}
\usepackage{amsmath,amssymb,amsfonts}
\usepackage{mathtools} % \xmapsto, \coloneqq... souvent utilisés par les modèles
\usepackage{cancel} % simplifications barrées (bilans de forces)
% \abs{z} (module, valeur absolue) et \norm{u} (norme), utilisés par les modèles
\providecommand{\abs}{}\let\abs\relax\DeclarePairedDelimiter{\abs}{\lvert}{\rvert}
\providecommand{\norm}{}\let\norm\relax\DeclarePairedDelimiter{\norm}{\lVert}{\rVert}
\usepackage[table]{xcolor} % [table] : fournit aussi \rowcolors (alternance de lignes)
\usepackage{pagecolor}
\usepackage{tikz}
\usetikzlibrary{arrows.meta,positioning,calc,shapes.geometric,shapes.misc,shapes.arrows,decorations.pathmorphing,decorations.pathreplacing,decorations.markings,patterns,shadows,mindmap,trees,fit,backgrounds,matrix,intersections,angles,quotes}
\usepackage{pgfplots}
\usepackage[version=4]{mhchem} % équations nucléaires \ce{^{235}_{92}U}
\usepackage[european,siunitx]{circuitikz} % schémas électriques
\pgfplotsset{compat=1.18}
\usepgfplotslibrary{fillbetween,groupplots} % aires entre courbes (calcul intégral)
\usepackage{enumitem}
\usepackage{fancyhdr}
% [most] charge toutes les bibliothèques tcolorbox (listings, minted, raster,
% documentation...) et gonfle inutilement la mémoire TeX sur un document long
% (~300 boîtes) : on ne charge que ce qui est réellement utilisé.
\usepackage[skins,breakable]{tcolorbox}
\usepackage{graphicx}
\usepackage{colortbl}
\usepackage{booktabs}
\usepackage{tabularx}
\usepackage{diagbox} % case d'en-tête diagonale des tableaux à double entrée
% Colonnes extensibles centrée / gauche / droite (C, L, R), souvent utilisées par les modèles
\newcolumntype{C}{>{\centering\arraybackslash}X}
\newcolumntype{L}{>{\raggedright\arraybackslash}X}
\newcolumntype{R}{>{\raggedleft\arraybackslash}X}
\usepackage{array}
\usepackage{multicol}
\usepackage{siunitx}
% parse-numbers=false : accepte aussi \SI{2,5 \times 10^{-3}}{...}
\sisetup{locale=FR,output-decimal-marker={,},group-minimum-digits=4,parse-numbers=false}
\DeclareSIUnit\ohm{\ensuremath{\symup{\Omega}}} % U+2126 absent de la police
\DeclareSIUnit\division{div} % réglages d'oscilloscope (ms/div, V/div)
\DeclareSIUnit\tr{tr} % tours (vitesse de rotation, ex. \SI{1500}{\tr\per\minute})
\DeclareSIUnit\ch{ch} % cheval-vapeur (puissance moteur, unité usuelle hors SI)
\DeclareSIUnit\franccfa{FCFA} % franc CFA (coûts, contexte camerounais)
\DeclareSIUnit\dioptre{\ensuremath{\delta}} % vergence d'une lentille (dioptrie)
\usepackage{qrcode}
% \degree : commande fréquemment utilisée par les modèles pour un angle ou une
% température en dehors de \SI{}{} (non fournie par les paquets chargés).
\providecommand{\degree}{^{\circ}}
% \qed (fin de démonstration), utilisé par les modèles sans amsthm
\providecommand{\qed}{\hfill$\square$}
% \barre : double barre des valeurs interdites dans un tableau de variations (tkz-tab)
\providecommand{\barre}{\ensuremath{\|}}
% environnement proof (amsthm non chargé) utilisé par les modèles
\ifdefined\proof\else\newenvironment{proof}[1][Démonstration]{\par\noindent\textit{#1.}\ }{\qed\par}\fi
\usepackage{lastpage}
\usepackage{hyperref}
\hypersetup{hidelinks}

% ============================================================
% Palette « Pop » OnBuch+ — mêmes hex que la classe P (plus_ui.dart)
% ============================================================
\definecolor{poporange}{HTML}{F59321}
\definecolor{popdark}{HTML}{E07A0C}
\definecolor{popcream}{HTML}{FFF6E8}
\definecolor{popink}{HTML}{1C1714}
\definecolor{poppurple}{HTML}{7A5AE0}
\definecolor{popgreen}{HTML}{2E9E6A}
\definecolor{poppink}{HTML}{E0457B}
\definecolor{popblue}{HTML}{2F7DE1}
\definecolor{popgold}{HTML}{C98A12}
\definecolor{popmuted}{HTML}{8A7F76}
\definecolor{popline}{HTML}{EADFD2}
\definecolor{popgray}{HTML}{8A7F76}
\colorlet{popgrey}{popgray}
% Alias tolérants (le modèle nomme parfois les couleurs différemment)
\colorlet{popwhite}{white}
\colorlet{popblack}{popink}
\colorlet{popred}{poppink}
\colorlet{popyellow}{popgold}
% Teintes claires pour fonds de tableaux / zones
\colorlet{popblueL}{popblue!10!white}
\colorlet{poporangeL}{poporange!14!white}
\colorlet{popgreenL}{popgreen!12!white}
\colorlet{poppurpleL}{poppurple!12!white}
\colorlet{poppinkL}{poppink!12!white}
\colorlet{popgoldL}{popgold!14!white}

\pagecolor{popcream}

% ============================================================
% Typographie
% ============================================================
\defaultfontfeatures{Ligatures=TeX}
\setmainfont{PlusJakartaSans-Medium.ttf}[Path=../../../fonts/,BoldFont=PlusJakartaSans-Bold.ttf]
% Maths en sans-serif (Fira Math) avec chiffres et lettres droites de Plus
% Jakarta, pour que les formules se fondent dans le texte.
\usepackage[math-style=ISO,bold-style=ISO]{unicode-math}
\setmathfont{FiraMath-Regular.otf}[Path=../../../fonts/]
\setmathfont{PlusJakartaSans-Medium.ttf}[Path=../../../fonts/,range={up/num,up/{Latin,latin},\mathup}]
\usepackage{icomma} % virgule décimale sans espace en mode math
% Caractères mathématiques Unicode absents des polices de texte (Plus Jakarta,
% Archivo Black) : rendus par leur équivalent mathématique, en texte comme en
% formule — sinon ils s'affichent en carré vide (ex. « Arithmétique dans ℤ »).
\usepackage{newunicodechar}
\newunicodechar{ℝ}{\ensuremath{\mathbb{R}}}
\newunicodechar{ℤ}{\ensuremath{\mathbb{Z}}}
\newunicodechar{ℂ}{\ensuremath{\mathbb{C}}}
\newunicodechar{ℕ}{\ensuremath{\mathbb{N}}}
\newunicodechar{ℚ}{\ensuremath{\mathbb{Q}}}
\newunicodechar{𝔻}{\ensuremath{\mathbb{D}}}
\newunicodechar{α}{\ensuremath{\alpha}}
\newunicodechar{β}{\ensuremath{\beta}}
\newunicodechar{γ}{\ensuremath{\gamma}}
\newunicodechar{δ}{\ensuremath{\delta}}
\newunicodechar{ε}{\ensuremath{\varepsilon}}
\newunicodechar{θ}{\ensuremath{\theta}}
\newunicodechar{λ}{\ensuremath{\lambda}}
\newunicodechar{μ}{\ensuremath{\mu}}
\newunicodechar{σ}{\ensuremath{\sigma}}
\newunicodechar{τ}{\ensuremath{\tau}}
\newunicodechar{φ}{\ensuremath{\varphi}}
\newunicodechar{ψ}{\ensuremath{\psi}}
\newunicodechar{ω}{\ensuremath{\omega}}
\newunicodechar{Σ}{\ensuremath{\Sigma}}
% majuscules produites par \MakeUppercase dans les titres (α-aminés -> Α)
\newunicodechar{Α}{\ensuremath{\alpha}}
\newunicodechar{Β}{\ensuremath{\beta}}
\newunicodechar{Γ}{\ensuremath{\Gamma}}
\newunicodechar{Δ}{\ensuremath{\Delta}}
\newunicodechar{Ε}{\ensuremath{\varepsilon}}
\newunicodechar{Θ}{\ensuremath{\Theta}}
\newunicodechar{Λ}{\ensuremath{\Lambda}}
\newunicodechar{Μ}{\ensuremath{\mu}}
\newunicodechar{Τ}{\ensuremath{\tau}}
\newunicodechar{Φ}{\ensuremath{\Phi}}
\newunicodechar{Ψ}{\ensuremath{\Psi}}
\newunicodechar{Ω}{\ensuremath{\Omega}}
\newunicodechar{↦}{\ensuremath{\mapsto}}
\newunicodechar{∈}{\ensuremath{\in}}
\newunicodechar{∉}{\ensuremath{\notin}}
\newunicodechar{∧}{\ensuremath{\wedge}}
\newunicodechar{⇔}{\ensuremath{\Leftrightarrow}}
\newunicodechar{⟺}{\ensuremath{\Longleftrightarrow}}
\newunicodechar{⇒}{\ensuremath{\Rightarrow}}
\newunicodechar{⟹}{\ensuremath{\Longrightarrow}}
\newunicodechar{∘}{\ensuremath{\circ}}
\newunicodechar{∪}{\ensuremath{\cup}}
\newunicodechar{∩}{\ensuremath{\cap}}
\newunicodechar{≡}{\ensuremath{\equiv}}
\newunicodechar{⊂}{\ensuremath{\subset}}
\newunicodechar{⊥}{\ensuremath{\perp}}
\newunicodechar{∃}{\ensuremath{\exists}}
\newunicodechar{∀}{\ensuremath{\forall}}
\newunicodechar{∛}{\ensuremath{\sqrt[3]{\,}}}
\newunicodechar{∜}{\ensuremath{\sqrt[4]{\,}}}
\newunicodechar{⃗}{\ensuremath{{}^{\to}}}
\newunicodechar{ⁿ}{\ifmmode^{n}\else\textsuperscript{\ensuremath{n}}\fi}
\newunicodechar{ˣ}{\ifmmode^{x}\else\textsuperscript{\ensuremath{x}}\fi}
\newunicodechar{ᵇ}{\ifmmode^{b}\else\textsuperscript{\ensuremath{b}}\fi}
\newunicodechar{ʳ}{\ifmmode^{r}\else\textsuperscript{\ensuremath{r}}\fi}
\newunicodechar{ᵘ}{\ifmmode^{u}\else\textsuperscript{\ensuremath{u}}\fi}
\newunicodechar{ʸ}{\ifmmode^{y}\else\textsuperscript{\ensuremath{y}}\fi}
\newunicodechar{ᵃ}{\ifmmode^{a}\else\textsuperscript{\ensuremath{a}}\fi}
\newunicodechar{ᵉ}{\ifmmode^{e}\else\textsuperscript{\ensuremath{e}}\fi}
\newunicodechar{ᵏ}{\ifmmode^{k}\else\textsuperscript{\ensuremath{k}}\fi}
\newunicodechar{ᵖ}{\ifmmode^{p}\else\textsuperscript{\ensuremath{p}}\fi}
\newunicodechar{ᴺ}{\ifmmode^{N}\else\textsuperscript{\ensuremath{N}}\fi}
\newunicodechar{ᶜ}{\ifmmode^{c}\else\textsuperscript{\ensuremath{c}}\fi}
\newunicodechar{ʰ}{\ifmmode^{h}\else\textsuperscript{\ensuremath{h}}\fi}
\newunicodechar{ᵠ}{\ifmmode^{q}\else\textsuperscript{\ensuremath{q}}\fi}
\newunicodechar{ₙ}{\ifmmode_{n}\else\textsubscript{\ensuremath{n}}\fi}
\newunicodechar{ₐ}{\ifmmode_{a}\else\textsubscript{\ensuremath{a}}\fi}
\newunicodechar{ᵢ}{\ifmmode_{i}\else\textsubscript{\ensuremath{i}}\fi}
\newunicodechar{ᵣ}{\ifmmode_{r}\else\textsubscript{\ensuremath{r}}\fi}
\newunicodechar{ₖ}{\ifmmode_{k}\else\textsubscript{\ensuremath{k}}\fi}
\newunicodechar{ᵦ}{\ifmmode_{\beta}\else\textsubscript{\ensuremath{\beta}}\fi}
\newunicodechar{ₓ}{\ifmmode_{x}\else\textsubscript{\ensuremath{x}}\fi}
\newfontfamily\popdisplay{ArchivoBlack-Regular.ttf}[Path=../../../fonts/]
\newfontfamily\poplogo{SpaceGrotesk-Bold.ttf}[Path=../../../fonts/]
\newfontfamily\popsemi{PlusJakartaSans-SemiBold.ttf}[Path=../../../fonts/]
\newfontfamily\popxbold{PlusJakartaSans-ExtraBold.ttf}[Path=../../../fonts/]
\newfontfamily\popxboldls{PlusJakartaSans-ExtraBold.ttf}[Path=../../../fonts/,LetterSpace=3.0]

\setlist[enumerate]{leftmargin=1.7em,itemsep=5pt,topsep=4pt}
\setlist[itemize]{leftmargin=1.7em,itemsep=4pt,topsep=4pt,label={\color{poporange}\textbullet}}
\setlength{\parindent}{0pt}
\setlength{\parskip}{5pt}
\renewcommand{\arraystretch}{1.25}

% pgfplots : style Pop par défaut
\pgfplotsset{
  every axis/.append style={axis line style={popink,line width=1pt},tick style={popink},
    grid style={popline,line width=0.5pt},label style={font=\small},tick label style={font=\footnotesize}},
  popcurve/.style={poporange,line width=1.8pt,no markers,smooth},
}
% Même style disponible dans un \draw TikZ simple (hors pgfplots)
\tikzset{popcurve/.style={poporange,line width=1.8pt}}
\tikzset{popgrid/.style={popline,very thin}}
\colorlet{popcurve}{poporange} % le nom du style est parfois utilisé comme couleur

% ============================================================
% Pill / squiggle
% ============================================================
\newcommand{\poppill}[3]{%
  \tikz[baseline=-0.55ex]\node[draw=popink,line width=1.2pt,rounded corners=2.6pt,%
    fill=#2,inner xsep=6.5pt,inner ysep=3pt]{%
    \popxboldls\fontsize{7}{7}\selectfont\color{#3}\MakeUppercase{#1}};%
}
\newcommand{\popsquiggle}[1]{%
  \tikz[baseline=-0.4ex]{
    \draw[#1,line width=2.6pt,line cap=round]
      plot [smooth] coordinates {(0,0.09) (0.34,0.01) (0.68,0.21) (1.02,0.01) (1.36,0.10)};
  }%
}
% Mot-clé surligné (marqueur orange sous le texte)
\newcommand{\cle}[1]{{\popxbold\color{popdark}#1}}

% ============================================================
% En-tête / pied
% ============================================================
\pagestyle{fancy}
\fancyhf{}
\renewcommand{\headrulewidth}{0pt}
\renewcommand{\footrulewidth}{0pt}
\lhead{\raisebox{-0.15ex}{\poplogo\fontsize{13}{13}\selectfont\color{popink}OnBuch\color{poporange}+}}
\rhead{\poppill{\DOCMATIERE\ \textbullet\ \DOCNIVEAU\ \textbullet\ Leçon \DOCLECON}{white}{popink}}
\lfoot{\popxbold\fontsize{7.6}{7.6}\selectfont\color{popmuted}\MakeUppercase{Cours signé OnBuch+}\ \textbullet\ {\popsemi\DOCTITRE}}
\rfoot{\tikz[baseline=-0.6ex]\node[rounded corners=8.5pt,fill=popink,minimum height=17pt,minimum width=17pt,inner xsep=5pt,inner sep=0]{\popxbold\fontsize{8}{8}\selectfont\color{popcream}\thepage\,/\,\pageref*{LastPage}};}

% ============================================================
% Titres
% ============================================================
\newcommand{\chaptitre}[2]{%
  \begin{tcolorbox}[enhanced,colback=poporange,colframe=popink,boxrule=2.6pt,arc=11pt,
      drop shadow=popink,left=15pt,right=15pt,top=12pt,bottom=11pt,before skip=3pt,after skip=18pt]
    \poppill{#2}{popink}{popcream}\par\vspace{9pt}
    {\raggedright\hyphenpenalty=10000\exhyphenpenalty=10000\popdisplay\fontsize{22}{24}\selectfont\color{popcream}\MakeUppercase{#1}\par}
  \end{tcolorbox}
}
% Section numérotée : pastille orange avec le numéro + titre Archivo
\newcounter{popsec}
\newcommand{\coursec}[1]{%
  \stepcounter{popsec}\setcounter{popsub}{0}%
  \par\vspace{16pt}\noindent
  \begin{minipage}{\linewidth}
  \tikz[baseline=-0.7ex]\node[draw=popink,line width=1.6pt,fill=poporange,rounded corners=4pt,minimum size=22pt,inner sep=0]{\popdisplay\fontsize{12}{12}\selectfont\color{popcream}\thepopsec};%
  \hspace{8pt}{\popdisplay\fontsize{15}{17}\selectfont\color{popink}\MakeUppercase{#1}}\par
  \vspace{4pt}\noindent{\color{poporange}\rule{36pt}{3.6pt}}
  \end{minipage}\par\nopagebreak\vspace{8pt}
}
\newcounter{popsub}
\newcommand{\courssub}[1]{%
  \stepcounter{popsub}%
  \par\vspace{9pt}\noindent{\popxbold\fontsize{11.5}{13}\selectfont\color{popink}{\color{poporange}\thepopsec.\thepopsub}\hspace{6pt}#1}\par\nopagebreak\vspace{4pt}
}

% ============================================================
% Boîtes pédagogiques — même recette : fond blanc, bordure épaisse colorée,
% ombre dure encre, badge plein. Titre optionnel [#1].
% ============================================================
\tcbset{popbase/.style={breakable,enhanced,colback=white,boxrule=2.3pt,arc=9pt,
  shadow={3pt}{-3pt}{0pt}{popink},left=14pt,right=14pt,top=11pt,bottom=11pt,before skip=11pt,after skip=15pt,
  fontupper=\popsemi\fontsize{10.6}{14.5}\selectfont\color{popink},
  fontlower=\popsemi\fontsize{10.6}{14.5}\selectfont\color{popink}}}
% Coche dessinée (le glyphe ✓ n'existe pas dans les polices du document)
\newcommand{\popcheck}[1]{\tikz[baseline=-0.5ex]\draw[#1,line width=1.6pt,line cap=round,line join=round](-0.13,0.0)--(-0.04,-0.09)--(0.13,0.1);}
\providecommand{\coche}{\popcheck{popgreen}}
\providecommand{\resultat}[1]{\par\smallskip\noindent\fcolorbox{popgreen}{popgreenL}{\parbox{\dimexpr\linewidth-2\fboxsep-2\fboxrule\relax}{\textbf{Résultat.} #1}}\par\smallskip}
\newcommand{\popboxhead}[3]{\poppill{#1}{#2}{white}\ifx\relax#3\relax\else\hspace{6pt}{\popxbold\fontsize{10.6}{12}\selectfont\color{popink}#3}\fi\par\vspace{7pt}}

\newtcolorbox{aretenir}[1][]{popbase,colframe=popgreen,before upper={\popboxhead{À retenir}{popgreen}{#1}}}
\newtcolorbox{exemplebox}[1][]{popbase,colframe=poporange,before upper={\popboxhead{Exemple}{poporange}{#1}}}
\newtcolorbox{definition}[1][]{popbase,colframe=popblue,before upper={\popboxhead{Définition}{popblue}{#1}}}
\newtcolorbox{propriete}[1][]{popbase,colframe=poppurple,before upper={\popboxhead{Propriété}{poppurple}{#1}}}
\newtcolorbox{methode}[1][]{popbase,colframe=popgold,before upper={\popboxhead{Méthode}{popgold}{#1}}}
\newtcolorbox{attention}[1][]{popbase,colframe=poppink,before upper={\popboxhead{Attention}{poppink}{#1}}}
\newtcolorbox{experience}[1][]{popbase,colframe=popblue,colback=popblueL,before upper={\popboxhead{Expérience}{popblue}{#1}}}
% « Le savais-tu ? » : carte violette pleine (contexte, culture, Cameroun)
\newtcolorbox{savaistu}[1][]{popbase,colback=poppurple,colframe=popink,
  fontupper=\popsemi\fontsize{10.4}{14.2}\selectfont\color{white},
  before upper={\poppill{Le savais-tu\,?}{white}{poppurple}\ifx\relax#1\relax\else\hspace{6pt}{\popxbold\color{white}#1}\fi\par\vspace{7pt}}}
% Exercice résolu : énoncé en haut, \tcblower puis la solution sur fond crème
\newcounter{popexo}\newcounter{popexr}
\newtcolorbox{exoresolu}[1][]{popbase,colframe=poporange,colbacklower=popcream,
  segmentation style={popink,line width=1.2pt,dashed},
  before upper={\stepcounter{popexr}\popboxhead{Exercice résolu \thepopexr}{poporange}{#1}},
  before lower={\poppill{Solution}{popink}{popcream}\par\vspace{6pt}}}
% Exercice (non résolu en place) — la correction est dans la partie Corrigés
\newtcolorbox{exercice}[1][]{popbase,colframe=popink,
  before upper={\stepcounter{popexo}\popboxhead{Exercice \thepopexo}{popink}{#1}}}
% Corrigé
\newtcolorbox{corrige}[1][]{popbase,colframe=popgreen,colback=popgreenL,
  before upper={\popboxhead{Corrigé}{popgreen}{#1}}}
% Objectifs / prérequis / bilan
\newtcolorbox{objectifs}{popbase,colframe=popink,colback=poporangeL,
  before upper={\popboxhead{Objectifs de la leçon}{popink}{}}}
\newtcolorbox{prerequis}{popbase,colframe=popink,colback=popblueL,
  before upper={\popboxhead{Prérequis}{popblue}{}}}
\newtcolorbox{bilan}{popbase,colframe=popink,colback=white,boxrule=2.8pt,
  before upper={\popboxhead{Fiche bilan}{poporange}{L'essentiel à connaître pour l'examen}}}
% Figure encadrée avec légende
\newtcolorbox{popfigure}[1][]{enhanced,breakable=false,colback=white,colframe=popline,boxrule=1.4pt,arc=8pt,
  left=10pt,right=10pt,top=10pt,bottom=8pt,before skip=10pt,after skip=12pt,halign=center,
  lower separated=false,halign lower=center,
  fontlower=\popsemi\fontsize{9}{11.5}\selectfont\color{popmuted},
  #1}
\newcommand{\legende}[1]{\par\vspace{6pt}{\popsemi\fontsize{9}{11.5}\selectfont\color{popmuted}#1}}

% ============================================================
% Signature OnBuch+
% ============================================================
\newcommand{\poplogoinline}[1]{{\poplogo\fontsize{#1}{#1}\selectfont\color{popink}OnBuch\color{poporange}+}}
\newcommand{\signatureonbuch}{%
  \begin{tcolorbox}[enhanced,colback=popink,colframe=popink,boxrule=2.2pt,arc=10pt,
      drop shadow=poporange,left=14pt,right=14pt,top=10pt,bottom=10pt,before skip=0pt,after skip=0pt]
    \tikz[baseline=-0.7ex]\node[circle,fill=popgreen,draw=popcream,line width=1.3pt,minimum size=22pt,inner sep=0]{\popcheck{white}};%
    \hspace{9pt}\begin{minipage}[c]{0.62\linewidth}
      {\popxbold\fontsize{12}{13}\selectfont\color{popcream}Signé\ \textbullet\ {\poplogo OnBuch\color{poporange}+}}\par\vspace{2pt}
      {\raggedright\popsemi\fontsize{8}{10.5}\selectfont\color{popline}\MakeUppercase{Équipe pédagogique OnBuch+}\\\MakeUppercase{Conforme au programme officiel MINESEC}\par}
    \end{minipage}\hfill
    {\poplogo\fontsize{22}{22}\selectfont\color{popcream}OnBuch\color{poporange}+}
  \end{tcolorbox}%
}

% ============================================================
% Page de garde
% #1 titre · #2 matière · #3 niveau · #4 module · #5 n° leçon · #6 durée ·
% #7 description / objectifs (texte ou itemize)
% ============================================================
\newcommand{\pagegarde}[7]{%
  \thispagestyle{empty}
  \newgeometry{a4paper,margin=1.7cm,top=1.6cm,bottom=1.4cm}%
  \begin{tikzpicture}[remember picture,overlay]
    \fill[poporange!18] ([xshift=-3.2cm,yshift=-3.6cm]current page.north east) circle (4.6cm);
    \fill[poppurple!14] ([xshift=2.4cm,yshift=7.6cm]current page.south west) circle (3.8cm);
  \end{tikzpicture}%
  {\poplogo\fontsize{30}{30}\selectfont\color{popink}OnBuch\color{poporange}+}\hfill
  \raisebox{0.6ex}{\poppill{Cours signé \textbullet\ Premium}{popink}{popcream}}\par
  \vspace{1pt}\popsquiggle{poppurple}\par
  \vspace{8pt}
  \begin{tcolorbox}[enhanced,colback=poporange,colframe=popink,boxrule=2.8pt,arc=13pt,
      drop shadow=popink,left=18pt,right=18pt,top=12pt,bottom=13pt,after skip=10pt]
    \poppill{#2 \textbullet\ #3}{popink}{popcream}\hspace{5pt}\poppill{#4}{white}{popink}\par\vspace{8pt}
    {\popdisplay\fontsize{14}{14}\selectfont\color{popink}LEÇON #5}\par\vspace{4pt}
    {\raggedright\hyphenpenalty=10000\exhyphenpenalty=10000\popdisplay\fontsize{25}{27}\selectfont\color{popcream}\MakeUppercase{#1}\par}
    \vspace{8pt}
    {\popxbold\fontsize{9.5}{9.5}\selectfont\color{popink}\MakeUppercase{Durée conseillée : #6}}
  \end{tcolorbox}
  \begin{tcolorbox}[enhanced,colback=white,colframe=popink,boxrule=2.2pt,arc=10pt,
      drop shadow=popink,left=16pt,right=16pt,top=10pt,bottom=10pt,after skip=8pt]
    \poppill{Dans cette leçon}{poppurple}{white}\par\vspace{5pt}
    \popsemi\fontsize{9.3}{12.6}\selectfont\color{popink}#7
  \end{tcolorbox}
  \noindent\begin{minipage}[t]{0.487\linewidth}
    \begin{tcolorbox}[enhanced,colback=popgreen,colframe=popink,boxrule=2.2pt,arc=10pt,
        drop shadow=popink,left=11pt,right=11pt,top=10pt,bottom=10pt,height=3.1cm,valign=center]
      \tcbox[colback=white,colframe=popink,boxrule=1.3pt,arc=5pt,boxsep=0pt,nobeforeafter,
        left=3pt,right=3pt,top=3pt,bottom=3pt,valign=center]{\qrcode[height=1.9cm]{https://play.google.com/store/apps/details?id=com.luvvix.onbuch&hl=fr}}%
      \hspace{8pt}\begin{minipage}[c]{0.5\linewidth}
        {\popdisplay\fontsize{11}{12.5}\selectfont\color{white}\MakeUppercase{Télécharge OnBuch}}\par\vspace{4pt}
        {\raggedright\popsemi\fontsize{8.4}{11}\selectfont\color{white}Gratuit \textbullet\ Google Play\\Tuteur IA, annales, concours, résultats\par}
      \end{minipage}
    \end{tcolorbox}
  \end{minipage}\hfill
  \begin{minipage}[t]{0.487\linewidth}
    \begin{tcolorbox}[enhanced,colback=poppurple,colframe=popink,boxrule=2.2pt,arc=10pt,
        drop shadow=popink,left=11pt,right=11pt,top=10pt,bottom=10pt,height=3.1cm,valign=center]
      \tikz[baseline=-0.7ex]\node[circle,fill=white,draw=popink,line width=1.3pt,minimum size=1.6cm,inner sep=0]{\popdisplay\fontsize{24}{24}\selectfont\color{poppurple}+};%
      \hspace{8pt}\begin{minipage}[c]{0.6\linewidth}
        {\popdisplay\fontsize{11}{12.5}\selectfont\color{white}\MakeUppercase{Passe à OnBuch+}}\par\vspace{4pt}
        {\raggedright\popsemi\fontsize{8.4}{11}\selectfont\color{white}Tous les cours signés, Tuteur IA illimité, fiches PDF — onglet OnBuch+ de l'app\par}
      \end{minipage}
    \end{tcolorbox}
  \end{minipage}
  \vspace{8pt}
  \popcontenu
  \vfill
  \signatureonbuch
  \restoregeometry
  \newpage
}

% Rangée « Ce cours contient » (page de garde)
\newcommand{\poptile}[3]{% #1 couleur #2 gros texte #3 légende
  \begin{tcolorbox}[enhanced,colback=#1,colframe=popink,boxrule=2pt,arc=9pt,shadow={2.5pt}{-2.5pt}{0pt}{popink},
      left=6pt,right=6pt,top=9pt,bottom=9pt,halign=center,valign=center,height=1.75cm,width=\linewidth,nobeforeafter]
    {\popdisplay\fontsize{12.5}{13}\selectfont\color{white}\MakeUppercase{#2}}\par\vspace{3pt}
    {\centering\popsemi\fontsize{7.8}{9.5}\selectfont\color{white}#3\par}
  \end{tcolorbox}}
\newcommand{\popcontenu}{%
  \par\noindent{\popxboldls\fontsize{7.5}{7.5}\selectfont\color{popmuted}CE COURS SIGNÉ CONTIENT}\par\vspace{7pt}
  \noindent\begin{minipage}[t]{0.235\linewidth}\poptile{popblue}{Cours}{complet, illustré, pas à pas}\end{minipage}\hfill
  \begin{minipage}[t]{0.235\linewidth}\poptile{poporange}{Méthodes}{et exercices résolus}\end{minipage}\hfill
  \begin{minipage}[t]{0.235\linewidth}\poptile{poppink}{Exercices}{type examen + corrigés}\end{minipage}\hfill
  \begin{minipage}[t]{0.235\linewidth}\poptile{popgreen}{Fiche bilan}{l'essentiel à retenir}\end{minipage}\par}

% Sommaire
\newtcolorbox{sommaire}{enhanced,colback=white,colframe=popink,boxrule=2.2pt,arc=10pt,
  drop shadow=popink,left=18pt,right=18pt,top=13pt,bottom=13pt,before skip=6pt,after skip=20pt,
  before upper={\poppill{Sommaire}{poppurple}{white}\par\vspace{9pt}\popsemi\fontsize{10.6}{14.5}\selectfont\color{popink}}}

% Signature de fin de cours — poussée tout en bas de la dernière page
\newcommand{\findecours}{%
  \par\vfill\nopagebreak
  \begin{center}\popsquiggle{poporange}\end{center}\vspace{4pt}
  \signatureonbuch
}
\AtBeginDocument{\def\llbracket{\mathopen{[\mkern-3.5mu[}}\def\rrbracket{\mathclose{]\mkern-3.5mu]}}} % intervalles d'entiers (glyphe ⟦ absent de la police)
% Glyphes absents de Fira Math : redéfinis à partir de glyphes disponibles
\AtBeginDocument{%
  \def\setminus{\mathbin{\backslash}}\let\smallsetminus\setminus
  \def\vdots{\mathord{\vbox{\baselineskip3.2pt\lineskiplimit0pt\kern4pt\hbox{.}\hbox{.}\hbox{.}}}}%
  \def\ddots{\mathinner{\mkern1mu\raise6.5pt\hbox{.}\mkern2mu\raise3.6pt\hbox{.}\mkern2mu\raise0.7pt\hbox{.}\mkern1mu}}%
  \def\triangle{\mathord{\scalebox{0.9}{$\symup{\Delta}$}}}%
  \def\nabla{\mathord{\raisebox{\depth}{\scalebox{1}[-1]{$\symup{\Delta}$}}}}%
  \def\checkmark{\popcheck{popdark}}%
  \def\star{\mathbin{\ast}}\def\bigstar{\mathord{\ast}}%
  \def\diamond{\mathbin{\scalebox{0.6}{$\Diamond$}}}%
  \def\bigcup{\mathop{\mathchoice{\raisebox{-0.3em}{\scalebox{1.7}{$\cup$}}}{\scalebox{1.25}{$\cup$}}{\cup}{\cup}}\displaylimits}%
  \def\bigcap{\mathop{\mathchoice{\raisebox{-0.3em}{\scalebox{1.7}{$\cap$}}}{\scalebox{1.25}{$\cap$}}{\cap}{\cap}}\displaylimits}%
  \def\lesseqqgtr{\mathrel{\lessgtr}}\def\bigoplus{\mathop{\oplus}\displaylimits}%
}
\providecommand{\ssi}{si et seulement si} % abréviation fréquente
\AtBeginDocument{\providecommand{\wideparen}{\overparen}} % arc de cercle

%% FIGFIX-BEGIN v5 — correctifs globaux des figures (docs/onbuch-generation/tools/figfix)
\usepackage{adjustbox}\usepackage{etoolbox}\tcbuselibrary{raster}
% toute figure TikZ/pgfplots plus large que la ligne est réduite à la largeur disponible
\BeforeBeginEnvironment{tikzpicture}{\begin{adjustbox}{max width=\linewidth}}
\AfterEndEnvironment{tikzpicture}{\end{adjustbox}}
% légendes pgfplots lisibles par-dessus les courbes
\pgfplotsset{every axis/.append style={legend style={fill=white,fill opacity=0.92,text opacity=1,draw=black!15,inner sep=3pt}}}
% étiquettes d'arêtes (syntaxe edge["..."]) avec fond blanc
\tikzset{every edge quotes/.append style={fill=white,inner sep=1.2pt}}
%% FIGFIX-END
\begin{document}

\pagegarde{Sensibilisation sur la permanence du renouvellement moléculaire des cellules}{SVT}{1ère D}{Module 1 : Le Monde Vivant}{ND02}{6 h}{Cette leçon explore le renouvellement permanent des molécules constitutives de la cellule (protéines, lipides, ADN) et ses mécanismes de synthèse et de dégradation. Elle montre comment le ralentissement de ce turnover conduit au vieillissement cellulaire et pourquoi ce processus est essentiel à la santé de l'organisme.
\vspace{4pt}
\begin{multicols}{2}\small
\begin{itemize}
\item À la fin de cette leçon, je sais définir le renouvellement moléculaire (turnover) et citer les grandes familles de molécules concernées
\item À la fin de cette leçon, je sais décrire les mécanismes cellulaires de synthèse et de dégradation des protéines, des lipides et de l'ADN
\item À la fin de cette leçon, je sais exploiter des données de marquage isotopique pour mettre en évidence le renouvellement moléculaire
\item À la fin de cette leçon, je sais calculer et interpréter une demi-vie moléculaire à partir d'un graphique de décroissance d'un traceur
\item À la fin de cette leçon, je sais relier le renouvellement moléculaire au maintien de l'intégrité et du bon fonctionnement de la cellule
\item À la fin de cette leçon, je sais expliquer les conséquences cellulaires du ralentissement du renouvellement moléculaire (vieillissement)
\end{itemize}\end{multicols}}

\begin{sommaire}
\begin{multicols}{2}
\begin{enumerate}
\item Le renouvellement moléculaire : une évidence expérimentale
\item Les mécanismes de synthèse moléculaire
\item Les mécanismes de dégradation moléculaire
\item Exploitation de données de marquage isotopique
\item Renouvellement moléculaire et intégrité cellulaire
\item Ralentissement du renouvellement et vieillissement cellulaire
\item Activité d'intégration
\item Méthodes et exercices résolus
\item Exercices
\item Corrigés des exercices
\item Fiche bilan
\end{enumerate}
\end{multicols}
\end{sommaire}

\begin{objectifs}
\begin{itemize}
\item À la fin de cette leçon, je sais définir le renouvellement moléculaire (turnover) et citer les grandes familles de molécules concernées
\item À la fin de cette leçon, je sais décrire les mécanismes cellulaires de synthèse et de dégradation des protéines, des lipides et de l'ADN
\item À la fin de cette leçon, je sais exploiter des données de marquage isotopique pour mettre en évidence le renouvellement moléculaire
\item À la fin de cette leçon, je sais calculer et interpréter une demi-vie moléculaire à partir d'un graphique de décroissance d'un traceur
\item À la fin de cette leçon, je sais relier le renouvellement moléculaire au maintien de l'intégrité et du bon fonctionnement de la cellule
\item À la fin de cette leçon, je sais expliquer les conséquences cellulaires du ralentissement du renouvellement moléculaire (vieillissement)
\item À la fin de cette leçon, je sais argumenter, à partir d'exemples concrets, sur l'importance du renouvellement moléculaire pour la santé de l'organisme
\end{itemize}
\end{objectifs}

\begin{prerequis}
\begin{itemize}
\item Constituants chimiques de la cellule : protéines, lipides, glucides, acides nucléiques (Seconde)
\item Structure et rôle de l'ADN, notion de gène et de synthèse des protéines (notions de base)
\item Organisation de la cellule eucaryote : noyau, cytoplasme, organites (ribosomes, réticulum, lysosomes)
\item Notion de métabolisme : anabolisme et catabolisme
\item Lecture et exploitation de graphiques et de tableaux de données expérimentales
\end{itemize}
\end{prerequis}

\newpage

\coursec{Le renouvellement moléculaire : une évidence expérimentale}

\courssub{Une situation d'accroche : Mama Odile et Junior}
Mama Odile, 65 ans, et son petit-fils Junior, 15 ans, se blessent au pied le même jour lors d'un match de quartier à Yaoundé. Tous deux mangent les mêmes plats riches en protéines — ndolé, arachides, poisson fumé — pourtant la plaie de Junior cicatrise en quelques jours alors que celle de Mama Odile met des semaines à refermer. \textbf{Comment expliquer cette différence de vitesse de réparation cellulaire ?} La réponse réside dans la capacité des cellules à renouveler sans cesse leurs molécules constitutives. Cette leçon lève le voile sur ce renouvellement permanent, véritable moteur de la vie cellulaire.

\courssub{Observation : une composition globale stable masquant un changement perpétuel}
Si l'on analyse la composition chimique d'une cellule à deux instants différents, les pourcentages en protéines, lipides, acides nucléiques et glucides restent quasiment identiques. Pourtant, au niveau moléculaire, chaque macromolécule est constamment dégradée et resynthétisée. C'est ce que l'on appelle un \cle{état stationnaire dynamique} : la cellule maintient sa composition globale stable grâce à un flux continu de molécules entrantes (synthèse) et sortantes (dégradation).

\begin{definition}[Renouvellement moléculaire (turnover)]
Ensemble des processus de synthèse et de dégradation qui renouvellent en permanence les molécules constitutives de la cellule (protéines, lipides membranaires, ADN, glucides de réserve) sans modifier sa composition chimique globale.
\end{definition}

\courssub{Les molécules concernées par le turnover}
\begin{itemize}
  \item \textbf{Protéines} : enzymes, protéines structurales (collagène, actine), protéines de signalisation, récepteurs membranaires.
  \item \textbf{Lipides membranaires} : phospholipides, cholestérol, glycolipides (renouvellement par remodelage des membranes).
  \item \textbf{ADN} : la séquence génétique est globalement conservée, mais des réparations locales (excision de bases, réparation par excision de nucléotides, recombinaison homologue) remplacent en permanence des segments endommagés ; lors de la division cellulaire, la réplication semi-conservative renouvelle les brins.
  \item \textbf{Glucides de réserve} : glycogène hépatique et musculaire, constamment mobilisé (glycogénolyse) et reconstitué (glycogenèse).
\end{itemize}

\begin{experience}[Expérience historique de Schoenheimer (1942) — Preuve du turnover protéique]
\textbf{Matériel :} Rats, leucine marquée à l'azote 15 (\ce{^{15}N}-leucine), spectromètre de masse.
\textbf{Protocole :}
\begin{enumerate}
  \item Injection intraveineuse de \ce{^{15}N}-leucine à des rats jeûnés.
  \item Prélèvement du foie à intervalles réguliers (0 h, 10 h, 24 h, 48 h, 72 h).
  \item Extraction des protéines hépatiques, hydrolyse acide, dosage de l'enrichissement en \ce{^{15}N} de la leucine protéique par spectrométrie de masse.
\end{enumerate}
\textbf{Observations :} L'enrichissement en \ce{^{15}N} de la leucine protéique augmente rapidement (pic vers 10 h), puis décroît exponentiellement pour revenir au niveau naturel (abondance isotopique \ce{^{15}N} \approx 0,37 \%) en quelques jours.
\textbf{Interprétation :} La leucine marquée est incorporée dans les protéines lors de la synthèse (pic). La décroissance prouve que ces protéines marquées sont dégradées (protéolyse) et remplacées par de nouvelles protéines synthétisées à partir de leucine non marquée (provenant de l'alimentation ou du recyclage). Les protéines corporelles ne sont pas des structures fixes : elles sont en \textbf{renouvellement constant}.
\textbf{Conclusion :} L'état stationnaire dynamique est une réalité expérimentale ; la masse protéique constante résulte de l'égalité temporaire entre taux de synthèse et taux de dégradation.
\end{experience}

\coursec{Les mécanismes de synthèse moléculaire}

\courssub{Synthèse des protéines : transcription et traduction}
Le renouvellement des protéines repose sur l'expression génique. Pour chaque protéine à renouveler :
\begin{enumerate}
  \item \textbf{Transcription} (noyau) : l'ADN sert de matrice pour synthétiser un ARNm mature (épissage, capage, polyadénylation).
  \item \textbf{Traduction} (cytoplasme, ribosomes) : l'ARNm est lu par les ribosomes ; les acides aminés (apportés par les ARNt) sont assemblés selon le code génétique. L'énergie est fournie par le GTP (facteurs d'initiation, d'élongation, de terminaison) et l'ATP (activation des acides aminés par les aminoacyl-ARNt synthétases).
  \item \textbf{Pliage et maturation} : chaperonnes (Hsp70, chaperonines), modifications post-traductionnelles (glycosylation, phosphorylation, clivage propeptidique), adressage vers la destination finale.
\end{enumerate}
La vitesse de synthèse est régulée principalement au niveau transcriptionnel (facteurs de transcription, épigénétique) et translationnel (voie mTOR, phosphorylation d'eIF2$\alpha$).

\courssub{Synthèse des lipides membranaires}
Les phospholipides sont synthétisés au niveau du réticulum endoplasmique lisse (REL) à partir de précurseurs (glycérol-3-phosphate, acyl-CoA, CDP-choline/éthanolamine). Le cholestérol est synthétisé dans le cytosol/REL via la voie du mévalonate (enzyme clé : HMG-CoA réductase). Le renouvellement membranaire implique un trafic vésiculaire constant (exocytose/endocytose) et un remodelage enzymatique (phospholipases, acyltransférases).

\courssub{Réparation et réplication de l'ADN}
Hors division, le renouvellement de l'ADN se limite à la \textbf{synthèse de réparation} : excision de bases endommagées (BER), excision de nucléotides (NER), réparation des brisures double brin (HR, NHEJ). Les polymérases $\delta$, $\varepsilon$ (réparation/réplication) et $\beta$ (BER) incorporent des dNTP neufs. En phase S, la réplication semi-conservative renouvelle l'intégralité du génome.

\coursec{Les mécanismes de dégradation moléculaire}

\courssub{Voies de dégradation des protéines}
Deux systèmes majeurs coexistent :
\begin{itemize}
  \item \textbf{Système ubiquitine-protéasome (UPS)} : dégradation sélective des protéines cytoplasmiques, nucléaires, mal repliées (ERAD) ou régulatrices (cyclines, facteurs de transcription). Marquage par une chaîne de polyubiquitine (liaison Lys48) $\rightarrow$ reconnaissance par le protéasome 26S $\rightarrow$ hydrolyse ATP-dépendante en peptides courts $\rightarrow$ recyclage des acides aminés. Demi-vies courtes (min à h).
  \item \textbf{Autophagie-lysosome} : dégradation en vrac (macroautophagie) ou sélective (mitophagie, ribophagie, chaperon-mediated autophagy). Formation d'un autophagosome $\rightarrow$ fusion avec le lysosome $\rightarrow$ hydrolyse par hydrolases acides (cathepsines). Gestion du stress (jeûne, hypoxie), turnover des organites, protéines d'agrégation. Demi-vies longues (h à jours).
\end{itemize}

\courssub{Dégradation des lipides et glucides}
\begin{itemize}
  \item \textbf{Lipides} : phospholipases (A1, A2, C, D) hydrolysent les phospholipides membranaires $\rightarrow$ acides gras libres + lysophospholipides (réutilisés ou $\beta$-oxydation). Cholestérol : efflux vers HDL (transport inverse) ou conversion en acides biliaires (foie).
  \item \textbf{Glycogène} : glycogénolyse (glycogène phosphorylase $\rightarrow$ glucose-1-P $\rightarrow$ glucose-6-P) ; dans le foie, glucose-6-phosphatase libère du glucose libre dans le sang ; dans le muscle, utilisation locale en glycolyse.
\end{itemize}

\begin{propriete}[Équation du turnover à l'état stationnaire]
Pour une population moléculaire de masse $M$ constante :
\[
\frac{dM}{dt} = v_{\text{synthèse}} - v_{\text{dégradation}} = 0 \quad \Rightarrow \quad v_{\text{synthèse}} = v_{\text{dégradation}}
\]
La demi-vie $t_{1/2}$ est liée à la constante de dégradation $k_d$ (loi du premier ordre) par $t_{1/2} = \frac{\ln 2}{k_d}$. Au steady-state, le flux de synthèse $J = k_d \cdot M$.
\end{propriete}

\coursec{Exploitation de données de marquage isotopique}

\courssub{Principe du marquage isotopique}
Pour rendre visible ce flux invisible, les biologistes utilisent des \cle{isotopes stables ou radioactifs} (\ce{^{15}N}, \ce{^{13}C}, \ce{^{14}C}, \ce{^{3}H}) incorporés dans les précurseurs métaboliques : acides aminés marqués pour les protéines, acides gras marqués pour les lipides, nucléotides marqués pour l'ADN. Après administration (injection, perfusion, alimentation), on suit la radioactivité ou l'enrichissement isotopique dans les macromolécules cellulaires au cours du temps.

\begin{popfigure}
\begin{tikzpicture}[scale=0.9, every node/.style={font=\small}]
  % Précurseur marqué
  \node[draw, rounded corners, fill=popblueL, thick] (prec) at (0,0) {Acide aminé marqué (\ce{^{15}N})};
  % Flèche incorporation
  \draw[->, thick, popblue] (prec.east) -- ++(1.5,0) node[midway, above] {Incorporation (synthèse)} coordinate (mid1);
  % Protéine marquée
  \node[draw, rounded corners, fill=popgreenL, thick] (prot) at (4,0) {Protéine marquée};
  % Flèche dégradation
  \draw[->, thick, poporange] (prot.east) -- ++(1.5,0) node[midway, above] {Dégradation (protéolyse)} coordinate (mid2);
  % Acide aminé libre marqué
  \node[draw, rounded corners, fill=poporangeL, thick] (free) at (7.5,0) {Acide aminé libre marqué};
  % Flèche réutilisation / excrétion
  \draw[->, thick, poppurple] (free.south) -- ++(0,-1) -| (prec.south) node[midway, left] {Réutilisation / Excrétion};
  % Légendes des compartiments
  \node[align=center, text width=2.5cm] at (0,-1.5) {Milieu extracellulaire / alimentation};
  \node[align=center, text width=2.5cm] at (4,-1.5) {Compartiment protéique cellulaire};
  \node[align=center, text width=2.5cm] at (7.5,-1.5) {Pool d'acides aminés libres};
\end{tikzpicture}
\legende{Schéma du principe du marquage isotopique : un acide aminé marqué (\ce{^{15}N}) est incorporé dans une protéine lors de la synthèse ; après dégradation de cette protéine, l'acide aminé marqué est libéré dans le pool libre et peut être réutilisé pour une nouvelle synthèse ou excrété.}
\end{popfigure}

\courssub{Notion de demi-vie moléculaire}
La \cle{demi-vie moléculaire} (notée $t_{1/2}$) est la durée nécessaire pour que la moitié des molécules marquées d'une population donnée soit remplacée par des molécules non marquées. Elle se lit directement sur une courbe de décroissance du marquage isotopique (phase de \emph{chase} ou de dilution).

\begin{definition}[Demi-vie moléculaire]
Durée au bout de laquelle 50 \% des molécules d'une population donnée ont été remplacées par des molécules neuves (non marquées). Elle reflète la vitesse de dégradation ($k_d = \ln 2 / t_{1/2}$).
\end{definition}

\courssub{Diversité des demi-vies : de quelques minutes à plusieurs mois}
\begin{itemize}
  \item \textbf{Protéines régulatrices / enzymes de contrôle} (ex. : phosphofructokinase-1, cyclines, HIF-1$\alpha$, p53) : demi-vie de quelques minutes à quelques heures — permet une adaptation métabolique et transcriptionnelle rapide.
  \item \textbf{Protéines structurales / d'abondance} (ex. : actine, tubuline, hémoglobine) : demi-vie de quelques jours à quelques semaines.
  \item \textbf{Protéines de la matrice extracellulaire} (ex. : collagène, élastine) : demi-vie de plusieurs mois à années — assure la stabilité tissulaire (cicatrisation lente).
  \item \textbf{Lipides membranaires} : demi-vie de quelques jours à semaines (remodelage constant).
  \item \textbf{ADN} : stabilité de la séquence (pas de turnover global hors division), mais turnover local des nucléotides lors des réparations (demi-vie des segments réparés : heures à jours).
\end{itemize}

\courssub{Illustration : courbe de décroissance exponentielle du marquage isotopique}
La figure ci-dessous montre l'évolution du pourcentage de marquage isotopique (\% \ce{^{15}N}) dans une protéine hépatique après une injection unique de leucine marquée (pulse-chase). La courbe suit une décroissance exponentielle $M(t)=M_0 e^{-k_d t}$. La demi-vie se lit graphiquement.

\begin{popfigure}
\begin{tikzpicture}
\begin{axis}[
  width=\linewidth,
  height=7cm,
  xlabel={Temps (jours)},
  ylabel={\% marquage \ce{^{15}N}},
  domain=0:5,
  samples=200,
  axis lines=middle,
  xmin=0, xmax=5.5,
  ymin=0, ymax=105,
  xtick={0,1,2,3,4,5},
  ytick={0,25,50,75,100},
  grid=major,
  grid style={popline},
  tick label style={font=\small},
  label style={font=\small},
  clip=false
]
% Courbe exponentielle : M0=100, k=ln(2)/1 => demi-vie 1 jour
\addplot[popcurve, thick, popblue] {100*exp(-ln(2)*x)};
% Ligne horizontale à 50%
\addplot[dashed, poporange, thick] coordinates {(0,50) (1,50)};
% Ligne verticale à t=1
\addplot[dashed, poporange, thick] coordinates {(1,0) (1,50)};
% Points
\addplot[only marks, mark=*, mark size=2, popblue] coordinates {(0,100) (1,50) (2,25) (3,12.5)};
% Annotations
\node[above right, font=\small, popblue] at (axis cs:0,100) {$M_0=100\%$};
\node[above left, font=\small, poporange] at (axis cs:1,50) {$t_{1/2}=1$ jour};
\node[below right, font=\small, popblue] at (axis cs:2,25) {$25\%$};
\end{axis}
\end{tikzpicture}
\legende{Courbe de décroissance du marquage isotopique (\ce{^{15}N}) dans une protéine hépatique après un \emph{pulse-chase}. La demi-vie $t_{1/2}$ se lit à l'intersection de la courbe avec la ligne horizontale à 50 \% du marquage initial $M_0$. Ici $t_{1/2}\approx 1$ jour.}
\end{popfigure}

\courssub{Méthode : exploiter une courbe de marquage isotopique}
\begin{methode}[Exploitation d'une courbe de marquage isotopique (pulse-chase)]
\begin{enumerate}
  \item Repérer le maximum d'incorporation ($M_0$) sur l'ordonnée (au temps $t=0$ du \emph{chase} ou au pic d'incorporation).
  \item Tracer une ligne horizontale à l'ordonnée $M_0/2$ (50 \% du maximum).
  \item Lire l'abscisse correspondante sur l'axe des temps : c'est la demi-vie $t_{1/2}$.
  \item Comparer les demi-vies de différentes molécules pour conclure sur leur vitesse de renouvellement (demi-vie courte = renouvellement rapide = protéine régulatrice ; demi-vie longue = renouvellement lent = protéine structurelle).
  \item Vérifier la cinétique : une droite en échelle semi-log (ln(marquage) vs temps) confirme une cinétique du premier ordre (dégradation proportionnelle à la quantité présente).
\end{enumerate}
\end{methode}

\courssub{Exemple chiffré : détermination de la demi-vie d'une protéine hépatique}
\begin{exemplebox}[Calcul de demi-vie à partir d'un résidu de marquage]
Une protéine hépatique présente 25 \% de marquage résiduel après 2 jours. On sait que la décroissance suit une loi exponentielle du premier ordre. Après une demi-vie, le marquage passe de 100 \% à 50 \%; après deux demi-vies, de 50 \% à 25 \%. Puisque 25 \% correspond à deux demi-vies écoulées en 2 jours, une demi-vie vaut $2\,\text{jours} / 2 = 1\,\text{jour}$.
\end{exemplebox}

\courssub{Erreur fréquente à éviter}
\begin{attention}[Piège : confusion stabilité compositionnelle et immobilité moléculaire]
Ne pas croire qu'une cellule « stable » contient les mêmes molécules toute sa vie. La stabilité de la composition globale (pourcentages en protéines, lipides, ADN) cache un renouvellement incessant. C'est précisément ce flux continu qui permet la réparation, l'adaptation et la réponse aux signaux. Au Probatoire, on exige la distinction entre \textbf{constance de la composition} (état stationnaire) et \textbf{flux de matière} (turnover).
\end{attention}

\coursec{Renouvellement moléculaire et intégrité cellulaire}

\courssub{Le turnover comme garant de la qualité protéique (protéostasie)}
Le renouvellement permanent permet d'éliminer les molécules endommagées (oxydation, glycation, déamidation, erreurs de pliage) et de les remplacer par des copies fonctionnelles.
\begin{itemize}
  \item \textbf{Contrôle qualité co-traductionnel} : ribosome-associated quality control (RQC), dégradation des chaînes nascentes aberrantes.
  \item \textbf{Contrôle qualité post-traductionnel} : chaperonnes (Hsp70, Hsp90) tentent le repliage ; échec $\rightarrow$ ubiquitination $\rightarrow$ protéasome.
  \item \textbf{Agrégats} : si la capacité de dégradation est dépassée (vieillissement, mutations), formation d'agrégats toxiques (maladies à conformation protéique : Alzheimer, Parkinson, Huntington).
\end{itemize}

\courssub{Adaptation aux signaux environnementaux}
Le turnover rapide des protéines régulatrices (facteurs de transcription, enzymes allostériques, récepteurs) permet à la cellule de modifier son phénotype en quelques minutes :
\begin{itemize}
  \item \textbf{Jeûne / réalimentation} : dégradation d'enzymes glycolytiques, synthèse d'enzymes néoglucogéniques (régulation transcriptionnelle + turnover).
  \item \textbf{Stress oxydant} : synthèse rapide d'enzymes antioxydantes (SOD, catalase, glutathion peroxydase) via Nrf2.
  \item \textbf{Réponse immunitaire} : turnover rapide des cytokines, chémokines, molécules du CMH.
\end{itemize}

\courssub{Exemple camerounais : cicatrisation et synthèse de collagène}
La cicatrisation de la plaie de Junior (15 ans) est rapide car ses fibroblastes dermiques ont un turnover élevé du collagène (synthèse massive de procollagène $\rightarrow$ maturation extracellulaire $\rightarrow$ réticulation). Chez Mama Odile (65 ans), la synthèse de collagène est ralentie (moins de facteurs de croissance, sénescence des fibroblastes, glycation du collagène existant le rendant résistant à la dégradation), d'où une cicatrisation retardée.

\coursec{Ralentissement du renouvellement et vieillissement cellulaire}

\courssub{Preuves expérimentales du ralentissement avec l'âge}
\begin{itemize}
  \item \textbf{Diminution de la synthèse protéique} : baisse de l'activité ribosomale, de l'initiation de la traduction (voie mTOR dysregulée), de la transcription (hétérochromatinisation).
  \item \textbf{Altération de la dégradation} : déclin de l'activité du protéasome (oxydation des sous-unités), dysfonction lysosomale (accumulation de lipofuscine), baisse de l'autophagie (gènes ATG moins exprimés).
  \item \textbf{Accumulation de dommages} : protéines carbonylées, AGEs (produits de glycation avancée), ADN muté, mitochondries dysfonctionnelles (ROS).
\end{itemize}

\courssub{Conséquences cellulaires et tissulaires}
\begin{itemize}
  \item \textbf{Perte de protéostasie} : agrégation protéique, dysfonction enzymatique, signalisation altérée.
  \item \textbf{Sénescence cellulaire} : arrêt de prolifération irréversible, phénotype sécrétoire pro-inflammatoire (SASP : IL-6, MMPs, TNF-$\alpha$) $\rightarrow$ inflammation chronique stérile (\emph{inflammaging}).
  \item \textbf{Épuisement des cellules souches} : turnover insuffisant pour maintenir le pool de cellules souches (intestin, peau, muscle, hématopoïèse) $\rightarrow$ atrophie tissulaire, fragilité.
  \item \textbf{Rigidification tissulaire} : collagène vieux, glycation croisée (pentosidine) $\rightarrow$ artères rigides (hypertension), peau ridée, cicatrisation fibreuse (kyste, chéloïde).
\end{itemize}

\begin{aretenir}[Renouvellement moléculaire et vieillissement]
Le vieillissement cellulaire n'est pas une simple usure passive, mais la conséquence d'un \textbf{déséquilibre progressif du turnover} : la synthèse ne compense plus la dégradation (perte de masse) ou la dégradation n'élimine plus les molécules endommagées (accumulation de déchets). Maintenir un turnover efficace (exercice, restriction calorique modérée, sommeil, antioxydants alimentaires) est une stratégie de prévention du vieillissement pathologique.
\end{aretenir}

\courssub{Le savais-tu ?}
\begin{savaistu}[Turnover corporel complet et datation au carbone 14]
Presque toutes les molécules de ton corps sont remplacées en quelques mois : protéines musculaires en ~2-3 mois, lipides membranaires en quelques semaines, glycogène en heures. Seuls l'ADN des neurones (non divisés) et l'émail dentaire (acellulaire) datent de ta naissance. La mesure du \ce{^{14}C} atmosphérique (pic des essais nucléaires 1955-1963) incorporé dans l'ADN lors de la dernière division permet de dater l'âge des cellules : les cardiomyocytes se renouvellent à ~1 \%/an à 25 ans, ~0,5 \%/an à 75 ans ; les hépatocytes ~1 an ; les adipocytes ~10 ans. Tu n'es littéralement plus « fait » des mêmes atomes qu'à la rentrée scolaire !
\end{savaistu}

\coursec{Exercices d'application}

\courssub{Exercice résolu : interprétation de données de marquage}
\begin{exoresolu}[Détermination de la demi-vie d'une enzyme régulatrice]
\textbf{Énoncé :} On suit le marquage au \ce{^{14}C} d'une enzyme hépatique après injection d'acétate marqué. Les mesures donnent : $t=0\,\text{h}$ : 100 \%; $t=2\,\text{h}$ : 70 \%; $t=4\,\text{h}$ : 50 \%; $t=6\,\text{h}$ : 35 \%; $t=8\,\text{h}$ : 25 \%. Déterminer la demi-vie de cette enzyme et commenter sa vitesse de renouvellement.

\textbf{Solution détaillée :}
\begin{enumerate}
  \item Le marquage initial $M_0 = 100\%$ (à $t=0$ du \emph{chase}).
  \item On cherche le temps où le marquage atteint $50\%$ ($M_0/2$). D'après le tableau, à $t=4\,\text{h}$ le marquage est exactement 50 \%.
  \item Donc la demi-vie $t_{1/2} = 4\,\text{h}$.
  \item Vérification cinétique : à $t=8\,\text{h}$ (deux demi-vies), marquage attendu = 25 \%. Mesuré = 25 \%. La cinétique est du premier ordre.
  \item \textbf{Commentaire :} Une demi-vie de 4 heures indique un renouvellement très rapide, typique des enzymes régulatrices du métabolisme intermédiaire (ex. : phosphofructokinase-1, pyruvate kinase) dont la concentration doit pouvoir varier rapidement en réponse aux signaux métaboliques (insuline, glucagon, AMPc, AMP).
\end{enumerate}
\end{exoresolu}

\courssub{Exercice proposé : comparaison de demi-vies}
\begin{exercice}[Comparaison du turnover de deux protéines hépatiques]
On a mesuré le marquage résiduel au \ce{^{15}N} de deux protéines hépatiques (A et B) après un \emph{pulse-chase} unique. Résultats :

\begin{center}
\begin{tabular}{|c|c|c|}\hline
Temps (jours) & Protéine A (\% marquage) & Protéine B (\% marquage) \\ \hline
0 & 100 & 100 \\
1 & 80 & 50 \\
2 & 64 & 25 \\
3 & 51 & 12,5 \\
4 & 41 & 6,25 \\ \hline
\end{tabular}
\end{center}

\textbf{Questions :}
\begin{enumerate}
  \item Tracer les courbes de décroissance des deux protéines sur un même graphique (papier millimétré ou logiciel).
  \item Déterminer graphiquement la demi-vie de chaque protéine.
  \item Classer les protéines A et B par vitesse de renouvellement croissante.
  \item Proposer une identité fonctionnelle probable pour A et B (enzyme régulatrice vs protéine structurelle) et justifier.
  \item Calculer la constante de dégradation $k_d$ (en jour$^{-1}$) pour chaque protéine.
\end{enumerate}
\end{exercice}

\begin{corrige}[Exercice : Comparaison du turnover de deux protéines hépatiques]
\begin{enumerate}
  \item \textbf{Graphique} : Courbes exponentielles décroissantes. Protéine B plus pentue que A.
  \item \textbf{Demi-vies} :
  \begin{itemize}
    \item Protéine A : $M_0=100$, $M_0/2=50$. Atteint 51 \% à $t=3$ j $\Rightarrow t_{1/2}(\text{A}) \approx 3,0$ jours.
    \item Protéine B : $M_0=100$, $M_0/2=50$. Atteint 50 \% à $t=1$ j $\Rightarrow t_{1/2}(\text{B}) = 1,0$ jour.
  \end{itemize}
  \item \textbf{Vitesse de renouvellement} : B (1 j) $>$ A (3 j). Protéine B se renouvelle 3 fois plus vite.
  \item \textbf{Identité probable} :
  \begin{itemize}
    \item Protéine B ($t_{1/2}=1$ j) : enzyme régulatrice (ex. : glucokinase, phosphofructokinase) ou protéine de signalisation à turnover rapide.
    \item Protéine A ($t_{1/2}=3$ j) : protéine d'abondance / structurelle modérée (ex. : albumine, transferrine, enzymes glycolytiques stables).
  \end{itemize}
  \item \textbf{Constantes de dégradation} : $k_d = \ln 2 / t_{1/2}$.
  \begin{align*}
    k_d(\text{A}) &= \frac{0,693}{3,0} \approx 0,231\,\text{jour}^{-1} \\
    k_d(\text{B}) &= \frac{0,693}{1,0} = 0,693\,\text{jour}^{-1}
  \end{align*}
\end{enumerate}
\end{corrige}

\coursec{Les mécanismes de synthèse moléculaire}

Dans la section précédente, nous avons mis en évidence, grâce aux expériences de marquage isotopique, que les molécules constitutives de la cellule — protéines, lipides, ADN — ne sont pas figées : elles disparaissent et réapparaissent en permanence. Une question fondamentale surgit alors : si les molécules sont dégradées, d'où proviennent les nouvelles ? Qui les fabrique, où, et à quel coût énergétique ? C'est l'objet de l'étude des \cle{mécanismes de synthèse moléculaire}, que nous allons explorer pas à pas.

\courssub{La synthèse moléculaire : un chantier permanent appelé anabolisme}

Imagine un quartier de Yaoundé en pleine mutation : des camions livrent des briques et du ciment, des maçons assemblent les parpaings, pendant que des engins démolissent les bâtiments vétustes pour laisser place au neuf. La cellule fonctionne sur ce même principe dynamique. Elle possède un chantier de construction permanent, l'\cle{anabolisme}, qui fabrique en continu de nouvelles molécules à partir d'unités simples.

\begin{definition}[Anabolisme]
L'\cle{anabolisme} est l'ensemble des réactions de \textbf{synthèse} des molécules constitutives de la cellule (protéines, lipides, glucides complexes, acides nucléiques) à partir de molécules plus simples appelées \cle{précurseurs} (acides aminés, acides gras, glycérol, nucléotides). Ces réactions sont \textbf{consommatrices d'énergie}, fournie sous forme d'\cle{ATP} (adénosine triphosphate). L'anabolisme s'oppose au \cle{catabolisme}, ensemble des réactions de dégradation qui, elles, libèrent de l'énergie.
\end{definition}

\begin{aretenir}[L'équilibre du renouvellement]
Le renouvellement moléculaire = anabolisme (synthèses) $+$ catabolisme (dégradations). Tant que les deux flux s'équilibrent, la cellule conserve sa masse et son intégrité structurelle tout en remplaçant ses constituants : c'est le \cle{turnover moléculaire}.
\end{aretenir}

\courssub{La synthèse des protéines : de l'ADN au ribosome}

\textbf{Une observation pour démarrer.} Un enfant qui grandit, une plaie qui se referme après une coupure, un sportif dont les muscles se développent : dans tous ces cas, l'organisme fabrique massivement des \cle{protéines}. Or les protéines sont des macromolécules complexes, formées de longues chaînes d'\cle{acides aminés} assemblés dans un ordre précis, dicté par l'information génétique. Comment la cellule connaît-elle cet ordre ? Où se fait l'assemblage ?

\textbf{Hypothèse.} L'information nécessaire à la fabrication de chaque protéine doit être stockée dans la cellule, puis transmise à un « atelier d'assemblage » situé hors du noyau.

\textbf{Les faits expérimentaux.} Les expériences classiques de marquage radioactif (acides aminés marqués au carbone 14, suivis par autoradiographie sur coupes cellulaires, menées notamment par \emph{Paul Zamecnik} et \emph{Mary Stephenson} dans les années 1950) ont montré que :
\begin{itemize}
\item l'information est stockée dans l'\cle{ADN}, situé dans le \textbf{noyau} ;
\item l'assemblage des acides aminés se fait sur les \cle{ribosomes}, petits granules ribonucléoprotéiques situés dans le cytoplasme (libres) ou accrochés à la face cytoplasmique du \cle{réticulum endoplasmique rugueux} (RER) ;
\item un intermédiaire, l'\cle{ARN messager} (ARNm), transporte l'information du noyau vers les ribosomes.
\end{itemize}

\textbf{Interprétation et conclusion.} La synthèse d'une protéine (plus rigoureusement : d'une chaîne polypeptidique) se déroule en deux grandes étapes spatialement et temporellement séparées chez les eucaryotes :

\begin{enumerate}
\item \textbf{La transcription} (dans le noyau) : la portion d'ADN correspondant à un \cle{gène} est recopiée sous forme d'une molécule d'ARN messager (pré-ARNm puis ARNm mature après maturation : épissage, coiffage, polyadénylation). C'est comme photocopier une page d'un livre de recettes (l'ADN, qui ne quitte jamais la bibliothèque du noyau) pour l'emporter en cuisine.
\item \textbf{La traduction} (sur les ribosomes, dans le cytoplasme) : le ribosome « lit » l'ARN messager selon le code génétique et assemble les acides aminés un à un, dans l'ordre dicté par la succession des codons, pour former la chaîne polypeptidique.
\end{enumerate}

\begin{propriete}[Le code génétique — propriété admise en Première D]
La synthèse d'une protéine suit le \cle{code génétique} : à chaque groupe de trois nucléotides de l'ARN messager (un \cle{codon}) correspond un acide aminé déterminé (ou un signal d'arrêt). \textbf{Chaque gène porte l'information nécessaire à la fabrication d'une chaîne polypeptidique (souvent appelée protéine par simplification).} Cette propriété est admise en Première D : le détail moléculaire du mécanisme (rôle des ARN de transfert, structure du ribosome, mutations, régulation) sera approfondi en classe de Terminale.
\end{propriete}

\begin{popfigure}
\begin{center}
\begin{tikzpicture}[font=\small, >=Stealth]
% ADN double brin
\draw[very thick, popblue] (0,0) -- (3,0);
\draw[very thick, popblue] (0,0.35) -- (3,0.35);
\foreach \x in {0.3,0.9,1.5,2.1,2.7} \draw[popblue!60] (\x,0) -- (\x,0.35);
\node[popblue, above] at (1.5,0.5) {\textbf{ADN (gène) — Noyau}};
% Flèche transcription
\draw[->, very thick, poporange] (3.3,0.175) -- (4.7,0.175) node[midway, above, popdark]{Transcription};
% ARNm (brin simple)
\draw[very thick, popgreen] (5,0.175) -- (7.5,0.175);
% Symboles nucléotides sur ARNm
\foreach \x in {5.2,5.8,6.4,7.0} \draw[popgreen!70!black] (\x,0.175) -- (\x,0.45);
\node[popgreen!70!black, above] at (6.25,0.55) {\textbf{ARN messager (ARNm)}};
% Flèche traduction
\draw[->, very thick, poporange] (7.8,0.175) -- (9.2,0.175) node[midway, above, popdark]{Traduction};
% Ribosome (deux sous-unités)
\draw[fill=poppurple!30, draw=poppurple] (10.1,0.5) circle (0.4);
\draw[fill=poppurple!30, draw=poppurple] (10.1,-0.15) circle (0.3);
\node[poppurple, below] at (10.1,-0.6) {\textbf{Ribosome}};
% Protéine naissante (bobine)
\draw[very thick, poppink, decorate, decoration={coil, aspect=0.4, segment length=3.5pt, amplitude=2.5pt}] (10.6,0.5) -- (12.6,0.5);
\node[poppink, above] at (11.6,0.7) {\textbf{Chaîne polypeptidique}};
\end{tikzpicture}
\end{center}
\legende{Schéma fléché simplifié de la synthèse d'une protéine : l'information du gène (ADN, dans le noyau) est transcrite en ARNm, qui migre dans le cytoplasme pour être traduit sur les ribosomes en une chaîne d'acides aminés (polypeptide).}
\end{popfigure}

\courssub{D'où viennent les acides aminés ? Le rôle de l'alimentation}

Le ribosome est l'atelier, l'ARN messager est le plan de montage, mais il manque la matière première : les \cle{acides aminés}. Il en existe 20 types différents codés par le génome. Certains peuvent être fabriqués par l'organisme (acides aminés non essentiels), mais 8 d'entre eux (9 ou 10 chez l'enfant en croissance), dits \cle{acides aminés essentiels} (ou indispensables), doivent obligatoirement être apportés par l'alimentation car l'organisme ne sait pas synthétiser leur chaîne carbonée.

Lors d'un repas composé de poisson braisé, de sauce arachide, de haricots (niébé) ou d'œufs — sources de protéines courantes sur nos tables camerounaises — les protéines alimentaires sont \textbf{digérées} (hydrolysées) dans l'estomac (pepsine) et l'intestin grêle (trypsine, chymotrypsine, peptidases). Elles sont découpées en acides aminés libres (et di/tripeptides). Ceux-ci sont absorbés par les entérocytes de l'intestin grêle, passent dans le sang (veine porte $\to$ circulation générale) qui les distribue à toutes les cellules du corps. Chaque cellule puise alors dans ce « stock circulant » (pool d'acides aminés plasmatiques et intracellulaires) les briques dont elle a besoin pour ses propres synthèses sur les ribosomes.

\begin{methode}[Relier alimentation et synthèse cellulaire — rédaction type Probatoire]
Pour expliquer le lien entre un aliment et le renouvellement moléculaire d'une cellule, structure ta réponse en trois étapes logiques :
\begin{enumerate}
\item \textbf{Identifier le nutriment} : l'aliment (poisson, arachides, haricots, œuf, chenilles) contient des protéines alimentaires.
\item \textbf{Suivre la digestion et l'absorption} : ces protéines sont hydrolysées par les enzymes digestives en \textbf{précurseurs} — les acides aminés — qui sont absorbés au niveau de l'intestin grêle et rejoignent la circulation sanguine.
\item \textbf{Montrer la réutilisation cellulaire} : les acides aminés sanguins pénètrent dans les cellules (transport actif) et servent à la synthèse de \emph{nouvelles} protéines cellulaires sur les ribosomes, remplaçant celles qui ont été dégradées (turnover).
\end{enumerate}
\textbf{Conclusion type :} « Les protéines du poisson ne s'installent pas telles quelles dans nos muscles ; elles sont démontées en acides aminés, qui servent de briques pour reconstruire nos propres protéines spécifiques. »
\end{methode}

\begin{attention}[Piège fréquent au Probatoire]
Ne dis jamais : « Les protéines alimentaires deviennent directement nos protéines » ou « passent dans le sang ». Une protéine étrangère (non-soi) ne peut pas être insérée telle quelle dans une cellule humaine : elle serait reconnue comme antigène et déclencherait une réponse immune, ou serait dégradée. Seuls les \textbf{acides aminés} (et quelques peptides courts) issus de la digestion complète sont réutilisés comme briques de construction universelles.
\end{attention}

\courssub{La synthèse des lipides et le renouvellement des membranes}

Les protéines ne sont pas les seules molécules renouvelées. Les \cle{lipides}, constituants essentiels des \textbf{membranes cellulaires} (membrane plasmique, membranes des organites : noyau, mitochondries, réticulum, Golgi, vésicules), sont eux aussi fabriqués en permanence.

Les lipides membranaires majeurs (phospholipides : phosphatidylcholine, phosphatidylethanolamine, etc. ; cholestérol) sont assemblés à partir de deux types de précurseurs : les \cle{acides gras} et le \cle{glycérol} (ou le glycérol-3-phosphate), issus pour partie de l'alimentation (huile de palme, arachides, avocats, produits laitiers) et pour partie de transformations métaboliques internes (néoglucogenèse, lipogenèse à partir du glucose). Cet assemblage a lieu principalement dans le \cle{réticulum endoplasmique lisse} (REL), un réseau de membranes du cytoplasme dépourvu de ribosomes — contrairement au réticulum endoplasmique rugueux (RER), couvert de ribosomes et spécialisé dans la synthèse des protéines destinées à la sécrétion ou aux membranes.

Les lipides neufs (et les protéines membranaires synthétisées au RER) sont ensuite transportés par vésiculation vers l'\cle{appareil de Golgi}, où ils peuvent être modifiés, puis adressés vers les différentes membranes cibles. Ainsi, la membrane plasmique d'une cellule n'est jamais « la même » d'une semaine à l'autre : ses lipides sont continuellement remplacés (flip-flop, échange vésiculaire), ce qui garantit son étanchéité, sa fluidité et sa capacité à intégrer de nouveaux récepteurs.

\courssub{La réplication et la réparation de l'ADN}

L'ADN est une molécule particulièrement stable chimiquement, mais elle aussi connaît deux formes de synthèse distinctes par leur ampleur, leur moment et leur finalité :

\begin{itemize}
\item \textbf{La réplication de l'ADN} : avant chaque \textbf{division cellulaire} (phase S du cycle cellulaire), la cellule copie \emph{intégralement} son génome. La double hélice s'ouvre localement (fourche de réplication) et chaque brin parental sert de moule pour la synthèse d'un brin nouveau et complémentaire, grâce à des enzymes complexes dont les \cle{ADN polymérases}. Résultat : deux molécules d'ADN identiques (semi-conservatif), une pour chaque cellule fille. La réplication est un événement \textbf{ponctuel}, strictement couplé au cycle cellulaire.
\item \textbf{La réparation de l'ADN} : tout au long de la vie de la cellule (phase G1, S, G2, G0), l'ADN subit des agressions physiques (rayons ultraviolets du soleil, rayonnements ionisants), chimiques (métaux lourds, hydrocarbures, chimiothérapie) ou endogènes (radicaux libres ROS issus de la respiration mitochondriale, erreurs de réplication) qui provoquent des lésions (cassures, bases modifiées, dimères de thymine). Des \textbf{systèmes enzymatiques de correction} (excision de base, excision de nucléotide, réparation d'appariement erroné, réparation par recombinaison homologue, jonction non homologue) patrouillent en permanence, détectent les anomalies, excisent la zone lésée et resynthétisent un fragment correct à partir du brin complémentaire intact. La réparation est un processus \textbf{continu}, vital pour la préservation de l'intégrité du génome.
\end{itemize}

\begin{attention}[Réplication ou réparation ? L'erreur classique à ne pas commettre]
Beaucoup d'élèves confondent ces deux mécanismes. Retiens bien la distinction opérationnelle :
\begin{itemize}
\item \textbf{Réplication} = copie \emph{complète} du génome, \emph{avant la division cellulaire} (phase S), pour doter les deux cellules filles. Enzyme clé : ADN polymérase $\delta$/$\varepsilon$ (eucaryotes).
\item \textbf{Réparation} = correction \emph{locale} des erreurs et lésions, \emph{en continu} (tout au long de l'interphase), pour préserver l'information génétique de la cellule \emph{existante}. Enzymes clés : glycosylases, endonucléases, ADN polymérase $\beta$, ligases.
\end{itemize}
Au Probatoire, préciser le \textbf{moment} (phase S vs continu), l'\textbf{ampleur} (génome entier vs fragment) et la \textbf{finalité} (division vs survie/intégrité) fait toute la différence pour obtenir les points.
\end{attention}

\courssub{Localisation cellulaire et finition : le trajet d'une protéine}

Les synthèses ne se font pas n'importe où : la cellule eucaryote est une usine compartimentée en ateliers spécialisés. Suivons le trajet complet d'une protéine destinée à être sécrétée (ex: insuline, anticorps, enzymes digestives) ou insérée dans la membrane plasmique (récepteurs, canaux) — on parle de \textbf{voie sécrétrice} :

\begin{enumerate}
\item \textbf{Noyau} : transcription du gène en pré-ARNm, maturation (épissage, cap 5', queue poly-A 3') $\to$ ARNm mature exporté par les pores nucléaires.
\item \textbf{Ribosomes} attachés au \textbf{Réticulum Endoplasmique Rugueux (RER)} : traduction co-traductionnelle ; la chaîne naissante pénètre dans la lumière du RER (signal peptide reconnu par le SRP). Premiers repliements (chaperonnes) et modifications (glycosylation N-liée initiale, formation ponts disulfure).
\item \textbf{Vésicules de transport COPII} : bourgeonnement du RER $\to$ appareil de Golgi (face \emph{cis}).
\item \textbf{Appareil de Golgi} : \textbf{finition} (glycosylation O-liée, troncature/remodelage chaînes oligosaccharidiques, sulfatation, clivage pro-protéine) et \textbf{adressage} (tri dans la face \emph{trans} : étiquetage Man-6-P pour lysosomes, vésicules constitutives vs régulées pour membrane/extérieur).
\item \textbf{Vésicules sécrétoires} : acheminement le long des microtubules (kinesine) $\to$ fusion avec la membrane plasmique (exocytose) ou intégration dans la membrane.
\end{enumerate}

\begin{popfigure}
\begin{center}
\begin{tikzpicture}[font=\footnotesize, >=Stealth]
% Contour cellule
\draw[very thick, popblue, fill=popblueL!30] (0,0) ellipse (6.5 and 3.6);
% Noyau
\draw[thick, poppurple, fill=poppurpleL] (-2.8,0.6) circle (1.4);
\draw[poppurple, fill=poppurple!60] (-2.8,0.6) circle (0.4);
\node[poppurple, align=center] at (-2.8,-1.1) {\textbf{1. Noyau}\\(Transcription\\ADN $\to$ ARNm)};
% RE Rugueux (serpent)
\draw[thick, poporange, decorate, decoration={snake, amplitude=1.5pt, segment length=7pt}] (-1.0,1.5) .. controls (0.8,2.2) and (2.5,1.9) .. (3.8,1.3);
\foreach \p in {(-0.4,1.7),(0.5,2.0),(1.4,1.95),(2.3,1.8),(3.2,1.55)} \fill[popdark] \p circle (0.08);
\node[poporange!80!black, align=center] at (1.5,2.6) {\textbf{2. RE Rugueux + Ribosomes}\\(Traduction, entrée dans la lumière)};
% Appareil de Golgi (cisternae)
\foreach \i in {0,1,2,3} {
  \draw[thick, popgreen!70!black, fill=popgreenL!50] (4.2,-0.5-0.3*\i) .. controls (4.8,-0.5-0.3*\i) and (5.5,-0.3-0.3*\i) .. (5.8,-0.3-0.3*\i) .. controls (6.0,-0.3-0.3*\i) and (6.0,-0.7-0.3*\i) .. (5.8,-0.7-0.3*\i) .. controls (5.5,-0.7-0.3*\i) and (4.8,-0.5-0.3*\i) .. (4.2,-0.5-0.3*\i) -- cycle;
}
\node[popgreen!50!black, align=center] at (5.5,-2.0) {\textbf{3. Appareil de Golgi}\\(Finition, Tri, Adressage)};
% Vésicules sécrétoires
\foreach \v in {(5.8,0.2), (6.2,-0.2), (5.5,-1.0), (6.0,-1.4)} \draw[popgreen!70!black, fill=popgreenL] \v circle (0.15);
% Mitochondrie
\draw[thick, poppink, fill=poppinkL] (-0.5,-2.2) ellipse (1.1 and 0.5);
\draw[poppink] (-1.4,-2.2) .. controls (-1.0,-1.9) and (-0.6,-2.5) .. (-0.2,-2.2) .. controls (0.1,-1.9) and (0.3,-2.5) .. (0.5,-2.2);
\node[poppink!70!black, align=center] at (-0.5,-3.0) {\textbf{Mitochondrie}\\(Production ATP)};
% Flèches trajet
\draw[->, very thick, popdark] (-1.5,1.0) -- (-0.6,1.4) node[midway, left, popdark, font=\tiny]{Export ARNm};
\draw[->, very thick, popdark] (3.8,1.1) -- (4.2,0.2) node[midway, right, popdark, font=\tiny]{Vésicules COPII};
\draw[->, very thick, popdark] (5.8,-0.2) -- (6.5,0.5) node[midway, right, popdark, font=\tiny]{Exocytose};
\draw[->, very thick, popdark, dashed] (5.8,-1.0) -- (6.5,-1.5) node[midway, right, popdark, font=\tiny]{Membrane/Extérieur};
% Membrane plasmique label
\node[popblue, align=center] at (0,3.3) {\textbf{Membrane plasmique}\\(Lipides \& Protéines renouvelés)};
\end{tikzpicture}
\end{center}
\legende{Cellule eucaryote et sites de synthèse moléculaire. Le trajet d'une protéine de la voie sécrétrice : (1) Transcription dans le noyau, (2) Traduction sur les ribosomes du RER et translocation, (3) Finition, glycosylation terminale et tri dans l'appareil de Golgi, (4) Transport vésiculaire vers la membrane plasmique (exocytose) ou l'espace extracellulaire. La mitochondrie fournit l'ATP nécessaire à toutes ces étapes énergivores.}
\end{popfigure}

\courssub{Le coût énergétique des synthèses : la facture ATP}

Aucun chantier ne fonctionne sans énergie. Assembler des acides aminés en protéines (liaisons peptidiques), des acides gras en triglycérides/phospholipides (liaisons ester), des nucléotides en ADN/ARN (liaisons phosphodiester) : chaque liaison chimique formée \textbf{consomme de l'énergie libre} (endergonique). Cette énergie est fournie par l'hydrolyse de l'\cle{ATP} (adénosine triphosphate) en ADP + Pi, et parfois du GTP, UTP, CTP.

L'ATP est régénérée essentiellement par les \cle{mitochondries} via la \textbf{respiration cellulaire} (glycolyse + cycle de Krebs + phosphorylation oxydative) : l'oxydation complète du glucose en présence de dioxygène libère de l'énergie, couplée à la synthèse d'ATP (chimiosmose). Bilan global simplifié :

\[ \ce{C6H12O6 + 6O2 -> 6CO2 + 6H2O} \quad \Delta G^\circ \approx -2870 \text{ kJ/mol} \quad \text{(environ 30 à 32 ATP récupérés)} \]

\begin{aretenir}[La chaîne énergétique du renouvellement]
Aliments (glucose, acides gras, acides aminés) $\xrightarrow{\text{digestion/absorption}}$ Sang $\xrightarrow{\text{transport}}$ Cellule $\xrightarrow{\text{respiration mitochondriale}}$ ATP $\xrightarrow{\text{hydrolyse}}$ Énergie pour synthèses (anabolisme) $\to$ Renouvellement des molécules cellulaires. C'est pourquoi une alimentation suffisante, équilibrée et une ventilation/pulmonaire efficaces (apport en $\ce{O2}$) sont indispensables au maintien de l'intégrité cellulaire.
\end{aretenir}

\begin{savaistu}[Des milliers de protéines chaque seconde — Contexte camerounais]
On estime qu'une cellule humaine active (hépatocyte, fibroblastes) synthétise \textbf{plusieurs milliers de molécules de protéines par seconde} ; certaines cellules très spécialisées, comme les cellules $\beta$ du pancréas (insuline) ou les plasmocytes (anticorps), en produisent des millions par minute. À l'échelle de l'organisme entier, \textbf{200 à 300 g de protéines} sont hydrolysées et resynthétisées chaque jour (turnover protéique total), bien au-delà de l'apport alimentaire journalier (50-70 g). Cela explique le \textbf{besoin quotidien en protéines alimentaires} (recommandation OMS : $\approx$ 0,83 g/kg/jour chez l'adulte, 1,2 g/kg/jour chez la femme enceinte/allaitante et l'enfant en croissance). Les campagnes de nutrition au Cameroun (MINSANTÉ, ONG) insistent à juste titre sur la consommation de sources de protéines accessibles et locales : haricots, niébé, arachides, soja, poisson fumé/séché, œufs, chenilles (protéines d'insectes très riches) — car sans apport régulier d'acides aminés essentiels, le chantier cellulaire tourne au ralenti : fonte musculaire, immunodépression, retard de croissance, œdèmes (kwashiorkor).
\end{savaistu}

\begin{exoresolu}[Exploiter des données : le coût énergétique de la synthèse protéique]
\textbf{Énoncé.} On estime que la formation d'une liaison peptidique entre deux acides aminés (activation de l'acide aminé $\to$ AA-AMP + translocation ribosome) consomme l'équivalent énergétique d'environ 4 molécules d'ATP (2 ATP pour l'activation par l'ARNt-synthétase, 2 GTP pour l'élongation/translocation sur le ribosome). Une protéine moyenne est formée de 300 acides aminés (soit 299 liaisons peptidiques).

1. Calcule le nombre de molécules d'ATP (équivalent) nécessaires à la synthèse d'une molécule de cette protéine.

2. Une cellule musculaire active synthétise $10^6$ molécules de cette protéine par minute. Calcule la consommation totale d'équivalents ATP par minute pour cette seule synthèse.

3. Sachant que la respiration cellulaire produit au maximum environ 30 molécules d'ATP par molécule de glucose oxydée, combien de molécules de glucose la cellule doit-elle oxyder par minute pour couvrir ce seul besoin ? Conclus sur l'importance du métabolisme énergétique.
\tcblower
\textbf{Solution détaillée.}

1. Nombre d'équivalents ATP par molécule de protéine :
\[ 299 \text{ liaisons peptidiques} \times 4 \text{ ATP/liaison} = \num{1196} \text{ ATP} \approx \num{1,2 \times 10^3} \text{ ATP par protéine.} \]

2. Consommation totale pour $10^6$ protéines par minute :
\[ 10^6 \times \num{1196} = \num{1,196 \times 10^9} \approx \num{1,2 \times 10^9} \text{ ATP par minute.} \]

3. Nombre de molécules de glucose à oxyder (rendement 30 ATP/glucose) :
\[ \frac{\num{1,2 \times 10^9} \text{ ATP}}{30 \text{ ATP/glucose}} = \num{4 \times 10^7} \text{ molécules de glucose par minute.} \]

\textbf{Conclusion.} La synthèse des protéines représente un coût énergétique considérable : \textbf{40 millions de molécules de glucose par minute} doivent être oxydées par les mitochondries d'une \emph{seule} cellule musculaire, rien que pour assurer le renouvellement d'un seul type de protéine. Cela justifie l'abondance exceptionnelle des mitochondries dans les cellules à forte activité synthétique (muscles cardiaque/squelettique, foie, rein, pancréas) et confirme que l'anabolisme est un processus \textbf{fortement consommateur d'énergie}, dépendant directement de la disponibilité des substrats énergétiques (glucose, acides gras) et du dioxygène.
\end{exoresolu}

En résumé, la cellule dispose d'une véritable industrie de synthèse compartimentée : l'ADN du noyau détient les plans (gènes), les ribosomes (libres ou sur RER) assemblent les polypeptides, le réticulum endoplasmique lisse fabrique les lipides, l'appareil de Golgi assure la finition post-traductionnelle et l'adressage vésiculaire, et les mitochondries paient la facture énergétique en ATP. Mais toute usine performante doit aussi gérer ses déchets, recycler ses matériaux usés et démolir ses structures obsolètes : c'est le rôle des mécanismes de dégradation moléculaire (protéasomes, lysosomes, autophagy), que nous étudierons dans la section suivante.

\coursec{Les mécanismes de dégradation moléculaire}

Dans la section précédente, nous avons vu comment la cellule fabrique en permanence de nouvelles molécules : synthèse des protéines à partir des acides aminés, synthèse des lipides, réplication de l'ADN. Mais une question se pose immédiatement : si la cellule synthétise sans cesse, pourquoi ne gonfle-t-elle pas indéfiniment ? Pourquoi sa composition reste-t-elle stable ? La réponse tient en un mot : la \cle{dégradation}. À chaque molécule synthétisée correspond, en moyenne, une molécule détruite. C'est ce second versant du renouvellement moléculaire que nous allons maintenant explorer.

\courssub{Pourquoi la cellule dégrade-t-elle ses propres molécules ?}

\textbf{Une situation concrète pour commencer.} Dans un hôpital de Yaoundé, le service de maintenance ne se contente pas d'acheter du matériel neuf : il répare ou élimine les appareils défectueux, faute de quoi les salles seraient encombrées de machines hors d'usage qui gêneraient le travail des soignants. La cellule fonctionne exactement ainsi. Ses molécules, soumises en permanence à des agressions (chaleur, radicaux libres, erreurs de synthèse), s'usent, se déforment ou deviennent inutiles. Une protéine mal repliée peut même devenir \emph{toxique} en s'agrégeant avec d'autres : c'est le cas dans la maladie d'Alzheimer, où l'accumulation de protéines anormales non dégradées détruit progressivement les neurones.

\begin{attention}[La dégradation n'est pas un gaspillage]
Erreur fréquente dans les copies : présenter la dégradation comme une perte, un « gaspillage » d'énergie. C'est faux. La dégradation est un mécanisme de \textbf{contrôle qualité} (élimination des molécules défectueuses) et de \textbf{recyclage} (récupération des acides aminés et acides gras pour de nouvelles synthèses). Sans elle, la cellule s'empoisonnerait avec ses propres déchets.
\end{attention}

La cellule dégrade donc ses molécules pour trois raisons essentielles :
\begin{itemize}
\item \textbf{éliminer les molécules usées ou endommagées} (protéines oxydées, membranes altérées) ;
\item \textbf{détruire les molécules défectueuses} dès leur fabrication (protéines mal repliées) ;
\item \textbf{supprimer les molécules devenues inutiles} (enzymes d'une voie métabolique arrêtée, protéines d'une phase du cycle cellulaire terminée).
\end{itemize}

\begin{definition}[Catabolisme]
Le \cle{catabolisme} est l'ensemble des réactions de dégradation des molécules de la cellule. Ces réactions, \textbf{exergoniques (libératrices d'énergie)}, produisent des précurseurs recyclables (acides aminés, acides gras, nucléotides) qui alimentent les nouvelles synthèses (anabolisme).
\end{definition}

\courssub{La dégradation des protéines : le protéasome et les lysosomes}

\textbf{Observation.} Au microscope électronique, le cytoplasme des cellules contient de petites vésicules denses, les \cle{lysosomes}, et des structures protéiques en forme de tonneau, les \cle{protéasomes}. Des expériences de marquage montrent que les protéines destinées à être détruites transitent par l'une ou l'autre de ces deux voies.

\textbf{La voie du protéasome : la destruction ciblée.} La cellule dispose d'un système d'étiquetage remarquable : une petite protéine, l'\cle{ubiquitine}, est fixée sur les protéines à détruire, comme une étiquette « à jeter ». Une protéine marquée par une chaîne de plusieurs molécules d'ubiquitine (polyubiquitination) est reconnue par le \cle{protéasome} (complexe 26S), un complexe enzymatique cylindrique qui la déroule, la déplie et la découpe en petits peptides puis en acides aminés. Cette voie concerne surtout les protéines intracellulaires de courte durée de vie : protéines mal repliées, protéines régulatrices à éliminer rapidement (cyclines, facteurs de transcription).

\textbf{La voie lysosomiale : l'estomac de la cellule.} Le \cle{lysosome} est une vésicule entourée d'une membrane unique, contenant une cinquantaine d'\cle{enzymes digestives} (protéases, lipases, nucléases) actives en milieu acide (pH voisin de 5, maintenu par des pompes à protons \ce{H+} ATPase). Il digère les matériaux englobés dans des vésicules qui fusionnent avec lui : particules ingérées par endocytose (hétérophagie), mais aussi constituants propres de la cellule (autophagie).

\begin{definition}[Autophagie (macroautophagie)]
L'\cle{autophagie} (du grec \emph{auto} : soi-même, et \emph{phagein} : manger) est le processus par lequel la cellule digère ses propres constituants altérés — organites usés (mitochondries, réticulum), agrégats de protéines — en les enfermant dans une vésicule à double membrane, l'\cle{autophagosome}, qui fusionne ensuite avec un lysosome pour former un \cle{autophagolysosome}.
\end{definition}

\textbf{Démarche d'investigation : comment a-t-on découvert l'autophagie ?}
\begin{itemize}
\item \emph{Observation} : des cellules privées de nutriments (jeûne) survivent plusieurs jours tout en voyant leur cytoplasme se remplir de vésicules à double membrane contenant des organites.
\item \emph{Hypothèse} : la cellule digère une partie de ses propres organites pour recycler leurs constituants et survivre.
\item \emph{Expérience} : Yoshinori Ohsumi (années 1990) utilise des levures mutantes incapables de dégrader le contenu de leurs vacuoles (équivalent des lysosomes). Privées de nutriments, ces levures accumulent des vésicules non digérées, visibles au microscope, ce qui permet d'identifier les gènes de l'autophagie (\emph{ATG genes}).
\item \emph{Interprétation} : l'autophagie est un mécanisme génétiquement programmé, indispensable à la survie en cas de carence et à l'élimination des organites usés.
\item \emph{Conclusion} : l'autophagie est une voie majeure de dégradation et de recyclage, conservée chez tous les eucaryotes, y compris l'Homme.
\end{itemize}

\begin{savaistu}[Un prix Nobel pour l'autophagie]
Le prix Nobel de physiologie ou médecine 2016 a récompensé le Japonais \textbf{Yoshinori Ohsumi} pour ses travaux sur les mécanismes de l'autophagie. On sait aujourd'hui qu'un défaut d'autophagie est impliqué dans la maladie de Parkinson, certains cancers et le vieillissement. Des recherches explorent même si le jeûne intermittent, en stimulant l'autophagie, pourrait avoir des effets bénéfiques sur la santé.
\end{savaistu}

\begin{popfigure}
\begin{center}
\begin{tikzpicture}[scale=1, every node/.style={font=\small}]
% Autophagosome avec organite usé (double membrane)
\draw[thick, poporange, fill=poporange!15] (0,0) circle (1.5);
\draw[thick, poporange!70, dashed] (0,0) circle (1.62);
% organite usé (mitochondrie)
\draw[thick, poppurple, fill=poppurple!15] (0,0) ellipse (0.9 and 0.45);
\draw[poppurple, thick] (-0.7,0.1) .. controls (-0.4,0.35) and (-0.1,-0.15) .. (0.2,0.15) .. controls (0.4,0.35) and (0.6,0.05) .. (0.72,0.12);
\node[popdark] at (0,-1.05) {\footnotesize organite usé};
% Lysosome
\draw[thick, popblue, fill=popblue!15] (4.5,0) circle (0.85);
\foreach \x/\y in {4.2/0.3, 4.7/0.1, 4.4/-0.35, 4.75/-0.3, 4.25/-0.1}
  \fill[popblue] (\x,\y) circle (0.07);
\node[popdark] at (4.5,1.25) {\footnotesize lysosome};
\node[popdark] at (4.5,1.55) {\footnotesize (enzymes digestives)};
% Flèche de fusion
\draw[-{Stealth[length=3mm]}, very thick, popgreen] (1.75,0) -- (3.5,0) node[midway, above, popdark]{\footnotesize fusion};
% Autolysosome résultat
\draw[thick, popgreen, fill=popgreen!15] (8,0) circle (1.15);
\foreach \x/\y in {7.7/0.3, 8.2/0.15, 7.9/-0.3, 8.3/-0.25, 7.6/-0.1}
  \fill[popgreen!70!popdark] (\x,\y) circle (0.06);
\draw[popmuted, thick, dashed] (8,0) ellipse (0.55 and 0.28);
\node[popdark] at (8,-1.5) {\footnotesize autophagolysosome : digestion};
% Produits
\draw[-{Stealth[length=3mm]}, very thick, popgold] (8,-1.9) -- (8,-2.7);
\node[popdark, align=center] at (8,-3.15) {\footnotesize acides aminés, acides gras\\ \footnotesize recyclés ou éliminés};
% Légende autophagosome
\node[popdark] at (0,2.05) {\footnotesize autophagosome};
\draw[thin, popmuted] (0,1.95) -- (0,1.65);
\end{tikzpicture}
\end{center}
\legende{L'autophagie : un organite usé (mitochondrie) est enfermé dans une vésicule à double membrane (autophagosome) qui fusionne avec un lysosome ; les enzymes digestives dégradent l'organite et les produits (acides aminés, acides gras) sont recyclés dans le cytoplasme.}
\end{popfigure}

\courssub{La dégradation des lipides}

Les lipides membranaires usés sont également renouvelés. Des enzymes appelées \cle{lipases} (lipase acide lysosomiale, lipases cytoplasmiques), présentes notamment dans les lysosomes (lipophagie) et le cytoplasme, hydrolysent les lipides (triglycérides, phospholipides) en \cle{acides gras} et glycérol. Ces acides gras ont deux destins possibles :
\begin{itemize}
\item être \textbf{recyclés} dans la synthèse de nouveaux lipides membranaires (voie de Kennedy, voie des phosphoinositides) ;
\item être \textbf{oxydés} dans les mitochondries (\cle{$\beta$-oxydation}) pour fournir de l'énergie (ATP) — dégradation libératrice d'énergie, typique du catabolisme.
\end{itemize}
Les membranes cellulaires sont ainsi entièrement renouvelées en quelques jours, ce qui garantit leur fluidité et leur imperméabilité sélective.

\courssub{Le destin des produits de dégradation : recyclage ou élimination}

Que deviennent les « briques » libérées par la dégradation ? La cellule, excellente gestionnaire, applique une règle simple : \textbf{rien ne se perd, tout se transforme}.
\begin{itemize}
\item Les \textbf{acides aminés} issus de l'hydrolyse des protéines rejoignent le « stock » cytoplasmique (pool d'acides aminés) et servent à la synthèse de nouvelles protéines. Seul l'excédent est éliminé : chez l'Homme, l'azote excédentaire (groupe amino) est converti en \ce{NH3} (toxique) puis en urée (cycle de l'urée dans le foie) et excrétée par les reins.
\item Les \textbf{acides gras} sont réutilisés pour fabriquer de nouveaux lipides ou brûlés pour produire de l'énergie (\ce{CO2} + \ce{H2O}).
\item Les \textbf{nucléotides} issus de la dégradation des acides nucléiques sont partiellement recyclés (voie de sauvetage / \emph{salvage pathway}) ou dégradés en acide urique (excréteur chez l'Homme).
\end{itemize}

\begin{exemplebox}[Un chiffre qui surprend : 200 à 300 g de protéines par jour]
Un adulte dégrade et resynthétise environ \textbf{200 à 300 g de protéines par jour}, alors que son alimentation ne lui apporte souvent que 60 à 80 g de protéines quotidiennes. D'où vient la différence ? Du \textbf{recyclage des acides aminés} issus de la dégradation interne. Sans ce recyclage, nos besoins alimentaires en protéines (viande, poisson, niébé, arachides) seraient impossibles à couvrir. Le renouvellement moléculaire repose donc sur une économie circulaire à l'échelle du corps.
\end{exemplebox}

\courssub{L'équilibre synthèse--dégradation : la clé de la stabilité cellulaire}

Nous pouvons maintenant rassembler les deux versants du renouvellement moléculaire dans un schéma d'ensemble.

\begin{popfigure}
\begin{center}
\begin{tikzpicture}[node distance=1.1cm, every node/.style={font=\small}]
\node[draw, thick, popblue, fill=popblue!15, rounded corners, minimum width=3.2cm, minimum height=1cm, align=center] (prec) {Précurseurs\\ (acides aminés, acides gras,\\ nucléotides)};
\node[draw, thick, popgreen, fill=popgreen!15, rounded corners, minimum width=3.2cm, minimum height=1cm, align=center, right=2.2cm of prec] (mol) {Molécules\\ fonctionnelles\\ (protéines, lipides, ADN)};
\node[draw, thick, poporange, fill=poporange!15, rounded corners, minimum width=3.6cm, minimum height=1cm, align=center, below=1.6cm of mol] (deg) {Dégradation\\ Protéasome / Lysosome\\ (Autophagie, Hétérophagie)};
\draw[-{Stealth[length=3mm]}, very thick, popgreen] (prec) -- (mol) node[midway, above, popdark]{Synthèse (Anabolisme)};
\draw[-{Stealth[length=3mm]}, very thick, poporange] (mol) -- (deg) node[midway, right, popdark, align=center]{Usure, Marquage\\ (Ubiquitine)};
\draw[-{Stealth[length=3mm]}, very thick, popblue] (deg.west) -| (prec.south) node[pos=0.25, below, popdark]{Recyclage};
\draw[-{Stealth[length=3mm]}, thick, popmuted] (deg.east) -- ++(1.6,0) node[right, popdark, align=left]{\footnotesize Élimination\\ \footnotesize (\ce{NH3} $\to$ urée, \ce{CO2})};
\draw[-{Stealth[length=3mm]}, thick, popgold] (prec.north) -- ++(0,0.9) node[above, popdark]{\footnotesize Apports alimentaires};
\end{tikzpicture}
\end{center}
\legende{Le cycle complet du renouvellement moléculaire : les précurseurs (d'origine alimentaire ou du recyclage) servent aux synthèses ; les molécules fonctionnelles usées sont dégradées par le protéasome ou les lysosomes ; les produits sont recyclés ou éliminés.}
\end{popfigure}

\begin{propriete}[État stationnaire de la cellule (admis)]
La composition de la cellule reste constante parce que la vitesse de synthèse est égale à la vitesse de dégradation : on dit que la cellule est en \cle{état stationnaire} (ou équilibre dynamique). Toute rupture prolongée de cet équilibre — dégradation excessive (fonte musculaire lors de dénutrition) ou insuffisante (accumulation de protéines anormales, maladies neurodégénératives) — menace la survie cellulaire.
\end{propriete}

\begin{exoresolu}[Raisonner sur l'équilibre synthèse--dégradation]
\textbf{Énoncé.} Chez un enfant souffrant de kwashiorkor (malnutrition protéique, fréquente lors des sevrages précoces au Cameroun), les apports alimentaires en acides aminés sont très insuffisants. Pourtant, les protéines musculaires continuent d'être dégradées. (1) Que devient l'équilibre synthèse--dégradation dans les muscles ? (2) Explique la fonte musculaire observée.
\tcblower
\textbf{Solution.} (1) Les protéines musculaires sont dégradées en permanence (renouvellement normal), mais faute d'acides aminés alimentaires, le stock de précurseurs cytoplasmiques s'épuise et la synthèse de nouvelles protéines ralentit fortement. La vitesse de dégradation devient supérieure à la vitesse de synthèse : l'état stationnaire est rompu. (2) Le bilan net est une perte de protéines : la masse musculaire diminue, c'est la fonte musculaire. Les acides aminés libérés sont détournés vers la synthèse de protéines vitales (enzymes, anticorps, albumine du sang) au détriment des protéines structurelles du muscle. Cet exemple montre que l'équilibre dynamique synthèse = dégradation est indispensable au maintien de l'intégrité des tissus.
\end{exoresolu}

\begin{aretenir}[L'essentiel de la section]
\begin{itemize}
\item La dégradation (\cle{catabolisme}) élimine les molécules usées, défectueuses ou inutiles ; ce n'est pas un gaspillage mais un contrôle qualité assorti d'un recyclage.
\item Les protéines marquées par l'\cle{ubiquitine} (polyubiquitination) sont détruites par le \cle{protéasome} (26S) ; les organites usés et agrégats sont digérés par \cle{autophagie} dans les \cle{lysosomes} (voie autophagosome $\to$ autophagolysosome).
\item Les lipases dégradent les lipides (lipophagie) ; acides aminés, acides gras et nucléotides sont recyclés dans de nouvelles synthèses ou éliminés (urée, \ce{CO2}, acide urique).
\item L'égalité entre vitesses de synthèse et de dégradation maintient l'\cle{état stationnaire} de la cellule ; sa rupture prolongée compromet la survie cellulaire (fonte musculaire, maladies neurodégénératives).
\end{itemize}
\end{aretenir}

\coursec{Exploitation de données de marquage isotopique}

Dans la section précédente, nous avons identifié les machineries moléculaires chargées de dégrader les protéines (protéasomes, lysosomes) et les acides nucléiques (nucléases), et nous avons vu que cette dégradation est indissociable d'une synthèse de remplacement. Mais comment \emph{prouver} qu'une molécule donnée dans une cellule vivante est effectivement détruite et remplacée par une neuve, alors que la cellule conserve son aspect et sa fonction ? L'intuition nous dit que si nous pouvions « étiqueter » les molécules présentes à un instant $t$, puis suivre le sort de cette étiquette au cours du temps, nous verrions l'étiquette disparaître des vieilles molécules pour apparaître dans les nouvelles. C'est exactement le principe du \textbf{marquage isotopique}, une démarche d'investigation historique qui a révélé le \cle{turnover moléculaire} (renouvellement permanent).

\courssub{Principes de la démarche expérimentale : le « pulse-chase »}

L'expérience type, dite \textbf{pulse-chase} (marquage-suivi), se déroule en trois temps logiques :

\begin{enumerate}
    \item \textbf{Le « pulse » (marquage) :} On administre à l'organisme (ou à une culture cellulaire) un \textbf{précurseur marqué} par un isotope (radioactif comme \ce{^{14}C}, \ce{^{3}H}, \ce{^{32}P} ; ou stable comme \ce{^{15}N}, \ce{^{13}C}). Ce précurseur est un constituant de base de la molécule d'intérêt (ex: leucine marquée pour les protéines, thymidine marquée pour l'ADN, acétate marqué pour les lipides). Pendant une courte période, les cellules incorporent ce précurseur \emph{à la place} de l'isotope naturel dans les molécules \emph{nouvellement synthétisées}. Seule la cohorte de molécules synthétisée pendant le pulse est marquée.
    \item \textbf{Le « chase » (poursuite) :} On arrête l'apport du précurseur marqué et on le remplace par le précurseur naturel non marqué (en excès). À partir de cet instant ($t=0$), \emph{aucune nouvelle molécule marquée n'est synthétisée}. Les molécules marquées existantes vont être dégradées selon leur cinétique propre, et remplacées par des molécules identiques mais \emph{non marquées}.
    \item \textbf{Les prélèvements et dosages :} À des intervalles de temps réguliers ($t_1, t_2, t_3...$), on prélève l'organe ou la cellule, on isole la molécule cible (par purification biochimique), et on mesure la \textbf{radioactivité résiduelle} (comptage scintillation pour \ce{^{14}C}, \ce{^{3}H}) ou le \textbf{rapport isotopique} (spectrométrie de masse pour \ce{^{15}N}, \ce{^{13}C}). Le résultat s'exprime en \textbf{\% de marquage initial} ou en \textbf{dpm/mg} (disintégrations par minute par milligramme de molécule).
\end{enumerate}

\begin{definition}[Traceur isotopique]
Un \textbf{traceur isotopique} est un atome dont le noyau possède un nombre de neutrons différent de l'isotope le plus abondant (ex: \ce{^{14}C} au lieu de \ce{^{12}C}), ce qui le rend détectable (radioactivité ou décalage de masse), tout en conservant des propriétés chimiques quasi identiques, lui permettant de suivre le même destin métabolique que l'isotope naturel.
\end{definition}

\courssub{Lecture et exploitation d'un tableau de données de marquage}

Considérons une expérience classique réalisée sur rat : injection unique de leucine-\ce{^{14}C} (pulse), puis suivi de la radioactivité dans deux protéines purifiées : une \textbf{protéine hépatique soluble} (enzymes du métabolisme intermédiaire) et le \textbf{collagène de tendon} (protéine structurale). Les données brutes sont présentées ci-dessous.

\begin{popfigure}
\legende{Tableau 1 : Données fictives d'une expérience de pulse-chase au \ce{^{14}C}-leucine chez le rat. Le marquage est exprimé en pourcentage du maximum atteint à J0 (fin du pulse).}
\begin{center}
\begin{tabularx}{\linewidth}{|c|X|X|}\hline
\rowcolor{popblueL}
\textbf{Temps (jours)} & \textbf{Protéine hépatique soluble (\% marquage initial)} & \textbf{Collagène de tendon (\% marquage initial)} \\ \hline
0 (fin du pulse) & 100 & 100 \\ \hline
1 & 65 & 99 \\ \hline
2 & 42 & 98 \\ \hline
3 & 27 & 97 \\ \hline
4 & 18 & 96 \\ \hline
5 & 11 & 95 \\ \hline
7 & 4,5 & 93 \\ \hline
10 & 1,2 & 90 \\ \hline
15 & 0,2 & 86 \\ \hline
20 & $\approx 0$ & 82 \\ \hline
30 & 0 & 74 \\ \hline
\end{tabularx}
\end{center}
\end{popfigure}

\textbf{Observation des données :}
\begin{itemize}
    \item Pour la \textbf{protéine hépatique}, le marquage chute brutalement : il est divisé par 2 entre J0 et J1 environ, puis devient quasi nul au bout de 10-15 jours. Cela signifie que la population de molécules marquée a été quasi entièrement dégradée et remplacée par des molécules non marquées.
    \item Pour le \textbf{collagène de tendon}, le marquage diminue à peine (100 $\to$ 82 en 20 jours). La cohorte marquée persiste très longtemps : le renouvellement est extrêmement lent.
\end{itemize}

\courssub{Construction et interprétation de la courbe de décroissance du marquage}

Pour visualiser la cinétique, on reporte le \% de marquage en fonction du temps. Comme la décroissance suit une loi exponentielle (voir propriété ci-après), l'utilisation d'une \textbf{échelle logarithmique en ordonnées} (graphique semi-logarithmique) redresse la courbe en ligne droite, facilitant la lecture de la demi-vie.

\begin{popfigure}
\legende{Figure 1 : Courbes de décroissance du marquage \ce{^{14}C} en échelle semi-logarithmique (ordonnées log). La protéine hépatique (courbe rouge, demi-vie courte) décroît vite ; le collagène (courbe bleue, demi-vie longue) décroît lentement. La pente est proportionnelle à la vitesse de renouvellement.}
\begin{tikzpicture}
\begin{semilogyaxis}[
    width=\linewidth,
    height=0.6\linewidth,
    xlabel={Temps (jours)},
    ylabel={\% Marquage initial},
    grid=both,
    minor grid style={dashed,popmuted},
    major grid style={solid,popline},
    xmin=0, xmax=20,
    ymin=0.5, ymax=120,
    xtick={0,2,4,6,8,10,12,14,16,18,20},
    ytick={1,2,5,10,20,50,100},
    legend pos=south west,
    legend cell align=left,
    domain=0:20,
    samples=200,
    popcurve/.style={thick, smooth},
]
% Protéine hépatique : demi-vie ~ 1.5 jours (rapide) -> k = ln(2)/1.5 = 0.462
% y = 100 * exp(-0.462*x)
\addplot[popred, popcurve] {100*exp(-0.462*x)};
\addlegendentry{Protéine hépatique ($t_{1/2} \approx 1,5$ j)}

% Collagène : demi-vie ~ 100 jours (lent) -> k = ln(2)/100 = 0.00693
% y = 100 * exp(-0.00693*x)
\addplot[popblue, popcurve] {100*exp(-0.00693*x)};
\addlegendentry{Collagène tendon ($t_{1/2} \approx 100$ j)}

% Ligne pour demi-vie protéine hépatique (50%)
\draw[popred, dashed, thin] (0,50) -- (1.5,50) -- (1.5,0.5);
\node[popred, anchor=north] at (1.5, 0.5) {$t_{1/2}$};

% Ligne pour 25% (2 demi-vies)
\draw[popred, dashed, thin] (0,25) -- (3,25) -- (3,0.5);
\node[popred, anchor=north] at (3, 0.5) {$2 t_{1/2}$};

\end{semilogyaxis}
\end{tikzpicture}
\end{popfigure}

\textbf{Interprétation :}
Sur un graphique semi-log, une décroissance exponentielle $M(t) = M_0 e^{-kt}$ donne une droite de pente $-k$ (constante de vitesse de dégradation). Plus la pente est raide (valeur absolue grande), plus $k$ est grand, plus le renouvellement est \textbf{rapide}. La protéine hépatique a une pente raide ; le collagène a une pente quasi horizontale.

\courssub{Détermination graphique de la demi-vie et comparaison de molécules}

La \textbf{demi-vie de renouvellement ($t_{1/2}$ ou $T_{1/2}$)} est le temps nécessaire pour que le marquage résiduel tombe à \textbf{la moitié} de sa valeur initiale. Sur la courbe semi-log, c'est l'abscisse du point où l'ordonnée vaut $M_0/2$.

\begin{methode}[Déterminer la demi-vie graphiquement et comparer les renouvellements]
\begin{enumerate}
    \item \textbf{Identifier le traceur et la cible :} Noter l'isotope utilisé (\ce{^{14}C}, \ce{^{3}H}...) et la molécule dosée (protéine, ADN, lipide) ainsi que le tissu d'origine.
    \item \textbf{Décrire l'évolution :} Lire le tableau ou la courbe : le marquage diminue-t-il ? Vite ou lentement ? Atteint-il un plateau non nul (marquage résiduel non renouvelé) ou tend-il vers 0 (renouvellement total) ?
    \item \textbf{Repérer $M_0$ :} Valeur du marquage à $t=0$ (fin du pulse). Souvent normalisé à 100\%.
    \item \textbf{Chercher $M_0/2$ :} Calculer la moitié de cette valeur (50\% si normalisé).
    \item \textbf{Lire l'abscisse correspondante :} Sur la courbe (linéaire ou semi-log), trouver le temps $t$ où l'ordonnée vaut $M_0/2$. C'est $t_{1/2}$.
    \item \textbf{Conclure et comparer :}
    \begin{itemize}
        \item $t_{1/2}$ court (heures à quelques jours) $\rightarrow$ \textbf{renouvellement rapide} (protéines régulatrices, enzymes hépatiques, ARNm).
        \item $t_{1/2}$ long (semaines à mois, années) $\rightarrow$ \textbf{renouvellement lent} (protéines structurales : collagène, kératine, cristallines du cristallin ; ADN nucléaire en phase $G_0$).
    \end{itemize}
\end{enumerate}
\end{methode}

\begin{propriete}[Loi de décroissance du marquage (admise au Probatoire)]
Dans un système au \textbf{steady-state} (état stationnaire : masse cellulaire constante, synthèse = dégradation), si l'on marque instantanément une cohorte de molécules à $t=0$ (pulse parfait) et qu'on suit son destin sans remise en circulation du traceur (chase parfait), la quantité de molécules marquées $M(t)$ décroît \textbf{exponentiellement} :
\[ M(t) = M_0 \cdot e^{-k_d t} \]
où $k_d$ est la constante de vitesse de dégradation (renouvellement). La demi-vie est liée à $k_d$ par $t_{1/2} = \frac{\ln 2}{k_d}$.
\textbf{Attention :} Cette loi suppose que la probabilité de dégradation est \emph{indépendante de l'âge} de la molécule (loi du hasard, cinétique du premier ordre), ce qui est vérifié pour la majorité des protéines solubles mais pas pour toutes (ex: hémoglobine dans l'érythrocyte qui vieillit de façon programmée). \textbf{Pour l'ADN}, il n'y a pas de dégradation/resynthèse globale de la double hélice : la perte de marquage reflète la \textbf{dilution par réplication semi-conservative} lors des divisions cellulaires (ou le turnover réparatif localisé).
\end{propriete}

\begin{exemplebox}[Exemple chiffré guidé : détermination de $t_{1/2}$]
On suit le marquage d'une protéine musculaire après pulse au \ce{^{3}H}-leucine.
Données : $t=0$ : 80\% ; $t=3$ j : 40\% ; $t=6$ j : 20\% ; $t=9$ j : 10\%.
\begin{enumerate}
    \item $M_0 = 80\%$.
    \item $M_0/2 = 40\%$.
    \item On lit le tableau : $M=40\%$ est atteint à $t=3$ jours.
    \item \textbf{Demi-vie $t_{1/2} = 3$ jours.}
    \item Vérification : à $t=6$ j ($2 t_{1/2}$), $M = 20\% = 80/4$ ; à $t=9$ j ($3 t_{1/2}$), $M = 10\% = 80/8$. La loi exponentielle est respectée.
    \item \textbf{Conclusion :} Cette protéine musculaire se renouvelle modérément (demi-vie de quelques jours), plus vite que le collagène, moins vite que les enzymes hépatiques.
\end{enumerate}
\end{exemplebox}

\courssub{Variabilité du renouvellement selon les tissus : l'exemple camerounais}

Le turnover n'est pas une constante universelle pour une molécule donnée ; il dépend du \textbf{type cellulaire} et de son \textbf{environnement physiologique}. Le tableau ci-dessous synthétise des valeurs typiques observées chez les mammifères (données transposables à l'homme et aux modèles d'élevage au Cameroun).

\begin{center}
\begin{tabularx}{\linewidth}{|l|X|X|X|}\hline
\rowcolor{popblueL}
\textbf{Tissu / Molécule} & \textbf{Demi-vie approximative ($t_{1/2}$)} & \textbf{Vitesse de renouvellement} & \textbf{Signification physiologique} \\ \hline
\textbf{Foie (enzymes solubles, albumine)} & 2 à 5 jours & \textbf{Très rapide} & Adaptation métabolique immédiate (jeûne/repas, détoxication). Le foie est l'usine chimique de l'organisme : il doit remplacer vite ses outils de travail. \\ \hline
\textbf{Muscle squelettique (protéines contractiles : actine, myosine)} & 10 à 20 jours & \textbf{Intermédiaire} & Plasticité musculaire : hypertrophie à l'effort (élevage bovin au Nord-Cameroun), atrophie à l'inactivité. Permet l'adaptation de la masse musculaire. \\ \hline
\textbf{Intestin (cellules de la muqueuse)} & 2 à 5 jours (turnover cellulaire complet) & \textbf{Extrêmement rapide} & Barrière face au lumen agressif (enzymes, bactéries, toxines). Renouvellement par division des cellules souches des cryptes. \\ \hline
\textbf{Peau (kératinocytes)} & 3 à 4 semaines & \textbf{Rapide} & Barrière physique, cicatrisation. \\ \hline
\textbf{Sang (hémoglobine dans globules rouges)} & $\sim$ 120 jours (durée de vie du GR) & \textbf{Programmé} & Pas de renouvellement moléculaire \emph{intra}-cellulaire (pas de noyau, pas de synthèse). Remplacement \emph{cellulaire} par érythropoïèse (moelle osseuse). \\ \hline
\textbf{Tendons / Cartilage / Os (collagène type I)} & Mois à années (collagène mature : $t_{1/2} > 1$ an) & \textbf{Très lent} & Fonction mécanique statique (résistance traction). Remplacement risqué (perte de solidité transitoire). Pathologie : tendinites chroniques, lenteur consolidation osseuse chez le sujet âgé. \\ \hline
\textbf{Cristallin (cristallines)} & \textbf{Nul} (durée de vie = vie de l'individu) & \textbf{Absent} & Transparence absolue requise. Pas de synthèse après la naissance. Dommages cumulés (UV, oxydation) $\rightarrow$ cataracte. \\ \hline
\textbf{ADN nucléaire (neurones, cardiomyocytes adultes)} & \textbf{Nul} (pas de division) & \textbf{Absent (sauf réparation)} & Stabilité du génome. Réparation excision (BER, NER) mais pas de réplication. \\ \hline
\end{tabularx}
\end{center}

\begin{savaistu}[Le turnover au service de la médecine nucléaire et de l'archéologie]
Le principe de la décroissance du marquage est exploité au quotidien :
\begin{itemize}
    \item \textbf{Scintigraphie osseuse (\ce{^{99m}Tc}) :} On injecte un diphosphonate marqué (\ce{^{99m}Tc}). Il se fixe sur l'hydroxyapatite osseuse là où le \textbf{turnover osseux} est actif (fracture en consolidation, métastase, ostéomyélite). L'image révèle les zones de \emph{renouvellement accéléré}.
    \item \textbf{Datation au Carbone 14 (\ce{^{14}C}) :} C'est l'exact inverse du pulse-chase ! Le \ce{^{14}C} atmosphérique (traceur naturel constant) est incorporé par la photosynthèse dans les végétaux, puis dans la chaîne alimentaire. À la mort de l'organisme, le « pulse » s'arrête (plus d'incorporation). Le \ce{^{14}C} radioactif décroît avec sa demi-vie physique ($t_{1/2} = 5730$ ans). Mesurer le \ce{^{14}C} résiduel dans un os ou du charbon de bois (ex: vestiges archéologiques de Shum Laka ou sédiments anciens du lac Nyos) permet de dater l'échantillon.
\end{itemize}
\end{savaistu}

\courssub{Limites et précautions d'interprétation : le piège du « turnover net »}

C'est un point crucial, source fréquente d'erreurs au Probatoire.

\begin{attention}[Erreur fréquente : Confondre disparition du marquage et disparition de la molécule]
\begin{itemize}
    \item \textbf{Ce qu'on mesure :} La disparition du \textbf{signal isotopique} (radioactivité ou excès d'isotope stable).
    \item \textbf{Ce que cela signifie :} La molécule \emph{initialement marquée} a été dégradée (hydrolysée en acides aminés, nucléotides, acides gras).
    \item \textbf{Ce que cela NE signifie PAS :} Que la \emph{fonction} ou la \emph{masse} de cette molécule a disparu. Elle a été \textbf{remplacée à l'identique} par une molécule neuve, non marquée, synthétisée à partir du pool de précurseurs non marqués (chase).
    \item \textbf{Piège « Turnover net vs Turnover brut » :} Si la synthèse et la dégradation sont égales (steady-state), la masse totale est constante. Le marquage ne mesure que le \textbf{flux de dégradation} (qui égale le flux de synthèse). Si la masse augmente (croissance, hypertrophie), synthèse $>$ dégradation : le marquage sous-estime la synthèse totale. Si la masse diminue (atrophie, jeûne), dégradation $>$ synthèse : le marquage surestime la synthèse nette.
    \item \textbf{Recyclage du traceur :} Chez l'animal entier, les acides aminés libérés par la dégradation des protéines marquées peuvent être \emph{réutilisés} pour synthétiser de nouvelles protéines (réincorporation du \ce{^{14}C}). Cela \textbf{allonge artificiellement} la demi-vie apparente mesurée. En culture cellulaire (milieu renouvelé), ce biais est éliminé.
\end{itemize}
\end{attention}

\courssub{Exercice résolu : Exploitation complète d'un protocole pulse-chase}

\begin{exoresolu}[Analyse du renouvellement de l'ADN hépatique vs ADN cérébral chez la souris]
\textbf{Énoncé :} On injecte de la thymidine-\ce{^{3}H} (précurseur spécifique de l'ADN) à des souris adultes (pulse court). On sacrifie des lots d'animaux à J0, J7, J30, J90, J180. On isole les noyaux du foie et du cerveau, on extrait l'ADN et on mesure la radioactivité (cpm/$\mu$g ADN). Résultats :

\begin{center}
\begin{tabular}{|c|c|c|}\hline
\rowcolor{popblueL}
\textbf{Temps} & \textbf{Foie (cpm/$\mu$g)} & \textbf{Cerveau (cpm/$\mu$g)} \\ \hline
J0 & 10 000 & 10 000 \\ \hline
J7 & 5 200 & 9 800 \\ \hline
J30 & 1 300 & 9 500 \\ \hline
J90 & 160 & 9 200 \\ \hline
J180 & 20 & 9 000 \\ \hline
\end{tabular}
\end{center}

\textbf{Questions :}
\begin{enumerate}
    \item Représenter graphiquement l'évolution du marquage pour les deux tissus (choisir l'échelle adaptée).
    \item Déterminer la demi-vie de l'ADN hépatique.
    \item Interpréter la différence entre foie et cerveau au niveau cellulaire.
    \item Pourquoi a-t-on utilisé de la thymidine marquée et non de la leucine marquée ?
\end{enumerate}

\tcblower

\textbf{Solution détaillée :}

\begin{enumerate}
    \item \textbf{Graphique :} Les valeurs hépatiques chutent de 10 000 à 20 (facteur 500) : échelle \textbf{semi-logarithmique} indispensable. Les valeurs cérébrales varient peu (10 000 $\to$ 9 000) : échelle \textbf{linéaire} possible, mais pour comparer sur le même graphique, on utilise le semi-log pour les deux.
    \begin{center}
    \begin{tikzpicture}
    \begin{semilogyaxis}[
        width=0.9\linewidth, height=0.5\linewidth,
        xlabel={Temps (jours)}, ylabel={Radioactivité (cpm/$\mu$g ADN)},
        grid=both, xmin=0, xmax=180, ymin=1, ymax=12000,
        legend pos=south west,
    ]
    \addplot[popred, mark=*, thick] coordinates {(0,10000) (7,5200) (30,1300) (90,160) (180,20)};
    \addlegendentry{Foie}
    \addplot[popblue, mark=square*, thick] coordinates {(0,10000) (7,9800) (30,9500) (90,9200) (180,9000)};
    \addlegendentry{Cerveau}
    \end{semilogyaxis}
    \end{tikzpicture}
    \end{center}

    \item \textbf{Demi-vie ADN hépatique :} $M_0 = 10\,000$. $M_0/2 = 5\,000$.
    Lecture sur la courbe (ou interpolation J0-J7) : 5 200 à J7 $\approx$ 5 000.
    \textbf{$t_{1/2} \approx 7$ jours.}
    (Vérification : J30 $\approx$ 4,3 $t_{1/2}$ $\rightarrow$ $10\,000 / 2^{4,3} \approx 10\,000 / 20 = 500$. Mesure 1 300. Ordre de grandeur cohérent, légère réincorporation ou population hétérogène).

    \item \textbf{Interprétation cellulaire :}
    \begin{itemize}
        \item \textbf{Foie :} Marquage divisé par 2 en 7 jours. L'ADN nucléaire ne se renouvelle pas par turnover moléculaire (pas de dégradation/resynthèse de la double hélice). La perte de marquage reflète la \textbf{dilution par division cellulaire (mitose)} : chaque division partage les brins marqués entre 2 cellules filles. Une demi-vie de 7 jours indique un \textbf{renouvellement cellulaire actif} (prolifération des hépatocytes) ou une forte activité de \textbf{réparation de l'ADN} (turnover nucléotidique par excision, non détectable comme dilution globale mais possible). Chez la souris adulte, le foie a un turnover cellulaire lent (mois), cette demi-vie courte suggère souvent une expérience sur animal jeune en croissance ou une réincorporation significative du \ce{^{3}H} libéré.
        \item \textbf{Cerveau :} Marquage stable ($\approx$ 95-100\% conservé sur 6 mois). Les neurones sont en \textbf{$G_0$ terminal (post-mitotiques)}. Pas de division $\rightarrow$ pas de dilution du marquage. L'ADN est \textbf{stable}, seul le turnover de réparation (faible) existe, invisible à cette échelle.
    \end{itemize}

    \item \textbf{Choix du traceur :} La \textbf{thymidine} (désoxyribonucléoside) est un précurseur \textbf{spécifique de l'ADN} (incorporée uniquement après phosphorylation en dTTP par la thymidine kinase, active en phase S). La leucine est incorporée dans les \textbf{protéines}. Pour étudier l'ADN, il faut un traceur de l'ADN.
\end{enumerate}
\end{exoresolu}

\begin{aretenir}[Synthèse : Ce que le marquage isotopique nous apprend]
\begin{itemize}
    \item Le \textbf{pulse-chase} est la méthode de référence pour quantifier le \textbf{turnover moléculaire} in vivo.
    \item La courbe de décroissance du marquage est \textbf{exponentielle} (loi du premier ordre) $\rightarrow$ caractérisation par une \textbf{demi-vie ($t_{1/2}$)}.
    \item $t_{1/2}$ court = renouvellement rapide (adaptation, régulation, défense) ; $t_{1/2}$ long = stabilité structurelle, économie d'énergie.
    \item \textbf{Le marquage suit la cohorte, pas l'individu :} Il mesure le \textbf{flux de remplacement net}, pas le destin d'une molécule unique.
    \item \textbf{Application médicale :} Imagerie du turnover osseux (scintigraphie \ce{^{99m}Tc}), datation archéologique (\ce{^{14}C}), diagnostic de tumeurs (marquage thymidine/\ce{^{18}F}-FDG pour prolifération).
\end{itemize}
\end{aretenir}

\coursec{Renouvellement moléculaire et intégrité cellulaire}

Dans la section précédente, l'exploitation des données de marquage isotopique (azote $^{15}$N, carbone $^{14}$C, acides aminés tritiés) nous a permis de démontrer un fait capital : les molécules de la cellule sont \emph{perpétuellement détruites et reconstruites}, même lorsque la cellule ne grandit pas et ne se divise pas. Une question surgit alors naturellement : \textbf{à quoi sert ce gaspillage apparent d'énergie ?} Pourquoi la cellule dépense-t-elle de l'ATP à fabriquer sans cesse des protéines et des lipides qui semblent déjà présents ? C'est à cette question que répond la présente section : le renouvellement moléculaire n'est pas un luxe, c'est la \cle{condition du maintien de l'intégrité structurale et fonctionnelle de la cellule}.

\courssub{Une observation pour commencer : pourquoi une blessure de la peau finit-elle par disparaître ?}

Un élève de Yaoundé se brûle légèrement la main en cuisinant. La peau rougit, une petite zone est endommagée. Quelques jours plus tard, la peau est redevenue normale : les cellules endommagées ont été réparées ou remplacées, et \emph{à l'intérieur même des cellules survivantes}, les protéines coagulées par la chaleur ont disparu et ont été remplacées par des protéines neuves et fonctionnelles.

De même, les sportifs camerounais qui s'entraînent sous le soleil de la saison sèche voient leur peau s'adapter ; les personnes qui se remettent d'une jaunisse voient leur foie retrouver son fonctionnement normal. Ces guérisons quotidiennes reposent sur un principe unique : \textbf{la cellule ne se contente pas de posséder des molécules, elle les surveille, élimine celles qui sont abîmées et les remplace}. Sans ce renouvellement, les molécules défectueuses s'accumuleraient, les fonctions se dégraderaient, et la cellule finirait par mourir.

\courssub{Le renouvellement élimine les molécules endommagées avant qu'elles ne perturbent la cellule}

\begin{definition}[Molécule endommagée]
Une \cle{molécule endommagée} (ou altérée) est une molécule dont la structure a été modifiée par des agressions physico-chimiques (chaleur, rayonnements UV, radicaux libres de l'oxygène, produits toxiques) de sorte qu'elle ne peut plus assurer correctement sa fonction. On distingue principalement :
\begin{itemize}
\item les \cle{protéines dénaturées} : leur repliement spatial (structure tridimensionnelle) est détruit, donc leur activité (enzymatique, structurale, de transport) est perdue ;
\item les \cle{lipides oxydés} : les acides gras insaturés des membranes sont attaqués par les radicaux libres (peroxydation lipidique), ce qui rigidifie et fragilise les membranes ;
\item les \cle{lésions de l'ADN} : cassures de brins, dimères de thymine provoqués par les UV, bases modifiées par des substances toxiques.
\end{itemize}
\end{definition}

La logique est celle d'une maison bien entretenue : si l'on ne remplace jamais les tôles rouillées ni les poutres attaquées par les termites, la maison finit par s'effondrer. De même, la cellule identifie en permanence ses molécules « rouillées » et les remplace.

\begin{propriete}[Propriété essentielle — à retenir et savoir argumenter]
Le \cle{renouvellement moléculaire} est la condition du maintien de l'\cle{intégrité structurale} (forme, membranes, organites) et de l'\cle{intégrité fonctionnelle} (activités enzymatiques, échanges, expression du patrimoine génétique) de la cellule. En remplaçant sans cesse les molécules usées ou endommagées par des molécules neuves, la cellule préserve ses structures et ses fonctions malgré les agressions permanentes du milieu et les erreurs de son propre métabolisme. \textbf{Tout ralentissement durable de ce renouvellement conduit à l'accumulation de molécules défectueuses, aux dysfonctionnements, puis à la mort cellulaire.}
\end{propriete}

\begin{methode}[Relier renouvellement moléculaire et intégrité cellulaire]
Face à un document ou à une question d'argumentation (type « montrez l'importance du renouvellement moléculaire »), procède toujours en trois étapes :
\begin{enumerate}
\item \textbf{Identifier la molécule altérée} : de quelle molécule s'agit-il (protéine, lipide membranaire, ADN) et quelle agression l'a endommagée (chaleur, UV, radicaux libres, toxique) ?
\item \textbf{Expliquer son remplacement} : la molécule défectueuse est dégradée (protéasome pour les protéines, lysosomes, systèmes d'excision pour l'ADN) et une molécule neuve est synthétisée (traduction pour les protéines, néosynthèse lipidique, réplication/réparation pour l'ADN).
\item \textbf{Conclure sur la fonction préservée} : grâce à ce remplacement, telle fonction (catalyse, perméabilité sélective, fidélité de l'information génétique) est maintenue, donc la cellule reste vivante et fonctionnelle.
\end{enumerate}
Cette démarche en trois temps « molécule altérée → remplacement → fonction préservée » est exactement celle attendue au Probatoire dans les questions d'argumentation.
\end{methode}

\courssub{Maintien des membranes : le renouvellement des lipides garantit fluidité et perméabilité sélective}

Toutes les membranes cellulaires (membrane plasmique, membranes du réticulum, des mitochondries, du noyau) sont des \cle{bicouches phospholipidiques} dans lesquelles flottent des protéines. Leur bon fonctionnement exige deux qualités :

\begin{itemize}
\item la \cle{fluidité membranaire} : les lipides doivent pouvoir se déplacer latéralement ; une membrane trop rigide ne permet ni les mouvements des protéines de transport, ni les déformations (endocytose, division cellulaire) ;
\item la \cle{perméabilité sélective} : la membrane contrôle les entrées et sorties de la cellule grâce à ses protéines (canaux, transporteurs, pompes).
\end{itemize}

Or les acides gras insaturés des phospholipides sont des cibles privilégiées des \cle{radicaux libres} produits par la respiration cellulaire elle-même. La peroxydation transforme ces lipides en lipides oxydés qui rigidifient la membrane et perturbent les protéines qu'elle contient. Le renouvellement permanent des lipides membranaires — dégradation des lipides oxydés et insertion de lipides néosynthétisés — maintient en permanence la fluidité et donc la perméabilité sélective. Les protéines membranaires, elles aussi, sont renouvelées : une pompe Na$^+$/K$^+$ dénaturée est remplacée, ce qui préserve les gradients ioniques indispensables à la vie de la cellule (potentiel de membrane, entrée du glucose).

\courssub{Réparation permanente de l'ADN : préserver l'information génétique}

L'ADN est la molécule la plus précieuse de la cellule : il porte l'\cle{information génétique}. Sa \emph{séquence} (la succession des nucléotides) est remarquablement stable, mais la \emph{molécule} d'ADN, elle, subit en permanence des agressions :

\begin{itemize}
\item les \textbf{rayons ultraviolets} du soleil provoquent la formation de \cle{dimères de thymine} (deux thymines voisines se lient anormalement), ce qui bloque la réplication et la transcription ;
\item les \textbf{produits toxiques} (goudrons du tabac, aflatoxines des arachides et maïs mal conservés, certains colorants) modifient chimiquement les bases ;
\item les \textbf{radicaux libres} et même l'eau provoquent des cassures et des pertes de bases (on estime des milliers de lésions spontanées par jour et par cellule).
\end{itemize}

La cellule possède des \cle{enzymes de réparation} qui patrouillent le long de l'ADN, détectent les lésions, excisent le segment défectueux et resynthétisent un segment correct en utilisant le brin complémentaire comme matrice. Ce mécanisme fait partie intégrante du renouvellement moléculaire : des fragments d'ADN sont constamment remplacés. Grâce à lui, l'information génétique reste fidèle et les mutations transmissibles restent rares.

\begin{exemplebox}[Les cellules de la peau et le soleil : réparer ou dégénérer]
Les kératinocytes de la peau d'un habitant de Maroua ou de Ngaoundéré, exposés chaque jour à un ensoleillement intense, subissent en permanence des dimères de thymine dans leur ADN. Chez une personne saine, les enzymes de réparation corrigent ces lésions au fur et à mesure : la peau reste saine malgré des années d'exposition. Chez les malades atteints de \emph{xeroderma pigmentosum} (maladie génétique où une enzyme de réparation est déficiente), les lésions s'accumulent, les mutations se multiplient, et des \textbf{cancers cutanés} apparaissent dès l'enfance. Cet exemple montre de manière frappante que la réparation permanente de l'ADN — une forme de renouvellement moléculaire — est une condition de survie : quand elle fait défaut, l'intégrité de l'information génétique est perdue et la cellule peut devenir cancéreuse.
\end{exemplebox}

\begin{attention}[Erreur fréquente : « l'ADN ne change jamais »]
Beaucoup d'élèves écrivent que « l'ADN est stable, donc il n'est pas concerné par le renouvellement moléculaire ». C'est une confusion entre \emph{séquence} et \emph{molécule}. La \textbf{séquence} des nucléotides (le message génétique) est effectivement stable et transmise fidèlement ; mais la \textbf{molécule} d'ADN subit en permanence des lésions (cassures, dimères, bases oxydées) qui sont réparées en continu : des nucléotides sont excisés et remplacés chaque jour. Le marquage isotopique le confirme : des nucléotides marqués s'incorporent dans l'ADN de cellules qui ne se divisent pas, preuve d'un renouvellement lié à la réparation. Ne jamais écrire au Probatoire que « l'ADN ne se renouvelle pas ».
\end{attention}

\courssub{Adaptation cellulaire : ajuster les quantités d'enzymes aux besoins}

Le renouvellement moléculaire n'est pas seulement un mécanisme de maintenance : c'est aussi un outil d'\cle{adaptation}. Puisque chaque protéine a une durée de vie limitée, la cellule peut moduler la \emph{quantité} de chaque enzyme en jouant sur l'équilibre synthèse/dégradation :

\begin{itemize}
\item après un \textbf{repas riche en protéines} (viande de bœuf, poisson fumé, niébé), le foie augmente la synthèse des enzymes de dégradation des acides aminés (transaminases, enzymes du cycle de l'urée) ; comme les anciennes enzymes sont continuellement dégradées, la proportion d'enzymes adaptées au nouveau régime augmente rapidement ;
\item après un \textbf{repas riche en glucides} (foufou de manioc, igname), les enzymes de la glycolyse et du stockage du glycogène sont davantage synthétisées ;
\item lors d'un \textbf{jeûne prolongé}, au contraire, les enzymes de la néoglucogenèse deviennent majoritaires.
\end{itemize}

Sans renouvellement, la cellule resterait figée avec son stock initial d'enzymes et serait incapable de s'adapter aux variations de l'alimentation ou de l'environnement. Le renouvellement confère donc à la cellule une \cle{plasticité fonctionnelle} indispensable à l'homéostasie de l'organisme.

\courssub{Conséquences d'un défaut de renouvellement : de l'accumulation à la mort cellulaire}

Que se passe-t-il si le renouvellement est bloqué ou ralenti ? La logique du phénomène se déroule en une chaîne de conséquences :

\begin{enumerate}
\item les molécules endommagées (protéines dénaturées, lipides oxydés, ADN lésé) ne sont plus éliminées : elles \textbf{s'accumulent} ;
\item les protéines défectueuses forment des \textbf{agrégats} qui encombrent le cytoplasme et perturbent les organites ; les membranes perdent leur fluidité ; les mutations non réparées s'additionnent ;
\item les fonctions cellulaires se \textbf{dégradent} : activités enzymatiques diminuées, échanges membranaires perturbés, expression génétique erronée ;
\item au-delà d'un seuil, la cellule devient inviable : c'est la \textbf{mort cellulaire}, ou pire, sa transformation cancéreuse si l'ADN lésé contrôle la division.
\end{enumerate}

\begin{popfigure}
\begin{center}
\begin{tikzpicture}[scale=0.92, every node/.style={font=\small}]
% ===== Cellule saine =====
\draw[thick, popgreen!70!black, fill=popgreenL] (0,0) ellipse (3.1 and 2.3);
\draw[thick, popgreen!70!black, fill=popgreen!25] (0,0.2) circle (0.75);
\node[font=\small\bfseries] at (0,0.2) {Noyau};
\node[font=\footnotesize, popgreen!50!black] at (0,-0.35) {ADN r\'epar\'e};
% molécules saines (petits hexagones réguliers)
\foreach \p in {(-1.9,1.1),(-1.2,-1.2),(1.4,1.3),(1.9,-0.9),(-2.2,-0.3),(2.3,0.4),(0.9,-1.6),(-0.9,1.7)}{
  \node[regular polygon, regular polygon sides=6, draw=popgreen!60!black, fill=popgreen!50, inner sep=1.6pt] at \p {};
}
% protéasome
\draw[thick, popblue, fill=popblueL] (-1.6,0.35) rectangle ++(0.5,0.9);
\node[font=\footnotesize, popblue!70!black] at (-1.35,1.55) {Prot\'easome};
% lysosome
\draw[thick, poppurple, fill=poppurpleL] (1.7,-1.5) circle (0.32);
\node[font=\footnotesize, poppurple!70!black, align=center] at (2.6,-1.85) {Lysosome};
% flèche recyclage
\draw[->, thick, popgreen!60!black] (-1.1,0.8) .. controls (-0.4,1.6) .. (0.6,1.5) node[midway, above, font=\footnotesize, align=center]{acides amin\'es\\recycl\'es};
\node[font=\bfseries, popgreen!50!black] at (0,-2.85) {Cellule \`a renouvellement efficace};
\node[font=\footnotesize, align=center] at (0,-3.45) {mol\'ecules saines, membrane fluide,\\ fonctions pr\'eserv\'ees};

% ===== Cellule bloquée =====
\begin{scope}[xshift=8.6cm]
\draw[thick, poporange!80!black, fill=poporangeL] (0,0) ellipse (3.1 and 2.3);
\draw[thick, poporange!80!black, fill=poporange!30] (0,0.2) circle (0.75);
\node[font=\small\bfseries] at (0,0.2) {Noyau};
\node[font=\footnotesize, poporange!80!black] at (0,-0.35) {ADN l\'es\'e};
% molécules défectueuses (formes irrégulières / étoiles)
\foreach \p in {(-1.9,1.1),(-1.2,-1.2),(1.4,1.3),(1.9,-0.9),(-2.2,-0.3),(2.3,0.4),(0.9,-1.6),(-0.9,1.7)}{
  \node[star, star points=5, draw=poporange!90!black, fill=poporange!70, inner sep=1.4pt] at \p {};
}
% agrégat
\draw[thick, poppink!80!black, fill=poppinkL, decorate, decoration={coil, segment length=3pt, amplitude=1.5pt}] (-1.5,0.5) circle (0.55);
\node[font=\footnotesize, poppink!80!black, align=center] at (-2.5,1.5) {Agr\'egat de\\prot\'eines\\d\'enatur\'ees};
% membrane fragilisée
\draw[thick, poppink!80!black] (2.6,1.5) -- (3.05,1.15);
\node[font=\footnotesize, poppink!80!black, align=center] at (2.9,1.85) {membrane\\rigidifi\'ee};
\node[font=\bfseries, poporange!80!black] at (0,-2.85) {Cellule \`a renouvellement bloqu\'e};
\node[font=\footnotesize, align=center] at (0,-3.45) {accumulation de mol\'ecules d\'efectueuses,\\ dysfonctionnements, mort possible};
\end{scope}
\end{tikzpicture}
\end{center}
\legende{Comparaison de deux cellules. À gauche, le renouvellement est efficace : le protéasome et les lysosomes dégradent les molécules usées, les acides aminés sont recyclés, les molécules (hexagones réguliers) sont saines, l'ADN est réparé. À droite, le renouvellement est bloqué : les molécules défectueuses (étoiles) s'accumulent, des agrégats de protéines dénaturées encombrent le cytoplasme, la membrane se rigidifie, l'ADN reste lésé : la cellule se dirige vers la mort.}
\end{popfigure}

\begin{savaistu}[Agrégats protéiques et maladies neurodégénératives]
La maladie d'Alzheimer et la maladie de Parkinson illustrent dramatiquement les conséquences d'un défaut d'élimination des protéines défectueuses : des protéines mal repliées (peptide bêta-amyloïde, alpha-synucléine) s'accumulent en agrégats dans les neurones, qui finissent par mourir. Avec l'allongement de l'espérance de vie au Cameroun, ces maladies concernent de plus en plus de familles : comprendre le renouvellement moléculaire, c'est aussi comprendre pourquoi la recherche biomédicale tente de stimuler les systèmes de dégradation cellulaire (protéasome, autophagie) pour traiter ces pathologies.
\end{savaistu}

\begin{exoresolu}[Argumenter le rôle du renouvellement moléculaire]
\textbf{Énoncé.} On cultive deux lots de cellules hépatiques (cellules du foie). Au lot A, on ajoute un inhibiteur du protéasome (qui bloque la dégradation des protéines) ; le lot B sert de témoin. Après 48 heures, on observe dans le lot A une accumulation de protéines anormales, une diminution de l'activité des enzymes respiratoires et 30 \% de cellules mortes, alors que le lot B est normal. Expliquez ces résultats et concluez sur le rôle du renouvellement moléculaire dans le maintien de l'intégrité cellulaire.
\tcblower
\textbf{Solution détaillée.}

\emph{Étape 1 — Identifier la molécule altérée.} Dans toute cellule vivante, des protéines sont continuellement dénaturées (chaleur métabolique, radicaux libres issus de la respiration). Normalement, ces protéines défectueuses sont dégradées par le \cle{protéasome}.

\emph{Étape 2 — Expliquer le défaut de remplacement.} Dans le lot A, l'inhibiteur bloque le protéasome : les protéines dénaturées ne sont plus éliminées. Elles s'accumulent dans le cytoplasme et forment des agrégats. La synthèse de protéines neuves continue, mais l'équilibre synthèse/dégradation est rompu : la proportion de protéines non fonctionnelles augmente sans cesse.

\emph{Étape 3 — Conclure sur la fonction perdue.} Les enzymes respiratoires endommagées n'étant plus remplacées efficacement, l'activité respiratoire chute ; la production d'ATP diminue ; privées d'énergie et encombrées d'agrégats, 30 \% des cellules meurent. Le lot témoin B, dont le protéasome fonctionne, élimine ses protéines défectueuses au fur et à mesure et conserve son intégrité.

\emph{Conclusion générale.} Cette expérience démontre que le renouvellement moléculaire — ici la dégradation des protéines usées couplée à leur resynthèse — est \textbf{indispensable au maintien de l'intégrité structurale et fonctionnelle de la cellule} : lorsqu'il est bloqué, les molécules défectueuses s'accumulent, les fonctions s'effondrent et la cellule meurt.
\end{exoresolu}

\begin{aretenir}[L'essentiel de la section]
\begin{itemize}
\item La cellule subit en permanence des agressions qui produisent des \cle{protéines dénaturées}, des \cle{lipides oxydés} et des \cle{lésions de l'ADN}.
\item Le \cle{renouvellement moléculaire} élimine ces molécules défectueuses (protéasome, lysosomes, excision de l'ADN) et les remplace par des molécules neuves : c'est la condition du maintien de l'\cle{intégrité structurale et fonctionnelle} de la cellule.
\item Le renouvellement des lipides et protéines membranaires garantit la \cle{fluidité} et la \cle{perméabilité sélective} des membranes.
\item La réparation permanente de l'ADN préserve l'\cle{information génétique} ; son défaut favorise les cancers (exemple des cancers cutanés).
\item Le renouvellement permet l'\cle{adaptation} : la cellule ajuste ses quantités d'enzymes aux besoins (après un repas, lors d'un jeûne).
\item Un défaut de renouvellement entraîne : accumulation de molécules défectueuses $\rightarrow$ dysfonctionnements $\rightarrow$ mort cellulaire (ou cancérisation).
\item Argumentation type en trois temps : \textbf{molécule altérée $\rightarrow$ remplacement $\rightarrow$ fonction préservée}.
\end{itemize}
\end{aretenir}

\begin{exercice}[Pour t'entraîner]
Des cellules de peau sont exposées aux UV puis cultivées. Chez des cellules normales, la radioactivité d'un précurseur marqué de l'ADN augmente transitoirement même en l'absence de division cellulaire ; chez des cellules d'un malade atteint de xeroderma pigmentosum, cette incorporation est quasi nulle et des mutations s'accumulent. Interprète ces résultats en utilisant la démarche « molécule altérée → remplacement → fonction préservée », et explique pourquoi ce malade doit se protéger strictement du soleil.
\end{exercice}

\begin{corrige}[Exercice]
Les UV provoquent des dimères de thymine : la molécule altérée est l'\textbf{ADN}. Chez les cellules normales, l'incorporation transitoire du précurseur marqué, sans division, traduit la \textbf{réparation} : les segments lésés sont excisés et resynthétisés (renouvellement local de l'ADN), ce qui préserve la fidélité de l'information génétique. Chez le malade, l'enzyme de réparation est déficiente : pas d'incorporation, les lésions ne sont pas corrigées, les mutations s'accumulent. La fonction préservée chez le sujet sain — l'intégrité du message génétique — est perdue chez le malade : ses cellules cutanées risquent la transformation cancéreuse. Il doit donc limiter strictement son exposition aux UV (vêtements couvrants, chapeau, crème solaire, horaires d'ensoleillement réduits) pour réduire le nombre de lésions que ses cellules ne peuvent pas réparer.
\end{corrige}

\coursec{Ralentissement du renouvellement et vieillissement cellulaire}

Dans la section précédente, nous avons vu que le renouvellement moléculaire permanent est la condition \emph{sine qua non} du maintien de l'intégrité cellulaire : synthèse et dégradation s'équilibrent pour remplacer les molécules abîmées avant qu'elles ne nuisent. Mais cet équilibre dynamique n'est pas figé dans le temps. Observons ce qui se passe lorsque l'horloge biologique avance.

\courssub{Constat expérimental : le turnover ralentit avec l'âge}

De nombreuses études, notamment par marquage isotopique (carbone 13, azote 15, eau lourde \ce{^2H2O}) chez l'humain et les mammifères, montrent une tendance universelle : la vitesse de synthèse des macromolécules diminue progressivement avec l'âge, tout comme la capacité de dégradation (protéasome, autophagie). Le \cle{turnover} global — le rapport entre la quantité renouvelée par unité de temps et le stock total — suit une courbe décroissante.

\begin{popfigure}
\begin{tikzpicture}
\begin{axis}[
    width=\linewidth,
    height=8cm,
    xlabel={Âge (années)},
    ylabel={Vitesse de renouvellement des protéines (\% du stock/jour)},
    domain=0:80,
    samples=100,
    xmin=0, xmax=80,
    ymin=0, ymax=6,
    axis lines=left,
    tick style={popdark},
    xlabel style={anchor=north, yshift=-5pt},
    ylabel style={anchor=south, xshift=-5pt},
    grid=major,
    grid style={popline},
]
% Courbe théorique décroissante exponentielle type
\addplot[popcurve, thick, domain=0:80] {5*exp(-x/30) + 0.5};
\addlegendentry{Tendance générale}
% Points indicatifs
\addplot[only marks, mark=*, mark size=2pt, poporange] coordinates {
    (10, 4.5)
    (30, 2.8)
    (50, 1.6)
    (70, 0.9)
};
\addlegendentry{Données expérimentales moyennes}
\end{axis}
\end{tikzpicture}
\legende{Évolution de la vitesse de renouvellement des protéines en fonction de l'âge. La courbe montre une décroissance exponentielle : le turnover est très actif pendant la croissance (enfance, adolescence), puis ralentit continuellement à l'âge adulte. Les points représentent des moyennes issues d'études de marquage isotopique (ex. eau lourde \ce{^2H2O}).}
\end{popfigure}

\begin{aretenir}[Constat clé]
Avec l'âge, les vitesses de synthèse (transcription, traduction) et de dégradation (protéasome, lysosome, autophagie) diminuent de concert. Le turnover moléculaire devient moins efficace : les molécules endommagées persistent plus longtemps avant d'être remplacées.
\end{aretenir}

\courssub{Conséquences moléculaires : accumulation de molécules défectueuses}

Lorsque la dégradation ne suit plus le rythme des dommages (stress oxydant, erreurs de réplication, chaleur, toxiques), et que la synthèse de neuf ne compense pas les pertes, trois types majeurs de dégâts s'accumulent :

\begin{itemize}
    \item \textbf{Protéines dénaturées et agrégats :} Les protéines mal repliées, normalement refoldées par les chaperonnes (Hsp70, Hsp90) ou dégradées par le protéasome, forment des agrégats insolubles (corps d'inclusion). Exemple : lipofuscine (pigment « d'usure »), agrégats d'$\alpha$-synucléine (maladie de Parkinson), plaques $\beta$-amyloïdes (Alzheimer).
    \item \textbf{Lipides oxydés :} Les membranes riches en acides gras polyinsaturés sont cibles des \cle{radicaux libres} (ROS : \ce{^.OH}, \ce{O2^{.-}}). La peroxydation lipidique produit du malondialdéhyde (MDA) et du 4-HNE, toxiques pour les protéines et l'ADN. Les membranes deviennent rigides, moins fluides, altérant les transporteurs et récepteurs.
    \item \textbf{Lésions d'ADN non réparées :} Brins cassés, bases oxydées (8-oxoguanine), adducts. Si les systèmes de réparation (BER, NER, HR, NHEJ) sont submergés ou moins expressifs, les mutations s'accumulent $\rightarrow$ instabilité génomique, risque cancéreux, dysfonction transcriptionnelle.
\end{itemize}

\begin{definition}[Vieillissement cellulaire (sénescence réplicative)]
État physiologique irréversible dans lequel une cellule, bien que métaboliquement active, a perdu sa capacité de prolifération et présente un phénotype altéré (morphologie agrandie, $\beta$-galactosidase active à pH 6, sécrétion de facteurs pro-inflammatoires : SASP \emph{Senescence-Associated Secretory Phenotype}). Elle résulte de l'incapacité du renouvellement moléculaire à maintenir l'intégrité du génome, du protéome et du lipidome.
\end{definition}

\courssub{Conséquences cellulaires : dysfonctionnement organellaire et entrée en sénescence}

L'accumulation de molécules défectueuses frappe au cœur des organites :

\begin{itemize}
    \item \textbf{Mitochondries :} Leur ADN (ADNmt), non protégé par des histones et proche des sites de production de ROS, s'endommage vite. Mutations $\rightarrow$ sous-unités respiratoires défectueuses $\rightarrow$ fuite d'électrons $\rightarrow$ plus de ROS (cercle vicieux). L'autophagie spécifique (mitophagie) ralentit $\rightarrow$ accumulation de mitochondries dysfonctionnelles $\rightarrow$ chute de l'ATP, activation de l'apoptose.
    \item \textbf{Réticulum endoplasmique :} Stress ER par accumulation de protéines mal repliées $\rightarrow$ réponse UPR (\emph{Unfolded Protein Response}) chronique $\rightarrow$ apoptose si non résolu.
    \item \textbf{Noyau :} Instabilité génomique, raccourcissement des télomères (horloge mitotique), altération de l'épigénome (horloge épigénétique) $\rightarrow$ arrêt du cycle cellulaire (p53/p21, p16/Rb) $\rightarrow$ \textbf{sénescence}.
\end{itemize}

La cellule sénescente ne meurt pas tout de suite ; elle sécrète le SASP (cytokines IL-6, IL-8, métalloprotéinases, facteurs de croissance) qui propage l'inflammation et la sénescence aux cellules voisines (\emph{bystander effect}).

\courssub{Conséquences à l'échelle de l'organisme : signes visibles du vieillissement}

Le vieillissement de l'organisme est la somme, non linéaire, des vieillissements cellulaires dans les tissus.

\begin{exemplebox}[Retour à Mama Odile et Junior]
\textbf{Situation :} Mama Odile (68 ans) et son petit-fils Junior (12 ans) se font une égratignure identique au genou en jouant au football dans la cour de la maison à Yaoundé. Chez Junior, la croûte tombe au bout de 5 jours, la peau est rose puis normale. Chez Mama Odile, la plaie suinte encore au 10\textsuperscript{e} jour, la cicatrice reste rouge, fine, fragile pendant des semaines.

\textbf{Explication moléculaire :} La cicatrisation repose sur la prolifération des fibroblastes et la synthèse massive de \cle{collagène} (protéine structurelle de la matrice extracellulaire).
\begin{itemize}
    \item Chez Junior : turnover du collagène rapide (demi-vie courte), fibroblastes prolifératifs, angiogenèse efficace $\rightarrow$ réparation nette.
    \item Chez Mama Odile : turnover du collagène ralenti (demi-vie de plusieurs années, voir \emph{Le savais-tu ?}), fibroblastes sénescents (moins de division, SASP pro-inflammatoire), vascularisation réduite $\rightarrow$ matrice mal déposée, cicatrisation lente, cicatrice atrophique.
\end{itemize}
\end{exemplebox}

\begin{center}
\begin{tabularx}{\linewidth}{|l|X|X|}
\hline
\rowcolor{popblueL} \textbf{Fonction / Tissu} & \textbf{Conséquence du ralentissement du turnover} & \textbf{Manifestation clinique} \\
\hline
Peau (derme) & Collagène/élastine renouvelés lentement, fragmentation, glycation & Rides, perte d'élasticité, peau fine, cicatrisation lente \\
\hline
Muscle squelettique & Synthèse protéique (myofibrilles) $<$ dégradation, perte motoneurones & \cle{Sarcopénie} (fonte musculaire), faiblesse, risque de chute \\
\hline
Système immunitaire & Renouvellement lent des lymphocytes (thymus involué), répertoire TCR réduit & \cle{Immunosénescence} : infections plus graves, vaccins moins efficaces, cancer \\
\hline
Os & Ostéoblastes moins actifs, matrice osseuse moins renouvelée & Ostéoporose, fractures \\
\hline
Cerveau & Agrégats protéiques, perte synapses, myéline altérée & Déclin cognitif, maladies neurodégénératives \\
\hline
\end{tabularx}
\end{center}

\courssub{Facteurs influençant la vitesse du renouvellement moléculaire}

Le programme génétique donne la tendance, mais l'environnement module l'amplitude. C'est la \cle{plasticité phénotypique} du vieillissement.

\begin{propriete}[Facteurs modulateurs du turnover moléculaire (admis)]
\begin{itemize}
    \item \textbf{Alimentation équilibrée :} Apport suffisant en acides aminés essentiels (leucine $\rightarrow$ activation mTOR $\rightarrow$ synthèse protéique), acides gras oméga-3 (fluidité membranaire, résolution inflammation), micronutriments (Zn, Se, Mg pour enzymes réparation/antioxydantes).
    \item \textbf{Activité physique régulière :} Stimule l'autophagie, la biogenèse mitochondriale (PGC-1$\alpha$), la synthèse protéique musculaire, la sensibilité à l'insuline. C'est le plus puissant « médicament » anti-âge connu.
    \item \textbf{Restriction calorique (sans malnutrition) :} Allonge la durée de vie chez de nombreuses espèces (levures, vers, mouches, rongeurs, primates) en activant AMPK, sirtuines (SIRT1), FOXO $\rightarrow$ réparation ADN, autophagie, défense antioxydante.
    \item \textbf{Toxiques :} \textbf{Tabac} (radicaux libres, adduits ADN, inhibition protéasome), \textbf{Alcool} (stress ER, peroxydation lipidique, carence vitamines B), \textbf{Rayonnements UV/ionisants} (brisures ADN, ROS).
    \item \textbf{Sommeil :} Pic de sécrétion GH (hormone de croissance) et mélatonine (antioxydant) pendant le sommeil profond $\rightarrow$ réparation tissulaire, nettoyage métabolique (système glymphatique cérébral).
\end{itemize}
\end{propriete}

\begin{savaistu}[Radicaux libres, antioxydants et fruits de chez nous]
Les \cle{radicaux libres} (espèces réactives de l'oxygène, ROS) sont produits en permanence par la chaîne respiratoire mitochondriale (fuite 1-2\% des électrons). À dose physiologique, ils sont des messagers (signalisation redox). En excès (stress oxydant), ils attaquent tout : protéines, lipides, ADN.
La cellule possède des défenses \textbf{enzymatiques} (SOD, catalase, glutathion peroxydase — nécessitent Zn, Cu, Mn, Se) et \textbf{non enzymatiques} (glutathion, vitamine C, vitamine E, caroténoïdes, polyphénols).
Au Cameroun, nos marchés regorgent d'antioxydants naturels peu coûteux : \textbf{goavy} (goyave, record de vitamine C), \textbf{papaye}, \textbf{ananas}, \textbf{agrumes} (vit C), \textbf{arachides}, \textbf{huile de palme rouge} (vit E, caroténoïdes), \textbf{légumes-feuilles} (éru, folong, amarante — polyphénols, Mg). Une assiette colorée protège le renouvellement moléculaire mieux qu'aucun complément en gélule.
\end{savaistu}

\courssub{Chaîne causale : du ralentissement moléculaire au vieillissement de l'organisme}

\begin{popfigure}
\begin{tikzpicture}[
    node distance=1.2cm and 1.5cm,
    box/.style={rectangle, rounded corners, draw=popdark, thick, fill=popcream, text width=3.2cm, align=center, minimum height=1.8cm},
    arrow/.style={-Stealth, thick, popblue},
    link/.style={-Stealth, dashed, poporange}
]
% Nœuds
\node[box, align=center] (turnover){Ralentissement du turnover moléculaire\\ (synthèse $\downarrow$, dégradation $\downarrow$)};
\node[box, right=of turnover, align=center] (accum){Accumulation de molécules défectueuses\\ (protéines agrégées, lipides oxydés, ADN lésé)};
\node[box, right=of accum, align=center] (dysfonc){Dysfonctionnement organellaire\\ (mitochondries, RE, noyau) \\ \& Stress cellulaire chronique};
\node[box, below=of dysfonc, align=center] (senesc){Sénescence cellulaire\\ (arrêt cycle, SASP, apoptose)};
\node[box, left=of senesc, align=center] (tissu){Altération tissulaire\\ (perte cellules fonctionnelles, fibrosis, inflammation)};
\node[box, left=of tissu, align=center] (orga){Vieillissement de l'organisme\\ (sarcopénie, immunodépression, rides, maladies)};

% Flèches principales
\draw[arrow] (turnover) -- (accum);
\draw[arrow] (accum) -- (dysfonc);
\draw[arrow] (dysfonc) -- (senesc);
\draw[arrow] (senesc) -- (tissu);
\draw[arrow] (tissu) -- (orga);

% Boucle de rétroaction (SASP)
\draw[link] (senesc) to[out=90, in=90, looseness=1.5] node[above, font=\small, poporange] {SASP propage la sénescence} (dysfonc);
\draw[link] (orga) to[out=180, in=180, looseness=1.5] node[left, font=\small, poporange] {Environnement systémique inflammatoire (inflammaging)} (turnover);
\end{tikzpicture}
\legende{Chaîne causale reliant le ralentissement du renouvellement moléculaire au vieillissement de l'organisme. Les flèches pleines indiquent la causalité directe ; les flèches pointillées indiquent les boucles de rétroaction positive (aggravation) : le SASP des cellules sénescentes aggrave le stress des voisines, et l'inflammation systémique (\emph{inflammaging}) freine encore le turnover.}
\end{popfigure}

\courssub{Importance pour la santé : argumenter pour une longévité en bonne santé}

Comprendre ces mécanismes permet de passer d'une vision fataliste (« vieillir, c'est s'user ») à une vision \textbf{préventive et proactive} : maintenir le turnover le plus longtemps possible.

\begin{methode}[Argumenter sur l'importance du renouvellement moléculaire pour la santé]
\textbf{Étape 1 : Partir d'un fait observable (clinique, vie courante).} Exemple : « Les personnes âgées cicatrisent mal et font plus d'infections. »
\textbf{Étape 2 : Remonter au mécanisme cellulaire et moléculaire.} « La cicatrisation lente = turnover du collagène ralenti + fibroblastes sénescents. Infections = immunosénescence = turnover lymphocytaire ralenti + thymus involué. »
\textbf{Étape 3 : Identifier les leviers modulables.} « Alimentation riche en protéines/vitamines C/E $\rightarrow$ substrats synthèse + protection antioxydante. Activité physique $\rightarrow$ stimulation autophagie/synthèse musculaire. Éviter tabac/alcool $\rightarrow$ moins de dommages à réparer. »
\textbf{Étape 4 : Formuler un conseil de vie concret, réaliste, positif.} « Manger varié (légumes-feuilles, fruits, poisson/arachides), marcher 30 min/jour, dormir 7-8h, ne pas fumer : c'est investir dans son renouvellement moléculaire quotidien pour vieillir en autonomie. »
\end{methode}

\begin{exoresolu}[Exploiter des données de marquage pour argumenter]
\textbf{Énoncé :} Une étude compare le taux de synthèse fractionnaire (FSR, \emph{Fractional Synthesis Rate}) de la myofibrille musculaire chez des sujets jeunes (25 ans) et âgés (70 ans), au repos et après un exercice de résistance + apport de 20 g de protéines whey. Résultats (moyennes) :
\begin{center}
\begin{tabular}{|c|c|c|}
\hline
\rowcolor{popblueL} \textbf{Groupe} & \textbf{FSR au repos (\%/h)} & \textbf{FSR post-exercice + protéines (\%/h)} \\
\hline
Jeunes (25 ans) & 0,045 & 0,095 \\
\hline
Âgés (70 ans) & 0,030 & 0,055 \\
\hline
\end{tabular}
\end{center}
\textbf{Questions :}
1. Calculer le \% d'augmentation du FSR induit par l'exercice+protéines chez les jeunes et chez les âgés.
2. Interpréter la différence de réponse entre les deux groupes au niveau moléculaire (signalisation).
3. En déduire une recommandation pratique pour un senior.

\tcblower
\textbf{Solution détaillée :}
1. \textbf{Calculs :}
Jeunes : $\frac{0.095 - 0.045}{0.045} \times 100 = \frac{0.050}{0.045} \times 100 \approx \textbf{111\%}$ d'augmentation.
Âgés : $\frac{0.055 - 0.030}{0.030} \times 100 = \frac{0.025}{0.030} \times 100 \approx \textbf{83\%}$ d'augmentation.

2. \textbf{Interprétation moléculaire :} Chez les jeunes, l'exercice mécanique + la leucine (whey) activent puissamment la voie \textbf{mTORC1} (cible de la rapamycine chez les mammifères) $\rightarrow$ phosphorylation de 4E-BP1 et S6K1 $\rightarrow$ initiation de la traduction massive $\rightarrow$ synthèse protéique musculaire. Chez les âgés, on observe une \textbf{résistance anabolique} : le même stimulus active moins mTORC1 (inflammation basale, stress ER, lipidémie, sédentarité), d'où une réponse synthétique atténuée (83\% vs 111\%). Le turnover de base est aussi plus bas (0,030 vs 0,045 \%/h), signe d'un équilibre déplacé vers le catabolisme.

3. \textbf{Recommandation :} « Même si la réponse est moins forte, elle reste significative (+83\% !). Un senior \textbf{doit} combiner exercice de résistance (port de charges, exercices au poids du corps) \textbf{ET} apport protéique de qualité (20-30 g, riche en leucine : œuf, poisson, soja, lait) \textbf{à chaque repas principal} (3 fois/jour) pour stimuler le turnover musculaire au maximum de sa capacité résiduelle et lutter contre la sarcopénie. »
\end{exoresolu}

\begin{attention}[Pièges à éviter au Probatoire]
\begin{itemize}
    \item \textbf{Erreur majeure :} Écrire « Les cellules vieillissent parce qu'elles s'usent » ou « Les molécules s'usent ». \textbf{Non.} Les molécules ne s'usent pas comme des pièces mécaniques ; elles sont \textbf{endommagées} (oxydation, erreurs) et \textbf{mal remplacées} car le turnover ralentit. Le vieillissement est un \textbf{défaut de maintenance active}, pas une usure passive.
    \item \textbf{Confusion :} Sénescence cellulaire $\neq$ mort cellulaire (apoptose/nécrose). La cellule sénescente est vivante, métabolique, mais ne divise plus et sécrète des facteurs délétères.
    \item \textbf{Raccourci abusif :} « Les antioxydants arrêtent le vieillissement. » Faux. Les antioxydants \emph{limitent les dommages} (ralentissent l'accumulation), ils ne \emph{rétablissent pas} le turnover ralenti. L'exercice et l'apport protéique sont plus puissants pour \emph{relancer} la synthèse.
    \item \textbf{Oubli de l'échelle :} Ne pas relier l'échelle moléculaire (protéine) à l'échelle organisme (sarcopénie, cicatrisation). Le programme exige la \textbf{mise en relation des échelles}.
\end{itemize}
\end{attention}

\begin{savaistu}[Le collagène, cette protéine qui a le temps long]
Le collagène de type I (peau, tendon, os) a une \textbf{demi-vie de renouvellement de 10 à 15 ans} chez l'adulte (mesurée au \ce{^14C} des essais nucléaires atmosphériques). C'est l'une des protéines les plus stables de l'organisme. Conséquence : une glycation (liaison sucre-protéine non enzymatique) ou une cross-link (pontage) anormale persiste des années, rigidifiant le tissu. C'est pourquoi les rides et la raideur s'installent lentement, et pourquoi le contrôle de la glycémie (diabète) et l'apport en vitamine C (cofacteur hydroxylation proline/lysine) sont cruciaux \emph{bien avant} l'apparition des signes visibles.
\end{savaistu}

\begin{exercice}[Entraînement Probatoire - Argumentation complète]
M. Kamga, 65 ans, agriculteur à l'Ouest-Cameroun, consulte pour une fatigue inhabituelle, une perte de force aux champs et une plaie au mollet qui ne cicatrise pas depuis 3 semaines. Il mange surtout du foufou-maïs et du ndolé (peu de viande/poisson), ne fait plus d'exercice hors champs, fume la pipe depuis 40 ans.
À partir de vos connaissances sur le renouvellement moléculaire, rédigez un argumentaire structuré (mécanismes moléculaires $\rightarrow$ conséquences cellulaires $\rightarrow$ signes cliniques $\rightarrow$ conseils) pour expliquer son état et lui proposer un plan d'action réaliste.
\end{exercice}

\begin{corrige}[Exercice : M. Kamga]
\textbf{1. Mécanismes moléculaires :} Tabac (ROS, adduits ADN) + apport protéique insuffisant (substrats synthèse limités) + sédentarité relative (pas de stimulus mTORC1 musculaire) $\rightarrow$ \textbf{ralentissement marqué du turnover} (synthèse $\downarrow\downarrow$, dégradation $\downarrow$ par protéasome/autophagie moins stimulée). Accumulation de protéines carbonylées, lipides peroxydés, lésions ADN non réparées.
\textbf{2. Conséquences cellulaires :} Fibroblastes dermiques sénescents (collagène non renouvelé $\rightarrow$ matrice fragile), myocytes atrophiés (sarcopénie : synthèse myofibrillaire $<$ dégradation), lymphocytes T naïfs rares (thymus involué + turnover lymphocytaire lent) $\rightarrow$ immunosénescence.
\textbf{3. Signes cliniques :} Plaie chronique (collagène non déposé, angiogenèse défaillante, infection sur terrain immunodéprimé), fatigue/marche difficile (sarcopénie, mitochondries dysfonctionnelles $\rightarrow$ ATP $\downarrow$), infections répétées.
\textbf{4. Plan d'action réaliste :}
\begin{itemize}
    \item \textbf{Arrêt tabac} (priorité n°1 : source majeure ROS).
    \item \textbf{Alimentation :} Ajouter source protéique à chaque repas (œufs, poisson sec, haricots/soja, arachides) + légumes-feuilles (folong, éru) + fruits (goavy, papaye) pour vit C/E/Zn/Se.
    \item \textbf{Activité :} Marcher 30 min/jour + exercices de renforcement (lever de chaise, porter eau) pour stimuler mTORC1 musculaire.
    \item \textbf{Soins plaie :} Nettoyage au sérum physiologique / eau bouillie refroidie, pansement humide (vaseline gaze), surveillance infection (signes inflammatoires), consultation infirmière/infirmerie scolaire/CSI.
\end{itemize}
\end{corrige}

\newpage

\coursec{Activité d'intégration}

\courssub{Objectifs de l'activité}

À l'issue de cette activité d'intégration, tu seras capable de :
\begin{itemize}
\item exploiter des données de \cle{marquage isotopique} pour mettre en évidence le renouvellement permanent des molécules constitutives de la cellule ;
\item tracer et interpréter des courbes de décroissance du marquage, puis déterminer graphiquement la \cle{demi-vie} d'une protéine ;
\item comparer les vitesses de renouvellement de deux protéines d'un même organe ;
\item argumenter, à partir de données chiffrées, sur l'importance du \cle{turnover} moléculaire pour la santé de l'organisme, en particulier chez la personne âgée.
\end{itemize}

\courssub{Situation}

Au laboratoire de biochimie de la Faculté de Médecine et des Sciences Biomédicales de l'Université de Yaoundé I, une équipe de chercheurs s'intéresse au \cle{renouvellement moléculaire} des protéines du foie. Cet organe, très actif dans le métabolisme, est fréquemment atteint au Cameroun par les hépatites virales, et les médecins constatent que les patients âgés récupèrent plus lentement que les jeunes après une hépatite.

Pour comprendre ce phénomène, les chercheurs réalisent l'expérience suivante sur des rats de laboratoire :
\begin{itemize}
\item les rats reçoivent, dans leur alimentation, des \cle{acides aminés marqués} à l'azote 15 ($^{15}$N), isotope non radioactif et détectable ;
\item ces acides aminés marqués sont incorporés dans les protéines nouvellement synthétisées par les cellules du foie (hépatocytes) ;
\item après saturation du marquage (jour 0), les rats repassent à une alimentation normale (azote 14 uniquement) ;
\item on mesure alors, à différentes dates, le \cle{pourcentage de marquage résiduel} dans deux protéines du foie : la protéine A (une enzyme du cycle de l'urée) et la protéine B (une protéine structurale des membranes).
\end{itemize}

Les résultats obtenus sont consignés dans le tableau ci-dessous :

\begin{center}
\rowcolors{1}{popblueL}{white}
\begin{tabularx}{\linewidth}{|l|X|X|X|X|}
\hline
\rowcolor{popblue!40} \textbf{Temps (jours)} & \textbf{Jour 0} & \textbf{Jour 2} & \textbf{Jour 4} & \textbf{Jour 6} \\
\hline
Protéine A : marquage résiduel (\%) & 100 & 50 & 25 & 12,5 \\
\hline
Protéine B : marquage résiduel (\%) & 100 & 84 & 71 & 60 \\
\hline
\end{tabularx}
\end{center}

\textbf{Problème posé :} comment ces données prouvent-elles que les protéines de la cellule sont renouvelées en permanence ? Les deux protéines se renouvellent-elles à la même vitesse ? En quoi ces résultats permettent-ils d'expliquer pourquoi une personne âgée récupère moins bien après une hépatite ?

\courssub{Consignes}

Par groupes de quatre élèves, et à l'aide de ton papier millimétré, réponds aux questions suivantes :

\begin{enumerate}
\item Rappelle ce qu'est un \cle{isotope} et explique pourquoi le marquage au $^{15}$N permet de suivre le devenir des protéines dans la cellule.
\item Trace sur un même graphique les deux courbes donnant le pourcentage de marquage résiduel en fonction du temps, pour la protéine A et pour la protéine B (échelle conseillée : 2 cm pour 1 jour en abscisse ; 1 cm pour 10 \% en ordonnée).
\item Décris l'allure de chaque courbe. Que signifie la diminution du marquage résiduel au cours du temps, alors qu'aucune protéine n'a été retirée de la cellule ?
\item Détermine graphiquement la \cle{demi-vie} de chaque protéine, c'est-à-dire le temps au bout duquel le marquage résiduel est tombé à 50 \%. Vérifie la cohérence de ta lecture avec les données du tableau.
\item Compare les vitesses de renouvellement des deux protéines. Propose une explication au fait qu'une enzyme se renouvelle plus vite qu'une protéine structurale.
\item En t'appuyant sur les mécanismes de synthèse et de dégradation moléculaire étudiés dans la leçon, rédige un court paragraphe argumenté (10 à 15 lignes) expliquant pourquoi une personne âgée, dont le renouvellement moléculaire est ralenti, récupère moins bien qu'un jeune adulte après une hépatite.
\end{enumerate}

\courssub{Ressources à mobiliser}

\begin{aretenir}[Ce qu'il faut avoir en tête]
\begin{itemize}
\item \cle{Renouvellement moléculaire} (turnover) : les molécules constitutives de la cellule (protéines, lipides, ADN) sont continuellement synthétisées puis dégradées, même en l'absence de croissance.
\item \cle{Marquage isotopique} : un isotope (ici $^{15}$N) est incorporé aux molécules néoformées ; la disparition progressive du marquage traduit la dégradation des molécules marquées et leur remplacement par des molécules non marquées.
\item \cle{Demi-vie} $t_{1/2}$ : durée nécessaire pour que la moitié des molécules marquées soit remplacée ; plus la demi-vie est courte, plus le renouvellement est rapide.
\item \cle{Mécanismes} : synthèse des protéines (transcription de l'ADN, traduction des ARNm au niveau des ribosomes) ; dégradation (protéases des lysosomes, système ubiquitine-protéasome).
\item \cle{Vieillissement cellulaire} : le ralentissement du turnover entraîne l'accumulation de molécules altérées et une perte d'efficacité des cellules.
\end{itemize}
\end{aretenir}

\courssub{Démarche et corrigé commenté}

\begin{exoresolu}[Corrigé commenté de l'activité]
\textbf{Question 1.} Un isotope est une variété d'un élément chimique dont le noyau possède le même nombre de protons mais un nombre de neutrons différent : $^{15}$N possède un neutron de plus que $^{14}$N. Chimiquement identique à l'azote ordinaire, il est incorporé par la cellule dans les acides aminés puis dans les protéines exactement comme $^{14}$N, mais il est détectable par spectrométrie de masse. On peut donc « marquer » les protéines synthétisées pendant la période d'administration, puis suivre leur disparition : c'est le principe de la \cle{trace isotopique}.
\tcblower
\textbf{Question 2.} Les deux courbes sont tracées sur le même repère (temps en abscisse, marquage résiduel en ordonnée).

\begin{popfigure}
\begin{center}
\begin{tikzpicture}
\begin{axis}[
width=0.9\linewidth, height=7cm,
xlabel={Temps (jours)}, ylabel={Marquage résiduel (\%)},
xmin=0, xmax=8, ymin=0, ymax=105,
xtick={0,1,2,3,4,5,6,7,8}, ytick={0,10,20,30,40,50,60,70,80,90,100},
grid=both, grid style={popline!50},
legend style={at={(0.97,0.97)}, anchor=north east, font=\small}
]
\addplot[popcurve, color=popblue, mark=*, mark options={fill=popblue}] coordinates {(0,100) (2,50) (4,25) (6,12.5)};
\addlegendentry{Protéine A}
\addplot[popcurve, color=poporange, mark=square*, mark options={fill=poporange}] coordinates {(0,100) (2,84) (4,71) (6,60)};
\addlegendentry{Protéine B}
\addplot[dashed, color=popmuted] coordinates {(0,50) (2,50)};
\addplot[dashed, color=popmuted] coordinates {(2,0) (2,50)};
\end{axis}
\end{tikzpicture}
\end{center}
\legende{Courbes de décroissance du marquage au $^{15}$N dans deux protéines du foie de rat. La protéine A perd la moitié de son marquage en 2 jours (pointillés), la protéine B beaucoup plus lentement.}
\end{popfigure}

\textbf{Question 3.} Les deux courbes sont \cle{décroissantes} : le marquage résiduel diminue au cours du temps. Celle de la protéine A décroît rapidement (allure exponentielle), celle de la protéine B décroît lentement et presque linéairement sur la période étudiée. Or aucune protéine n'a été extraite de la cellule : la disparition du marquage signifie donc que les protéines marquées sont \textbf{dégradées} et \textbf{remplacées} par des protéines nouvelles, synthétisées à partir d'acides aminés non marqués ($^{14}$N). C'est la preuve expérimentale du \cle{renouvellement permanent} des molécules cellulaires.

\textbf{Question 4.} On lit graphiquement le temps correspondant à 50 \% de marquage résiduel :
\begin{itemize}
\item protéine A : la courbe coupe la droite des 50 \% à $t = 2$ jours, donc $t_{1/2}(\mathrm{A}) \approx 2$ jours. Vérification avec le tableau : au jour 2, le marquage vaut exactement 50 \% ; au jour 4, il vaut 25 \% (soit la moitié de 50 \%) ; au jour 6, 12,5 \% (moitié de 25 \%). Chaque période de 2 jours divise le marquage par deux : la demi-vie est bien de 2 jours ;
\item protéine B : la courbe atteint 50 \% au-delà du jour 6 ; en prolongeant la décroissance (perte d'environ 8 \% par jour, soit 40 \% en 6 jours), on estime $t_{1/2}(\mathrm{B}) \approx 8$ jours.
\end{itemize}

\textbf{Question 5.} La protéine A se renouvelle environ quatre fois plus vite que la protéine B ($8 / 2 = 4$). La protéine A est une \textbf{enzyme} : son activité catalytique est vitale et toute molécule altérée doit être rapidement remplacée pour maintenir l'efficacité du métabolisme ; de plus, le renouvellement rapide des enzymes permet à la cellule d'ajuster vite leurs quantités selon les besoins. La protéine B est \textbf{structurale} : elle assure un rôle de charpente stable dans les membranes, qui n'exige pas un remplacement fréquent. La vitesse de renouvellement est donc adaptée à la fonction de chaque protéine.

\textbf{Question 6.} Paragraphe argumenté (modèle de réponse) : « Le foie lésé par une hépatite virale doit remplacer massivement ses protéines endommagées et régénérer ses hépatocytes. Or la guérison repose sur le renouvellement moléculaire : synthèse de nouvelles protéines par traduction des ARNm et élimination des protéines altérées par les protéases lysosomales et le protéasome. Chez le sujet jeune, ce turnover est rapide : comme la protéine A de notre expérience, les protéines essentielles sont remplacées en quelques jours, et l'intégrité cellulaire est vite restaurée. Chez la personne âgée, le renouvellement moléculaire est ralenti : la synthèse protéique diminue et les systèmes de dégradation deviennent moins efficaces, de sorte que les molécules altérées s'accumulent dans les cellules, comme le montre la courbe lente de la protéine B. Les hépatocytes âgés retrouvent donc lentement leur fonctionnalité, ce qui explique la convalescence prolongée et les complications plus fréquentes observées après une hépatite chez les patients âgés. Le maintien d'un turnover efficace est ainsi une condition essentielle de la santé de l'organisme. »
\end{exoresolu}

\begin{attention}[Pièges à éviter dans cette activité]
\begin{itemize}
\item Ne pas confondre \textbf{disparition du marquage} et \textbf{destruction de la cellule} : les hépatocytes restent en place, seules leurs molécules sont remplacées.
\item La demi-vie se lit à \textbf{50 \%} du marquage initial, pas à 0 \%.
\item Attention aux échelles du graphique : une lecture graphique fausse de $t_{1/2}$ vient le plus souvent d'une échelle mal choisie ou mal reportée.
\item Ne pas écrire que « les protéines vieillissent » sans préciser le mécanisme : il faut parler de \textbf{ralentissement des synthèses et des dégradations}, avec accumulation de molécules altérées.
\end{itemize}
\end{attention}

\courssub{Bilan de l'activité}

Cette activité t'a permis de mettre en évidence, à partir de données expérimentales réelles de type marquage isotopique, que :
\begin{itemize}
\item les protéines de la cellule sont \textbf{renouvelées en permanence}, même dans un organe qui ne croît plus : la décroissance du marquage au $^{15}$N en est la preuve directe ;
\item la vitesse de renouvellement, mesurée par la \cle{demi-vie}, \textbf{varie selon les protéines} : environ 2 jours pour l'enzyme A, environ 8 jours pour la protéine structurale B, ce qui montre que le turnover est adapté à la fonction de chaque molécule ;
\item le renouvellement moléculaire, qui résulte de l'équilibre entre \textbf{synthèses} (transcription, traduction) et \textbf{dégradations} (lysosomes, protéasome), est indispensable au \textbf{maintien de l'intégrité cellulaire} ;
\item le \textbf{ralentissement du turnover avec l'âge} explique l'accumulation de molécules altérées, le vieillissement cellulaire et la moindre capacité de récupération des personnes âgées après une agression comme une hépatite.
\end{itemize}

Tu as ainsi mobilisé les trois savoir-faire visés par le programme : exploiter des données de marquage isotopique, relier le renouvellement moléculaire au maintien de l'intégrité cellulaire, et argumenter sur son importance pour la santé de l'organisme. Garde cette démarche en mémoire : au Probatoire, un exercice d'analyse de courbes de marquage se résout exactement de la même manière — lecture rigoureuse des données, détermination de la demi-vie, puis interprétation biologique argumentée.

\newpage

\coursec{Méthodes et exercices résolus}

\courssub{Méthodes et exercices résolus}

\begin{methode}[Exploiter des données de marquage isotopique (cinétique d'incorporation)]
\textbf{Objectif :} Déterminer le taux de renouvellement d'une macromolécule (protéine, ADN, lipide) à partir d'une courbe d'incorporation d'un isotope (ex. \ce{^{14}C}, \ce{^{3}H}, \ce{^{15}N}, \ce{^{13}C}).
\par
\textbf{Démarche en 5 étapes :}
\begin{enumerate}
    \item \textbf{Identifier le protocole :} Pulse (marquage court) ou Pulse-Chase (marquage puis chasse au substrat non marqué). Repérer l'axe des abscisses (temps) et l'axe des ordonnées (radioactivité ou abondance isotopique, souvent en \% du max ou en dpm/µg ADN).
    \item \textbf{Repérer la phase d'incorporation (Pulse) :} La courbe monte. La pente initiale reflète la vitesse de synthèse. Le plateau correspond au marquage maximal (pool précurseur saturé).
    \item \textbf{Repérer la phase de décroissance (Chase) :} La courbe descend. Elle traduit la dégradation des molécules marquées et leur remplacement par des molécules non marquées (turnover).
    \item \textbf{Calculer la demi-vie de renouvellement ($t_{1/2}$) :} Sur la phase de décroissance (ou sur la montée si pas de chase), trouver le temps nécessaire pour que le signal perde la moitié de sa valeur maximale (ou atteigne 50\% du plateau). \textbf{Formule :} $t_{1/2} = \frac{\ln 2}{k_d}$ où $k_d$ est la constante de dégradation.
    \item \textbf{Conclure sur le turnover :} Un $t_{1/2}$ court = renouvellement rapide (protéines régulatrices, ARN) ; $t_{1/2}$ long = renouvellement lent (protéines structurales, ADN nucléaire). Comparer aux valeurs de référence du cours.
\end{enumerate}
\textbf{Réflexe Probatoire :} Préciser l'unité de temps (heures, jours). Ne pas confondre « demi-vie radioactive » (physique) et « demi-vie biologique de turnover » (biologique).
\end{methode}

\begin{exoresolu}[Détermination du turnover de l'albumine hépatique par marquage au \ce{^{14}C}-leucine]
\textbf{Énoncé :} On réalise une expérience de \emph{pulse-chase} sur des hépatocytes en culture. On incorpore de la \ce{^{14}C}-leucine pendant 30 min (Pulse), puis on remplace le milieu par un milieu contenant un excès de leucine non marquée (Chase). On mesure la radioactivité spécifique de l'albumine immunoprécipitée à différents temps. Résultats :
\begin{center}
\begin{tabular}{|c|c|c|c|c|c|c|}\hline
Temps de chase (h) & 0 & 2 & 4 & 6 & 8 & 10 \\ \hline
Radioactivité spécifique (dpm/µg albumine) & 1200 & 950 & 750 & 590 & 470 & 370 \\ \hline
\end{tabular}
\end{center}
\begin{enumerate}
    \item Tracer la courbe de décroissance de la radioactivité spécifique en fonction du temps de chase.
    \item Déterminer graphiquement la demi-vie ($t_{1/2}$) de l'albumine dans ces conditions.
    \item Calculer la constante de dégradation $k_d$ (en h$^{-1}$).
    \item Interpréter biologiquement ce résultat : l'albumine est-elle une protéine à turnover rapide ou lent ? Justifier.
\end{enumerate}
\tcblower
\textbf{Solution détaillée :}
\begin{enumerate}
    \item \textbf{Tracé de la courbe :} (Voir figure ci-dessous). On utilise un repère orthogonal. Abscisses : Temps de chase (h) de 0 à 10. Ordonnées : Radioactivité spécifique (dpm/µg) de 0 à 1300. Les points s'alignent sur une courbe exponentielle décroissante.
    \begin{popfigure}
    \begin{tikzpicture}[scale=0.7]
    \begin{axis}[
        width=\linewidth, height=8cm,
        xlabel={Temps de chase (h)}, ylabel={Radioactivité sp. (dpm/$\mu$g)},
        xmin=0, xmax=11, ymin=0, ymax=1300,
        grid=major, grid style={popline},
        tick label style={font=\small},
        label style={font=\small},
        legend style={at={(0.98,0.98)}, anchor=north east, font=\small, fill=white, draw=popdark}
    ]
    \addplot[popcurve, mark=*, mark size=2pt, color=popblue] coordinates {
        (0,1200) (2,950) (4,750) (6,590) (8,470) (10,370)
    };
    \addlegendentry{Données exp.}
    % Fit exponentiel approximatif pour guide visuel : y = 1200 * exp(-0.11 * x)
    \addplot[poporange, dashed, domain=0:10, samples=100] {1200*exp(-0.11*x)};
    \addlegendentry{Ajustement $y=1200e^{-0.11t}$}
    % Ligne demi-vie
    \draw[popgreen, dashed] (0,600) -- (6.3,600) node[midway, above, font=\tiny] {$R_0/2$};
    \draw[popgreen, dashed] (6.3,0) -- (6.3,600) node[midway, right, font=\tiny] {$t_{1/2} \approx 6{,}3$ h};
    \end{axis}
    \end{tikzpicture}
    \legende{Courbe de décroissance de la radioactivité spécifique de l'albumine lors du chase. La demi-vie graphique est d'environ 6,3 h.}
    \end{popfigure}
    \item \textbf{Détermination graphique de $t_{1/2}$ :}
    Radioactivité initiale $R_0 = 1200$ dpm/µg. Demi-valeur = $600$ dpm/µg.
    Sur la courbe (ou par interpolation entre 6h (590) et 4h (750)), le temps correspondant à 600 dpm/µg est \textbf{$t_{1/2} \approx 6{,}3$ h}.
    \item \textbf{Calcul de $k_d$ :}
    \[ k_d = \frac{\ln 2}{t_{1/2}} = \frac{0{,}693}{6{,}3\,\text{h}} \approx \textbf{0{,}11 h}^{-1} \]
    \item \textbf{Interprétation biologique :}
    Une demi-vie de 6,3 heures indique un \textbf{turnover rapide}. L'albumine est une protéine sécrétée massivement par le foie (jusqu'à 15 g/jour chez l'adulte). Son renouvellement rapide permet d'ajuster la concentration plasmatique (pression oncotique, transport) aux besoins physiologiques variables. Cela contraste avec le collagène (demi-vie de mois/années) ou l'ADN (quasi nul en phase G0).
\end{enumerate}
\textbf{Conclusion :} L'albumine hépatique a une demi-vie de turnover d'environ 6,3 h ($k_d \approx 0{,}11$ h$^{-1}$), caractéristique d'une protéine à renouvellement rapide, adaptée à sa fonction de protéine plasmatique majeure dont la synthèse est très régulée.
\end{exoresolu}

\begin{methode}[Calculer le flux de renouvellement (turnover) à l'état stationnaire]
\textbf{Objectif :} Relier la masse molaire d'un pool moléculaire, sa demi-vie et le flux de synthèse/dégradation quotidien.
\par
\textbf{Démarche en 4 étapes :}
\begin{enumerate}
    \item \textbf{Vérifier l'état stationnaire :} La masse du pool $M$ est constante $\Rightarrow$ Vitesse de synthèse $v_s$ = Vitesse de dégradation $v_d$ = Flux $J$.
    \item \textbf{Utiliser la relation fondamentale :} $J = k_d \times M = \frac{\ln 2}{t_{1/2}} \times M$.
    \item \textbf{Convertir les unités :} $M$ en mol (ou g), $t_{1/2}$ en jours (ou heures), $J$ en mol/jour (ou g/jour).
    \item \textbf{Interpréter le flux :} Un flux élevé avec un pool petit = turnover très rapide (ex. ATP, protéines régulatrices). Un flux faible avec un pool grand = turnover lent (ex. collagène, cristallines du cristallin).
\end{enumerate}
\textbf{Formules à retenir :}
\[ J\,(\text{mol/jour}) = \frac{0{,}693}{t_{1/2}\,(\text{jour})} \times M\,(\text{mol}) \]
\[ \text{Taux de renouvellement journalier (\%/jour)} = \frac{0{,}693}{t_{1/2}\,(\text{jour})} \times 100 \]
\textbf{Attention :} $t_{1/2}$ doit être en \textbf{jours} pour un flux journalier.
\end{methode}

\begin{exoresolu}[Flux de renouvellement de l'ATP musculaire au repos et à l'effort]
\textbf{Énoncé :} Chez un sportif, le pool d'ATP musculaire est d'environ 5 mmol/kg de muscle. Au repos, la demi-vie de l'ATP est de 1 minute. Lors d'un effort intense, elle chute à 1 seconde.
\begin{enumerate}
    \item Calculer le flux de turnover de l'ATP (en mmol/kg/min) au repos et à l'effort.
    \item Exprimer ce flux en \% du pool renouvelé par minute dans les deux cas.
    \item Expliquer comment la cellule maintient la concentration [ATP] quasi constante malgré ce flux énorme à l'effort.
\end{enumerate}
\tcblower
\textbf{Solution détaillée :}
\begin{enumerate}
    \item \textbf{Flux au repos :} $t_{1/2} = 1$ min. $M = 5$ mmol/kg.
    $k_d = \frac{\ln 2}{1\,\text{min}} = 0{,}693\,\text{min}^{-1}$.
    $J_{\text{repos}} = k_d \times M = 0{,}693 \times 5 = \textbf{3{,}47 mmol/kg/min}$.
    \textbf{Flux à l'effort :} $t_{1/2} = 1$ s = $1/60$ min.
    $k_d = \frac{\ln 2}{1/60} = 0{,}693 \times 60 = 41{,}58\,\text{min}^{-1}$.
    $J_{\text{effort}} = 41{,}58 \times 5 = \textbf{207{,}9 mmol/kg/min}$.
    \item \textbf{\% renouvelé/min :}
    Repos : $\frac{J}{M} \times 100 = k_d \times 100 = \textbf{69{,}3\%/min}$ (le pool entier tourne en ~1,4 min).
    Effort : $41{,}58 \times 100 = \textbf{4158\%/min}$ (le pool entier tourne 41 fois par minute !).
    \item \textbf{Maintien de l'homéostasie [ATP] :} La concentration [ATP] reste stable (~5 mM) car \textbf{synthèse = dégradation} à chaque instant (état stationnaire dynamique). À l'effort, l'activation de la glycolyse, de la phosphocréatine kinase et de la respiration mitochondriale augmente $v_s$ pour égaler exactement $v_d$ augmentée. Les systèmes de rétroaction (ADP, AMP, Ca$^{2+}$ activant les kinases) assurent ce couplage fin. Si $v_s < v_d$ ne serait-ce que quelques secondes, [ATP] chuterait, entraînant la rigor mortis (contracture irréversible).
\end{enumerate}
\textbf{Conclusion :} Le flux de turnover de l'ATP passe de ~3,5 à ~208 mmol/kg/min (facteur 60), illustrant la réactivité extrême du métabolisme énergétique. L'intégrité cellulaire repose sur l'égalité stricte et instantanée Synthèse = Dégradation.
\end{exoresolu}

\begin{methode}[Relier renouvellement moléculaire et intégrité cellulaire (argumentation structurée)]
\textbf{Objectif :} Construire une réponse argumentée (type dissertation courte ou QROC) reliant turnover, contrôle qualité, adaptation et pathologie.
\par
\textbf{Plan type en 4 arguments (méthode « 4 piliers ») :}
\begin{enumerate}
    \item \textbf{Pilier 1 : Élimination des macromolécules endommagées (Contrôle Qualité).} 
    Renouvellement = dégradation des protéines mal repliées (système ubiquitine-protéasome, autophagy), des lipides peroxydés, de l'ADN muté (réparation/excision). \emph{Mots-clés :} Protéostase, stress oxydant, agrégats toxiques (maladies neurodégénératives).
    \item \textbf{Pilier 2 : Adaptation phénotypique (Plasticité).} 
    Changement du protéome/lipidome sans changer le génome. Ex. : induction CYP450 (xénobiotiques), récepteurs hormonaux, enzymes glycolytiques (effort), protéines de choc thermique. Permet la réponse aux signaux environnementaux.
    \item \textbf{Pilier 3 : Régulation métabolique et signalisation (Contrôle des flux).} 
    Turnover rapide des messagers (cAMP, Ca$^{2+}$, IP3, NO, protéines phosphorylées, ARNm courts) = réversibilité du signal. Turnover des enzymes clés (phosphofructokinase, HMG-CoA réductase) = ajustement des flux métaboliques.
    \item \textbf{Pilier 4 : Économie énergétique et allocation des ressources.} 
    Coût ATP élevé (4 ATP/liaison peptide, coût transcription/réplication). Le turnover est sélectif : protéines structurelles stables, protéines régulatrices instables (séquences PEST, boîtes de destruction). Déséquilibre = pathologie (cachexie, atrophie, fibrose).
\end{enumerate}
\textbf{Rédaction Probatoire :} Utiliser les connecteurs logiques : \emph{« En effet », « Par conséquent », « Cependant », « De plus »}. Citer des exemples moléculaires précis (ubiquitine, protéasome, chaperonnes, ARN interférents).
\end{methode}

\begin{exoresolu}[Argumenter : Le ralentissement du turnover protéique dans le vieillissement musculaire (sarcopénie)]
\textbf{Énoncé :} Chez le sujet âgé, on observe une diminution de la masse musculaire (sarcopénie) associée à une « résistance anabolique ». À partir de vos connaissances sur le renouvellement moléculaire, expliquez comment un ralentissement du turnover des protéines myofibrillaires contribue à la perte de fonction musculaire. Structurez votre réponse en 3 arguments distincts.
\tcblower
\textbf{Solution détaillée (Modèle de rédaction attendue) :}

\textbf{Argument 1 : Accumulation de protéines endommagées et dysfonctionnement contractile.}
Le ralentissement du turnover (allongement de la demi-vie des protéines myofibrillaires comme l'actine, la myosine, la titine) réduit l'efficacité du système de contrôle qualité (ubiquitine-protéasome et autophagy/mitophagie). Les protéines oxydées (carbo-nylation par ROS mitochondriaux), glyquées (produits de Maillard) ou mal repliées ne sont plus éliminées assez vite. Elles forment des agrégats insolubles qui perturbent l'architecture du sarcomère, diminuent la force spécifique (force/section) et augmentent la raideur passive du muscle. \emph{Conséquence :} Perte de qualité contractile avant même la perte de quantité.

\textbf{Argument 2 : Incapacité d'adaptation plastique aux stimuli mécaniques et nutritionnels.}
Le muscle adulte conserve une plasticité phénotypique : hypertrophie (voie mTOR/IGF-1) ou atrophie (voie FoxO/Atrogin-1-MuRF1). Cette plasticité nécessite un turnover basal élevé pour permettre le remplacement rapide des isoformes (ex. passage MyHC-IIx vers IIa à l'entraînement) et l'élimination des composants obsolètes. Un turnover ralenti « fige » le protéome : la synthèse de nouvelles protéines fonctionnelles est freinée (résistance anabolique : signal mTOR affaibli, inflammation chronique TNF-$\alpha$/IL-6), tandis que la dégradation des protéines endommagées est insuffisante. Le muscle ne peut plus s'adapter à la charge ni se réparer après une microlésion.

\textbf{Argument 3 : Déséquilibre du bilan azoté et coût énergétique délétère.}
Le maintien d'un pool protéique stable avec un turnover ralenti masque un déséquilibre : la synthèse chute plus vite que la dégradation (ou inversément en cas d'inflammation aiguë). Le coût énergétique du turnover (environ 20\% de la dépense énergétique de base) diminue, mais au prix d'une accumulation de « déchets » moléculaires. De plus, la résistance anabolique empêche l'utilisation efficace des apports protéiques (leucine) pour stimuler la synthèse. Le bilan azoté devient négatif $\rightarrow$ fonte musculaire progressive.

\textbf{Synthèse :} Le ralentissement du turnover n'est pas une simple conséquence passive du vieillissement, mais un \textbf{moteur actif} de la sarcopénie : il couple perte de qualité (protéines endommagées), perte de plasticité (incapacité d'adaptation) et bilan catabolique net.
\end{exoresolu}

\begin{methode}[Interpréter une courbe de Pulse-Chase classique (ADN, Protéine, Lipide)]
\textbf{Objectif :} Différencier visuellement et quantitativement les profils de renouvellement : ADN (semi-conservatif, pas de turnover en G0), Protéines (turnover variable), Lipides (turnover rapide membranes).
\par
\textbf{Grille de lecture en 3 temps :}
\begin{enumerate}
    \item \textbf{Phase Pulse (Montée) :} 
    \begin{itemize}
        \item \textbf{Pente raide + Plateau haut + Rapide :} Précurseur incorporé vite, pool tournant vite (ARNm, protéines régulatrices, phospholipides membranaires).
        \item \textbf{Pente douce + Plateau bas + Lent :} Pool stable, dilution par division cellulaire seulement (ADN en phase S, histones, cristallines).
    \end{itemize}
    \item \textbf{Phase Chase (Descente) :} 
    \begin{itemize}
        \item \textbf{Descente exponentielle nette :} Turnover actif (dégradation + resynthèse non marquée). $t_{1/2}$ mesurable.
        \item \textbf{Plateau maintenu (pas de descente) :} Pas de turnover (ADN en G0/G1, protéines structurelles très stables). La radioactivité reste dans la molécule mère.
        \item \textbf{Descente par paliers (division cellulaire) :} Dilution par mitose (ADN semi-conservatif : 100\% $\to$ 50\% $\to$ 25\% par division).
    \end{itemize}
    \item \textbf{Comparer les demi-vies :} Classer les molécules : $t_{1/2}$(ARNm) $\ll$ $t_{1/2}$(Protéines régulatrices) $<$ $t_{1/2}$(Enzymes métaboliques) $<$ $t_{1/2}$(Protéines structurelles) $\ll$ $t_{1/2}$(ADN nucléaire).
\end{enumerate}
\textbf{Piège fréquent :} Confondre « pas de décroissance en chase » (molécule stable) et « décroissance par dilution mitotique » (ADN). L'ADN ne se dégrade pas (sauf apoptose), il se dilue.
\end{methode}

\begin{exoresolu}[Analyse comparative : ADN vs Protéine mitochondriale vs Phospholipide]
\textbf{Énoncé :} Trois expériences de Pulse-Chase au \ce{^{3}H}-thymidine (ADN), \ce{^{35}S}-méthionine (protéines) et \ce{^{14}C}-acetate (lipides) sont menées sur des fibroblastes en culture en phase de croissance exponentielle. Les courbes de radioactivité spécifique (normalisée à 100\% à la fin du pulse) sont obtenues.
\begin{enumerate}
    \item Pour chaque type de molécule, décrire la forme attendue de la courbe (montée + descente) et justifier par le mécanisme moléculaire.
    \item Préciser pour l'ADN si la décroissance observée reflète une dégradation ou une dilution. Calculer le \% de radioactivité restant après 2 divisions cellulaires.
    \item Une protéine mitochondriale (ex. cytochrome c) a un $t_{1/2}$ de 5 jours. Un phospholipide membranaire a un $t_{1/2}$ de 2 jours. Lequel a le turnover le plus rapide ? Calculer le rapport de leurs constantes de dégradation ($k_{d,\text{lipide}} / k_{d,\text{protéine}}$).
\end{enumerate}
\tcblower
\textbf{Solution détaillée :}
\begin{enumerate}
    \item \textbf{ADN (\ce{^{3}H}-thymidine) :} Montée uniquement en phase S (pic). \textbf{Chase : Plateau horizontal (pas de dégradation)}. Si les cellules divisent : décroissance par paliers de 50\% par mitose (dilution semi-conservative). \emph{Mécanisme :} Réplication semi-conservative, pas de turnover enzymatique de l'ADN nucléaire en conditions normales (réparation excisionnelle invisible à cette échelle).
    \textbf{Protéine mitochondriale (\ce{^{35}S}-Méthionine) :} Montée progressive (temps de traduction/import). \textbf{Chase : Décroissance exponentielle continue.} $t_{1/2} \approx 5$ j. \emph{Mécanisme :} Dégradation par protéases mitochondriales (Lon, ClpP) et mitophagie, resynthèse continue.
    \textbf{Phospholipide (\ce{^{14}C}-Acetate) :} Montée très rapide (voie métabolique courte). \textbf{Chase : Décroissance exponentielle rapide.} $t_{1/2} \approx 2$ j. \emph{Mécanisme :} Remodelage membranaire constant (phospholipases, acyltransférases), turnover des membranes (endocytose/exocytose, autophagie).
    \item \textbf{ADN : Dilution, pas dégradation.} Après 1 division : 50\% radioactivité. Après 2 divisions : $50\% \times 50\% = \textbf{25\%}$ de la radioactivité initiale par noyau (ou 25\% de l'ADN total est marqué sur les deux brins).
    \item \textbf{Comparaison turnover :}
    $k_{d,\text{protéine}} = \ln 2 / 5 = 0{,}1386\,\text{jour}^{-1}$.
    $k_{d,\text{lipide}} = \ln 2 / 2 = 0{,}3466\,\text{jour}^{-1}$.
    Le \textbf{phospholipide} a le turnover le plus rapide (demi-vie plus courte).
    Rapport $k_{d,\text{lipide}} / k_{d,\text{protéine}} = 0{,}3466 / 0{,}1386 = \textbf{2{,}5}$.
    Le phospholipide se renouvelle 2,5 fois plus vite que la protéine mitochondriale.
\end{enumerate}
\textbf{Conclusion :} La fluidité membranaire exige un turnover lipidique très rapide (2 j), tandis que les complexes respiratoires mitochondriaux sont plus stables (5 j), l'ADN nucléaire étant quasi stable hors division.
\end{exoresolu}

\begin{methode}[Argumenter l'importance du renouvellement pour la santé : Structure « Pathologie $\leftrightarrow$ Mécanisme $\leftrightarrow$ Thérapeutique »]
\textbf{Objectif :} Répondre à une question de synthèse type « Discutez l'importance du renouvellement moléculaire pour la santé » ou « Expliquez la base moléculaire d'une maladie liée au turnover ».
\par
\textbf{Structure en 3 colonnes (à faire mentalement ou en brouillon) :}
\begin{center}
\begin{tabular}{|l|l|l|}\hline
\textbf{Pathologie / Situation} & \textbf{Mécanisme : Défaut de renouvellement} & \textbf{Piste thérapeutique / Prévention} \\ \hline
Neurodégénérescence (Alzheimer, Parkinson, Huntington) & Accumulation d'agrégats (A$\beta$, $\alpha$-synucleine, Huntingtine) $\leftarrow$ Défaillance protéasome / autophagy / chaperonnes (turnover ralenti/saturé) & Inducteurs d'autophagie (rapamycine), chaperonnes pharmacologiques, anticorps monoclonaux \\ \hline
Cancer & Turnover déséquilibré : Protéines onco-géniques stabilisées (p53 muté non dégradé, Myc, cyclines) OU Suppresseurs de tumeurs dégradés trop vite (p53 WT par MDM2) & Inhibiteurs de protéasome (Bortézomib - myélome), Inhibiteurs MDM2, PROTACs (dégradation ciblée) \\ \hline
Maladies métaboliques (Diabète T2, Stéatose/NASH) & Turnover lipidique ralenti (lipogenèse > lipophagie/$\beta$-oxydation) $\rightarrow$ Steatose. Résistance à l'insuline $\rightarrow$ signalisation altérée (turnover récepteurs IRS) & Agonistes GLP-1 (améliorent autophagy hépatique), Jeûne intermittent (induit autophagy), Exercice (turnover mitochondrial) \\ \hline
Vieillissement / Sarcopénie / Immunosénescence & Turnover global ralenti (protéostase effondrée, mitophagie défaillante, thymus involué $\rightarrow$ renouvellement répertoire TCR ralenti) & Restriction calorique, Exercice, Sénolytiques, Thérapies géniques (telomérase) \\ \hline
Intoxications / Médicaments & Saturation voies de dégradation (ex. paracétamol $\rightarrow$ NAPQI déplete GSH, adducts protéiques non éliminés) & Antidots (N-acétylcystéine), Dialyse, Induction CYP \\ \hline
\end{tabular}
\end{center}
\textbf{Conseil rédaction :} Choisir 2 exemples précis. Pour chacun : \textbf{1. Molécule concernée} $\rightarrow$ \textbf{2. Type de défaut (synthèse ? dégradation ? reconnaissance ?)} $\rightarrow$ \textbf{3. Conséquence cellulaire} $\rightarrow$ \textbf{4. Symptôme clinique} $\rightarrow$ \textbf{5. Levier d'action}.
\end{methode}

\begin{exoresolu}[Base moléculaire de l'efficacité du bortézomib (Velcade\textsuperscript{\textregistered}) dans le myélome multiple]
\textbf{Énoncé :} Le myélome multiple est une hémopathie maligne des plasmocytes sécréteurs d'immunoglobulines monoclonales. Le bortézomib, inhibiteur réversible du protéasome 26S (sous-unité $\beta$5, activité chymotrypsine-like), est un traitement de référence. Expliquez, en mobilisant la notion de renouvellement moléculaire, pourquoi les cellules de myélome sont hypersensibles à l'inhibition du protéasome par rapport aux cellules saines.
\tcblower
\textbf{Solution détaillée :}

\textbf{1. Contexte physiologique : Charge protéique extrême.}
Les plasmocytes cancéreux (myélome) synthétisent et sécrètent des quantités massives d'immunoglobulines (Ig) monoclonales (paraprotéine), pouvant atteindre des dizaines de grammes par jour. Cela impose un \textbf{flux de synthèse protéique énorme} au niveau du réticulum endoplasmique (RE).

\textbf{2. Mécanisme : Dépendance au contrôle qualité RE-Associated Degradation (ERAD).}
\begin{itemize}
    \item Une fraction significative des chaînes lourdes/légères d'Ig ne se replie pas correctement ou n'est pas assemblée.
    \item Ces protéines mal repliées sont rétrotranslocées du RE vers le cytosol et \textbf{ubiquitinées} pour être dégradées par le \textbf{protéasome 26S} (voie ERAD).
    \item Dans le myélome, le \textbf{flux de substrats ERAD est saturé/maximal}. Le protéasome fonctionne à flux tendu pour éviter l'accumulation toxique d'agrégats dans le RE (stress RE).
\end{itemize}

\textbf{3. Effet de l'inhibition (Bortézomib) : Rupture de l'homéostasie (Protéostase).}
\begin{itemize}
    \item Inhibition du protéasome $\rightarrow$ \textbf{Blocage brutal de la dégradation} des protéines ubiquitinées (ERAD + turnover normal protéines courtes de vie : cyclines, p53, I$\kappa$B, NOXA).
    \item \textbf{Accumulation massive} de protéines mal repliées dans le RE $\rightarrow$ \textbf{Stress RE sévère non résolu} $\rightarrow$ Activation de la réponse aux protéines non repliées (UPR : PERK, IRE1, ATF6).
    \item Basculement de l'UPR adaptatif vers \textbf{apoptose} (via CHOP, caspase-12, JNK, montée de NOXA/PUMA non dégradées).
    \item Parallèlement : Stabilisation de p53 (arrêt cycle), inhibition NF-$\kappa$B (I$\kappa$B$\alpha$ non dégradé $\rightarrow$ survie bloquée), accumulation cyclines (catastrophe mitotique).
\end{itemize}

\textbf{4. Fenêtre thérapeutique (Sélectivité) :}
Les cellules saines (turnover protéique basal faible, pas de charge ERAD massive) tolèrent une inhibition partielle et transitoire du protéasome (récupération après clairance du drogue). Les cellules de myélome, \textbf{« accros » au protéasome} (concept \emph{proteasome addiction}), basculent rapidement en apoptose.

\textbf{Conclusion :} L'efficacité du bortézomib illustre dramatiquement que \textbf{l'intégrité cellulaire dépend d'un équilibre dynamique Synthèse = Dégradation}. Briser le bras « Dégradation » dans une cellule à flux de synthèse maximal provoque un effondrement catastrophique de la protéostase, tuant sélectivement la cellule cancéreuse.
\end{exoresolu}

\begin{methode}[Modéliser simplement la dynamique d'un pool moléculaire (Équation différentielle du turnover)]
\textbf{Objectif :} Formaliser mathématiquement le renouvellement pour prévoir l'évolution d'un pool (ex. prise de médicament, jeûne, pathologie).
\par
\textbf{Modèle de base (État stationnaire perturbé) :}
Soit $M(t)$ la masse (ou quantité) de la molécule à l'instant $t$.
\[ \frac{dM}{dt} = v_s(t) - k_d \cdot M(t) \]
\begin{itemize}
    \item $v_s(t)$ : Vitesse de synthèse (entrée), peut varier dans le temps (induction, répression).
    \item $k_d \cdot M(t)$ : Vitesse de dégradation (sortie), proportionnelle à la masse (loi du 1er ordre). $k_d = \ln 2 / t_{1/2}$.
\end{itemize}
\textbf{Cas 1 : État stationnaire (Normal).} $v_s = cte$, $dM/dt = 0 \Rightarrow M_{eq} = v_s / k_d$.
\textbf{Cas 2 : Arrêt brutal de la synthèse ($v_s = 0$, ex. Actinomycine D, Cycloheximide, Jeûne total).}
\[ \frac{dM}{dt} = -k_d M \Rightarrow M(t) = M_0 \cdot e^{-k_d t} = M_0 \cdot 2^{-t/t_{1/2}} \]
\textbf{Cas 3 : Induction de la synthèse ($v_s$ passe de $v_0$ à $v_1 > v_0$).}
\[ M(t) = M_{eq,1} + (M_0 - M_{eq,1}) e^{-k_d t} \quad \text{avec } M_{eq,1} = v_1/k_d \]
Le temps pour atteindre 90\% du nouveau plateau $\approx 3{,}3 \times t_{1/2}$.
\textbf{Application Probatoire :} On ne demande pas de résoudre l'équation différentielle, mais de \textbf{l'écrire}, d'\textbf{identifier les termes} et d'\textbf{interpréter les paramètres} ($k_d$, $v_s$, $M_{eq}$).
\end{methode}

\begin{exoresolu}[Modélisation de la disparition de l'hémoglobine glyquée (HbA1c) après normalisation de la glycémie]
\textbf{Énoncé :} L'HbA1c reflète la glycémie moyenne sur la durée de vie de l'érythrocyte (~120 jours). On considère que la formation de l'HbA1c est proportionnelle à la glycémie (synthèse « virtuelle » sur l'hémoglobine existante) et que sa disparition suit le turnover des globules rouges (pas de déglycation enzymatique significative). Un patient diabétique mal équilibré a une HbA1c à 9\%. Il met en place un traitement efficace qui normalise sa glycémie (cible HbA1c d'équilibre = 5\%) instantanément à $t=0$.
\begin{enumerate}
    \item Écrire l'équation différentielle régissant l'évolution du taux d'HbA1c ($H$) en fonction du temps, en définissant les paramètres.
    \item Exprimer la demi-vie de l'HbA1c ($t_{1/2, HbA1c}$) en fonction de la durée de vie moyenne des GR ($T_{GR}$). Justifier.
    \item Calculer le temps nécessaire pour que l'HbA1c passe de 9\% à 6,5\% (seuil diagnostique). Donner le résultat en jours.
    \item Pourquoi ce délai clinique est-il important pour le suivi thérapeutique ?
\end{enumerate}
\tcblower
\textbf{Solution détaillée :}
\begin{enumerate}
    \item \textbf{Équation différentielle :}
    Pas de synthèse nouvelle d'HbA1c (glycation non enzymatique stoppée car glycémie normale $\rightarrow$ taux de formation négligeable/constant bas, on modélise $v_s \approx 0$ pour l'excès pathologique).
    Dégradation = disparition des GR porteurs de l'HbA1c. Loi du 1er ordre.
    \[ \frac{dH}{dt} = -k_d \cdot H \quad \text{avec } H(t) \text{ en \%}, \quad k_d = \frac{\ln 2}{t_{1/2}} \text{ (jour}^{-1}\text{)} \]
    \item \textbf{Demi-vie de l'HbA1c :}
    Les GR ont une durée de vie moyenne $T_{GR} = 120$ jours. Leur élimination suit une loi exponentielle (senescence aléatoire). La demi-vie de la population de GR (et donc de l'HbA1c qu'ils portent) est :
    \[ t_{1/2, HbA1c} = T_{GR} \cdot \ln 2 \approx 120 \times 0{,}693 = \textbf{83{,}2 jours} \]
    \emph{Justification :} Pour une population à durée de vie fixe $T$ mais entrée continue, la distribution des âges est exponentielle de paramètre $1/T$. La demi-vie de la population est $T \ln 2$. (Ou plus simplement : élimination exponentielle de la cohorte de GR marqués).
    \item \textbf{Calcul du temps pour atteindre 6,5\% :}
    $H(t) = H_0 \cdot e^{-k_d t} = 9 \cdot 2^{-t / 83{,}2}$.
    On cherche $t$ tel que $H(t) = 6{,}5$.
    \[ 6{,}5 = 9 \cdot 2^{-t / 83{,}2} \iff 2^{-t / 83{,}2} = \frac{6{,}5}{9} \approx 0{,}722 \]
    \[ -\frac{t}{83{,}2} = \log_2(0{,}722) = \frac{\ln 0{,}722}{\ln 2} \approx \frac{-0{,}326}{0{,}693} \approx -0{,}470 \]
    \[ t = 0{,}470 \times 83{,}2 \approx \textbf{39{,}1 jours} \]
    \item \textbf{Intérêt clinique :} Ce délai d'environ \textbf{6 semaines} explique pourquoi on ne dose l'HbA1c que tous les 3 à 6 mois. Un dosage trop précoce (ex. 15 jours) ne refléterait pas l'efficacité du nouveau traitement (inertie du marqueur). Cela évite des changements thérapeutiques intempestifs.
\end{enumerate}
\textbf{Conclusion :} L'HbA1c, par son couplage au turnover érythrocytaire ($t_{1/2} \approx 83$ j), intègre la glycémie sur ~3 mois. La normalisation thérapeutique met ~39 jours pour faire passer l'HbA1c de 9\% à 6,5\%, validant l'intervalle de surveillance trimestriel.
\end{exoresolu}

\begin{attention}[Les 5 erreurs qui coûtent des points au Probatoire]
\begin{enumerate}
    \item \textbf{Confondre « Demi-vie radioactive » et « Demi-vie biologique (turnover) ».} 
    L'isotope (\ce{^{14}C}, \ce{^{3}H}, \ce{^{32}P}) a une demi-vie physique (années pour \ce{^{14}C}, 12 ans pour \ce{^{3}H}, 14 j pour \ce{^{32}P}). La molécule biologiquement marquée a une demi-vie de turnover (heures à jours). \textbf{Erreur :} Dire « la radioactivité diminue car l'isotope se désintègre ». \textbf{Correct :} « La radioactivité spécifique diminue car la molécule marquée est dégradée et remplacée par des molécules non marquées (turnover) ». La décroissance radioactive est négligeable sur l'échelle de l'expérience (heures/jours).
    \item \textbf{Oublier la condition « État stationnaire » pour utiliser $J = k_d M$.} 
    La formule Flux = $k_d \times$ Masse n'est valable que si la masse du pool est constante ($dM/dt = 0$). En phase de croissance, de jeûne, ou pathologie aiguë, $v_s \neq v_d$. \textbf{Réflexe :} Toujours écrire « \emph{Au voisinage de l'état stationnaire...} » ou vérifier que le pool est stable.
    \item \textbf{Interpréter une courbe de Chase plate (ADN) comme « pas de synthèse » au lieu de « pas de dégradation ».} 
    En Pulse-Chase sur l'ADN : Pulse (Phase S) $\rightarrow$ Marquage. Chase $\rightarrow$ Plateau horizontal. \textbf{Erreur :} « L'ADN n'est pas renouvelé car la courbe ne monte plus ». \textbf{Correct :} « L'ADN n'est pas dégradé (turnover nul), la molécule mère est conservée. La dilution par mitose est le seul mécanisme de perte du marquage. » Distinguer \emph{turnover enzymatique} (dégradation/resynthèse) et \emph{dilution réplicative}.
    \item \textbf{Écrire « Le renouvellement permet l'adaptation » sans citer de mécanisme moléculaire précis.} 
    \textbf{Sanction :} 0,5 pt / 2 pts. \textbf{Attendu :} Citer \emph{ubiquitine/protéasome, autophagy, chaperonnes (Hsp70/90), ARNm instables (AU-rich elements), phosphorylation/déphosphorylation, remodelage membranaire (phospholipases), réparation ADN (BER/NER)}. Un exemple précis (ex. « dégradation de l'HIF-1$\alpha$ par le protéasome via VHL en normoxie ») vaut mieux que 3 généralités.
    \item \textbf{Inverser les vitesses : « La dégradation est rapide donc la protéine s'accumule ».} 
    \textbf{Logique :} $dM/dt = v_s - v_d$. Si $v_d$ augmente (turnover accéléré) et $v_s$ constant $\rightarrow$ $M$ \textbf{diminue} (nouvel équilibre plus bas). Si $v_d$ diminue (turnover ralenti, ex. inhibition protéasome) $\rightarrow$ $M$ \textbf{augmente} (accumulation). Au Probatoire, on teste souvent : « Prédire l'évolution du taux de cycline B si on inhibe l'APC/C ». Réponse : Accumulation (car $v_d \to 0$).
\end{enumerate}
\end{attention}

\newpage

\coursec{Exercices}

\coursec{Sensibilisation sur la permanence du renouvellement moléculaire des cellules}

\courssub{Je vérifie mes connaissances}

\begin{exercice}[QCM : les acteurs du renouvellement moléculaire]
Pour chaque question, une seule réponse est exacte. Recopie la lettre de la bonne réponse.

\textbf{1.} Le renouvellement permanent des molécules d'une cellule s'appelle :
\begin{itemize}
\item a) la mitose
\item b) le \cle{turnover moléculaire}
\item c) la méiose
\item d) l'osmose
\end{itemize}

\textbf{2.} La dégradation des protéines cellulaires est assurée principalement par :
\begin{itemize}
\item a) les ribosomes et les mitochondries
\item b) les \cle{lysosomes} et le \cle{protéasome}
\item c) le noyau et la membrane plasmique
\item d) les chloroplastes et le réticulum
\end{itemize}

\textbf{3.} La \cle{demi-vie} d'une molécule est :
\begin{itemize}
\item a) le temps nécessaire à la synthèse de la moitié de la molécule
\item b) le temps au bout duquel la moitié des molécules marquées a été remplacée
\item c) la durée totale de vie de la cellule
\item d) le double du temps de synthèse
\end{itemize}

\textbf{4.} Le marquage isotopique (par exemple au \ce{^{14}C} ou aux acides aminés tritiés) permet de :
\begin{itemize}
\item a) colorer la cellule pour l'observer au microscope
\item b) \textbf{suivre le devenir des molécules et mesurer leur vitesse de renouvellement}
\item c) tuer les cellules cancéreuses
\item d) mesurer la taille des organites
\end{itemize}
\end{exercice}

\begin{exercice}[Vrai ou faux, à justifier]
Pour chaque affirmation, indique si elle est vraie ou fausse, puis justifie ta réponse en une ou deux phrases.
\begin{enumerate}
\item La composition chimique globale d'une cellule reste stable alors que ses molécules sont constamment détruites et resynthétisées.
\item Toutes les molécules d'une cellule se renouvellent à la même vitesse.
\item L'ADN des cellules ne subit jamais d'altérations et n'a donc pas besoin de mécanismes de réparation.
\item Le ralentissement du renouvellement moléculaire avec l'âge contribue au vieillissement cellulaire.
\item Les lipides des membranes cellulaires sont renouvelés en permanence.
\end{enumerate}
\end{exercice}

\begin{exercice}[Définitions essentielles]
Définis en deux ou trois phrases chacun des termes suivants :
\begin{enumerate}
\item \cle{Turnover moléculaire} (renouvellement moléculaire) ;
\item \cle{Demi-vie moléculaire} ;
\item \cle{Protéasome} ;
\item \cle{Autophagie} ;
\item \cle{Vieillissement cellulaire} (sénescence).
\end{enumerate}
\end{exercice}

\begin{exercice}[Calcul direct de demi-vie]
Une protéine du cytoplasme d'une cellule hépatique est marquée au \ce{^{14}C}. À l'instant initial (J0), la radioactivité mesurée est de \num{8000} cpm (coups par minute). On admet que la radioactivité est proportionnelle à la quantité de molécules marquées encore présentes.
\begin{enumerate}
\item Quelle sera la radioactivité mesurée après une demi-vie ? Après deux demi-vies ?
\item On mesure \num{2000} cpm à J12. Détermine la demi-vie de cette protéine.
\item Une autre protéine, du tissu conjonctif, présente une radioactivité de \num{6000} cpm à J0 et de \num{1500} cpm à J20. Calcule sa demi-vie et compare la vitesse de renouvellement des deux protéines.
\end{enumerate}
\end{exercice}

\courssub{Je m'entraîne}

\begin{exercice}[Exploitation d'un tableau de marquage]
Des chercheurs suivent le devenir d'une protéine musculaire marquée par un acide aminé radioactif chez un animal de laboratoire. Les résultats sont consignés dans le tableau ci-dessous.

\begin{center}
\begin{tabularx}{\linewidth}{|l|X|X|X|X|X|X|}
\hline
\rowcolor{popblueL} Temps (jours) & J0 & J2 & J4 & J6 & J8 & J10 \\
\hline
Radioactivité (cpm) & 6400 & 4800 & 3400 & 2400 & 1700 & 1200 \\
\hline
\end{tabularx}
\end{center}

\begin{enumerate}
\item Construis la courbe de décroissance de la radioactivité en fonction du temps (échelle : 1 cm pour 2 jours ; 1 cm pour 1000 cpm).
\item Détermine graphiquement la demi-vie de cette protéine.
\item Cette protéine est-elle renouvelée rapidement ou lentement ? Justifie.
\end{enumerate}
\end{exercice}

\begin{exercice}[Foie ou tendon : à toi d'identifier]
On a mesuré, chez un même animal, la décroissance de la radioactivité de protéines marquées dans deux tissus A et B. Les courbes obtenues sont schématisées ci-dessous : la courbe A a une demi-vie d'environ 2 jours, la courbe B une demi-vie d'environ 100 jours.

\begin{popfigure}
\begin{tikzpicture}
\begin{axis}[
    width=0.85\linewidth, height=6cm,
    xlabel={Temps (jours)}, ylabel={Radioactivité (\% de la valeur initiale)},
    xmin=0, xmax=12, ymin=0, ymax=110,
    legend style={at={(0.97,0.95)}, anchor=north east, font=\small},
    axis lines=left, grid=major, grid style={popline}]
\addplot[popcurve, domain=0:12, samples=100, color=popblue, thick] {100*exp(-0.3466*x)};
\addlegendentry{Tissu A ($t_{1/2} \approx 2$ j)}
\addplot[popcurve, domain=0:12, samples=100, color=poporange, thick] {100*exp(-0.00693*x)};
\addlegendentry{Tissu B ($t_{1/2} \approx 100$ j)}
\end{axis}
\end{tikzpicture}
\legende{Évolution de la radioactivité résiduelle (\% de J0) pour deux tissus aux cinétiques de renouvellement contrastées.}
\end{popfigure}

\begin{enumerate}
\item Sachant que le foie est un organe à forte activité métabolique alors que le tendon est un tissu conjonctif peu renouvelé, attribue à chaque tissu (foie, tendon) la courbe qui lui correspond.
\item Justifie ton attribution en utilisant les demi-vies.
\item Quelle conséquence pratique peux-tu dégager pour la vitesse de guérison d'une hépatite comparée à celle d'une tendinite ?
\end{enumerate}
\end{exercice}

\begin{exercice}[Lecture de document : l'autophagie, Prix Nobel 2016]
Lis le texte suivant, puis réponds aux questions.

\emph{« En 2016, le Prix Nobel de physiologie ou médecine a été attribué au biologiste japonais Yoshinori Ohsumi pour ses découvertes sur les mécanismes de l'autophagie. L'autophagie ("se manger soi-même") est un processus par lequel la cellule englobe ses propres constituants endommagés — organites usés, protéines défectueuses — dans des vésicules appelées autophagosomes, qui fusionnent ensuite avec les lysosomes où le contenu est digéré. Les petites molécules ainsi libérées (acides aminés, acides gras) sont réutilisées par la cellule pour fabriquer de nouvelles molécules. Ce processus, particulièrement actif lors de jeûne ou de stress, est essentiel au maintien de l'intégrité cellulaire ; son dysfonctionnement est impliqué dans le vieillissement et dans certaines maladies. »}

\begin{enumerate}
\item Définis l'autophagie à partir du texte.
\item Cite les deux types de structures cellulaires impliquées et précise le rôle de chacune.
\item Montre, en t'appuyant sur le texte, que l'autophagie participe au renouvellement moléculaire de la cellule.
\item Explique pourquoi un dysfonctionnement de l'autophagie peut favoriser le vieillissement cellulaire.
\end{enumerate}
\end{exercice}

\begin{exercice}[Bilan matière : synthèse contre dégradation]
Une cellule hépatique contient environ \num{8e9} molécules d'une protéine enzymatique dont la demi-vie est de 24 heures. On suppose que la concentration cellulaire de cette protéine reste constante.
\begin{enumerate}
\item Calcule le nombre de molécules de cette protéine dégradées en 24 heures.
\item Combien de molécules la cellule doit-elle synthétiser en 24 heures pour maintenir sa concentration constante ?
\item Formule en une phrase la notion d'\cle{équilibre dynamique} (état stationnaire) illustrée par cet exemple.
\item Cite deux besoins énergétiques ou métaboliques que ce renouvellement permanent impose à l'organisme.
\end{enumerate}
\end{exercice}

\courssub{Je me prépare au Probatoire}

\begin{exercice}[Type Probatoire : le renouvellement d'une protéine musculaire]
\textbf{(D'après un sujet de Probatoire D, adapté)}

Au Centre de recherche médicale de Yaoundé, des chercheurs étudient le renouvellement d'une protéine contractile du muscle squelettique chez le rat. Ils administrent à l'animal un acide aminé marqué au \ce{^{14}C}, qui s'incorpore dans la protéine au moment de sa synthèse. Ils mesurent ensuite la radioactivité de la protéine extraite du muscle à différentes dates. Les résultats sont les suivants :

\begin{center}
\begin{tabularx}{\linewidth}{|l|X|X|X|X|X|X|}
\hline
\rowcolor{popblueL} Temps (jours) & J0 & J3 & J6 & J9 & J12 & J15 \\
\hline
Radioactivité (cpm) & 9600 & 6800 & 4800 & 3400 & 2400 & 1700 \\
\hline
\end{tabularx}
\end{center}

\begin{enumerate}
\item Rappelle ce qu'on appelle renouvellement moléculaire (turnover) et nomme les deux grands types de mécanismes cellulaires qui en sont responsables. \hfill (2 pts)
\item Construis, sur papier millimétré, la courbe de la radioactivité en fonction du temps. Échelle : 1 cm pour 2 jours en abscisse ; 1 cm pour 1000 cpm en ordonnée. \hfill (3 pts)
\item Détermine graphiquement la demi-vie de cette protéine, en faisant apparaître les constructions sur le graphique. \hfill (2 pts)
\item La même expérience réalisée sur une protéine du cristallin de l'œil donne une radioactivité pratiquement constante sur 15 jours. Compare les vitesses de renouvellement des deux protéines et propose une interprétation. \hfill (2 pts)
\item Déduis de cette étude pourquoi une alimentation suffisante en protéines est indispensable tout au long de la vie, même lorsque la masse musculaire ne varie pas. \hfill (1 pt)
\end{enumerate}
\end{exercice}

\begin{exercice}[Type Probatoire : vieillissement et cicatrisation]
\textbf{(Contexte : Centre de santé de Bafoussam)}

Un médecin observe que les plaies cutanées des personnes âgées se referment beaucoup plus lentement que celles des jeunes adultes. Pour comprendre, il fait réaliser des prélèvements de peau et des mesures biochimiques dont voici les résultats :

\begin{center}
\begin{tabularx}{\linewidth}{|X|X|X|}
\hline
\rowcolor{popblueL} Paramètre mesuré & Jeunes adultes (20--30 ans) & Personnes âgées (70--80 ans) \\
\hline
Demi-vie du collagène dermique (jours) & 15 & 40 \\
\hline
Activité du protéasome (unités arbitraires) & 100 & 45 \\
\hline
Quantité de protéines défectueuses accumulées (\%) & 2 & 12 \\
\hline
Durée moyenne de cicatrisation d'une plaie type (jours) & 10 & 28 \\
\hline
\end{tabularx}
\end{center}

\begin{enumerate}
\item Analyse le tableau : dégage deux différences significatives entre les deux groupes. \hfill (2 pts)
\item Rappelle le rôle du protéasome dans la cellule. \hfill (1 pt)
\item Construis, sous forme de chaîne causale ordonnée (au moins quatre étapes), le raisonnement qui relie le ralentissement du renouvellement moléculaire au retard de cicatrisation observé chez la personne âgée. \hfill (3 pts)
\item Explique pourquoi l'accumulation de protéines défectueuses compromet la réparation du derme. \hfill (2 pts)
\item Propose deux conseils (hygiène de vie, alimentation) susceptibles de soutenir le renouvellement moléculaire chez les personnes âgées, en justifiant chacun. \hfill (2 pts)
\end{enumerate}
\end{exercice}

\begin{exercice}[Type Probatoire : question de synthèse et argumentation]
\textbf{Partie A — Question de synthèse.} \hfill (5 pts)

En une dizaine de lignes, explique comment la composition chimique d'une cellule peut rester globalement constante alors que ses molécules (protéines, lipides, ADN) se renouvellent en permanence. Ta réponse devra mobiliser obligatoirement les notions de synthèse, de dégradation, de demi-vie et d'\cle{équilibre dynamique}, et s'appuyer sur un exemple chiffré de ton choix.

\textbf{Partie B — Argumentation guidée.} \hfill (5 pts)

Lors d'une campagne de sensibilisation dans un lycée de Douala, un élève déclare : « Puisque notre corps se renouvelle sans cesse, les maladies liées au vieillissement ne devraient pas exister. »
\begin{enumerate}
\item Montre, en deux arguments précis portant sur l'ADN et sur les systèmes de dégradation, que cette affirmation est erronée.
\item Rédige un court paragraphe (5 à 8 lignes) expliquant à cet élève pourquoi le renouvellement moléculaire, bien que permanent, ne suffit pas à empêcher le vieillissement cellulaire.
\end{enumerate}
\end{exercice}

\newpage

\coursec{Corrigés des exercices}

\begin{corrige}[Exercice 1 — QCM : les acteurs du renouvellement moléculaire]
\textbf{1. Réponse b).} Le renouvellement permanent des molécules constitutives d'une cellule (protéines, lipides, ADN) s'appelle le \cle{turnover moléculaire}. La mitose (a) et la méiose (c) sont des divisions cellulaires, et l'osmose (d) est un déplacement d'eau à travers une membrane : aucun de ces termes ne désigne un renouvellement de molécules.

\textbf{2. Réponse b).} La dégradation des protéines est assurée par les \cle{lysosomes} (organites riches en enzymes digestives, intervenant notamment dans l'autophagie) et par le \cle{protéasome} (complexe enzymatique du cytoplasme qui détruit les protéines marquées par l'ubiquitine). Les ribosomes (a) synthétisent les protéines au lieu de les dégrader.

\textbf{3. Réponse b).} La \cle{demi-vie} moléculaire est le temps au bout duquel la moitié des molécules marquées initialement présentes a été dégradée et remplacée par des molécules neuves. Elle ne concerne ni la synthèse d'une seule molécule (a), ni la durée de vie de la cellule entière (c).

\textbf{4. Réponse b).} Le marquage isotopique (\ce{^{14}C}, acides aminés tritiés \ce{^{3}H}) rend les molécules « traçables » : en mesurant la décroissance de la radioactivité au cours du temps, on suit le devenir des molécules et on mesure leur vitesse de renouvellement. Il ne s'agit ni d'une coloration histologique (a), ni d'une technique thérapeutique (c).
\end{corrige}

\begin{corrige}[Exercice 2 — Vrai ou faux, à justifier]
\begin{enumerate}
\item \textbf{Vrai.} La cellule détruit et resynthétise sans cesse ses molécules à des vitesses équilibrées : c'est un \cle{équilibre dynamique} (état stationnaire). La quantité de chaque molécule reste globalement constante parce que la vitesse de synthèse compense la vitesse de dégradation.
\item \textbf{Faux.} Chaque type de molécule possède sa propre demi-vie : certaines enzymes hépatiques se renouvellent en quelques heures ou jours, alors que le collagène du tendon ou les protéines du cristallin se renouvellent très lentement, voire presque pas.
\item \textbf{Faux.} L'ADN subit en permanence des altérations (rayonnements UV, substances chimiques, erreurs de réplication). Il existe justement des \cle{mécanismes de réparation de l'ADN} ; sans eux, les mutations s'accumuleraient et provoqueraient cancers et mort cellulaire.
\item \textbf{Vrai.} Avec l'âge, l'activité des systèmes de synthèse et de dégradation (protéasome, autophagie) diminue : les molécules défectueuses s'accumulent et les structures ne sont plus remplacées assez vite, ce qui caractérise le \cle{vieillissement cellulaire} (sénescence).
\item \textbf{Vrai.} Les phospholipides et autres lipides des membranes (membrane plasmique, membranes des organites) sont constamment dégradés puis remplacés par des lipides néosynthétisés, ce qui maintient l'intégrité et la fluidité membranaires.
\end{enumerate}
\end{corrige}

\begin{corrige}[Exercice 3 — Définitions essentielles]
\begin{enumerate}
\item \textbf{Turnover moléculaire (renouvellement moléculaire) :} processus permanent par lequel les molécules constitutives de la cellule (protéines, lipides, ADN) sont continuellement dégradées puis remplacées par des molécules neuves synthétisées par la cellule. Il maintient la composition chimique globale de la cellule constante malgré le renouvellement incessant de ses constituants.
\item \textbf{Demi-vie moléculaire :} temps nécessaire pour que la moitié des molécules d'un type donné, présentes à un instant initial, soit dégradée et remplacée. Elle caractérise la vitesse de renouvellement : plus la demi-vie est courte, plus le renouvellement est rapide.
\item \textbf{Protéasome :} grand complexe enzymatique situé dans le cytoplasme et le noyau, qui dégrade les protéines cellulaires marquées au préalable par une petite molécule signal, l'ubiquitine. Il découpe ces protéines en peptides et acides aminés recyclables.
\item \textbf{Autophagie :} mécanisme par lequel la cellule englobe ses propres constituants usés ou endommagés (organites, agrégats de protéines) dans des vésicules appelées autophagosomes, qui fusionnent avec les lysosomes pour digérer leur contenu. Les petites molécules libérées sont réutilisées pour de nouvelles synthèses.
\item \textbf{Vieillissement cellulaire (sénescence) :} ensemble des altérations progressives d'une cellule liées au ralentissement du renouvellement moléculaire et à l'accumulation de molécules défectueuses et de lésions de l'ADN. Il se traduit par une perte d'efficacité des fonctions cellulaires.
\end{enumerate}
\end{corrige}

\begin{corrige}[Exercice 4 — Calcul direct de demi-vie]
\begin{enumerate}
\item Après une demi-vie, la moitié des molécules marquées a disparu :
\[ R_1 = \frac{8000}{2} = \num{4000} \text{ cpm}. \]
Après deux demi-vies, on divise encore par deux :
\[ R_2 = \frac{4000}{2} = \frac{8000}{4} = \num{2000} \text{ cpm}. \]

\item À J12, on mesure \num{2000} cpm, soit $8000 \rightarrow 4000 \rightarrow 2000$ : la radioactivité a été divisée par $4 = 2^2$, donc \textbf{deux demi-vies} se sont écoulées en 12 jours.
\[ t_{1/2} = \frac{12}{2} = 6 \text{ jours}. \]

\item Pour la protéine du tissu conjonctif : $6000 \rightarrow 3000 \rightarrow 1500$. La radioactivité est passée de \num{6000} à \num{1500} cpm, soit une division par $4 = 2^2$ : deux demi-vies en 20 jours.
\[ t_{1/2} = \frac{20}{2} = 10 \text{ jours}. \]
\textbf{Comparaison :} la protéine hépatique (demi-vie de 6 jours) se renouvelle plus rapidement que celle du tissu conjonctif (demi-vie de 10 jours). Le foie, organe à forte activité métabolique, renouvelle ses protéines plus vite que le tissu conjonctif, tissu peu actif.
\end{enumerate}
\end{corrige}

\begin{corrige}[Exercice 5 — Exploitation d'un tableau de marquage]
\begin{enumerate}
\item \textbf{Construction de la courbe.} Avec l'échelle imposée (1 cm pour 2 jours ; 1 cm pour 1000 cpm), on place les points : J0 $\rightarrow$ 6400 ; J2 $\rightarrow$ 4800 ; J4 $\rightarrow$ 3400 ; J6 $\rightarrow$ 2400 ; J8 $\rightarrow$ 1700 ; J10 $\rightarrow$ 1200, puis on trace une courbe régulière décroissante passant par ces points (décroissance de type exponentiel, rapide au début puis de plus en plus lente).

\item \textbf{Détermination graphique de la demi-vie.} La moitié de la valeur initiale est :
\[ \frac{6400}{2} = \num{3200} \text{ cpm}. \]
On trace la droite horizontale d'ordonnée 3200 cpm ; elle coupe la courbe en un point dont on lit l'abscisse. Celle-ci se situe entre J4 (3400 cpm) et J6 (2400 cpm), très près de J4 : on lit $t_{1/2} \approx \num{4,5}$ jours (toute valeur comprise entre 4 et 5 jours, cohérente avec la construction, est acceptable).

\item \textbf{Vitesse de renouvellement.} Une demi-vie d'environ 4 à 5 jours est courte : cette protéine musculaire est donc renouvelée \textbf{rapidement}. En moins de cinq jours, la moitié de son stock moléculaire a déjà été dégradé et remplacé, ce qui traduit l'activité métabolique importante du muscle.
\end{enumerate}
\end{corrige}

\begin{corrige}[Exercice 6 — Foie ou tendon : à toi d'identifier]
\begin{enumerate}
\item \textbf{Attribution :} la courbe A (demi-vie $\approx 2$ jours) correspond au \textbf{foie} ; la courbe B (demi-vie $\approx 100$ jours) correspond au \textbf{tendon}.

\item \textbf{Justification.} Le foie est un organe à forte activité métabolique (détoxication, synthèse de protéines plasmatiques, stockage) : ses protéines sont dégradées et resynthétisées très vite, d'où une demi-vie courte (2 jours) et une décroissance rapide de la radioactivité. Le tendon est un tissu conjonctif riche en collagène, peu vascularisé et à faible activité métabolique : ses protéines sont renouvelées très lentement, d'où une demi-vie longue (100 jours) et une radioactivité qui décroît à peine sur 12 jours.

\item \textbf{Conséquence pratique.} La guérison d'un tissu repose sur le remplacement des molécules et des cellules endommagées. Le foie, à renouvellement rapide, se régénère vite : une hépatite guérit généralement en quelques semaines. Le tendon, à renouvellement très lent, se répare mal : une tendinite exige une convalescence longue, souvent de plusieurs mois, avec repos prolongé.
\end{enumerate}
\end{corrige}

\begin{corrige}[Exercice 7 — Lecture de document : l'autophagie, Prix Nobel 2016]
\begin{enumerate}
\item \textbf{Définition.} D'après le texte, l'autophagie (« se manger soi-même ») est le processus par lequel la cellule englobe ses propres constituants endommagés (organites usés, protéines défectueuses) dans des vésicules, les autophagosomes, afin de les digérer et d'en recycler les constituants.

\item \textbf{Structures impliquées.} Deux types de structures interviennent :
\begin{itemize}
\item les \textbf{autophagosomes} : vésicules qui enferment et transportent les constituants à éliminer ;
\item les \textbf{lysosomes} : organites riches en enzymes digestives, qui fusionnent avec les autophagosomes et digèrent leur contenu en petites molécules (acides aminés, acides gras).
\end{itemize}

\item \textbf{Lien avec le renouvellement moléculaire.} Le texte précise que « les petites molécules ainsi libérées (acides aminés, acides gras) sont réutilisées par la cellule pour fabriquer de nouvelles molécules ». L'autophagie réalise donc l'étape de \textbf{dégradation} du turnover : elle élimine les molécules et organites usés tout en fournissant la matière première des \textbf{synthèses} de remplacement.

\item \textbf{Autophagie et vieillissement.} Si l'autophagie fonctionne mal, les organites usés et les protéines défectueuses ne sont plus éliminés : ils s'accumulent dans le cytoplasme, perturbent le fonctionnement cellulaire et ne sont plus remplacés par des constituants neufs. Ce défaut de « nettoyage » et de recyclage est précisément l'une des causes du vieillissement cellulaire, ce que confirme le texte : « son dysfonctionnement est impliqué dans le vieillissement et dans certaines maladies ».
\end{enumerate}
\end{corrige}

\begin{corrige}[Exercice 8 — Bilan matière : synthèse contre dégradation]
\begin{enumerate}
\item La demi-vie est de 24 h : en 24 heures, la moitié du stock est dégradée.
\[ N_{\text{dégradées}} = \frac{\num{8e9}}{2} = \num{4e9} \text{ molécules}. \]

\item Puisque la concentration reste constante, la synthèse compense exactement la dégradation : la cellule doit synthétiser $\num{4e9}$ molécules en 24 heures.

\item \textbf{Équilibre dynamique (état stationnaire) :} la quantité d'une molécule dans la cellule reste constante parce que sa vitesse de synthèse est égale à sa vitesse de dégradation — le stock est stable alors que les molécules qui le composent changent sans cesse.

\item \textbf{Besoins imposés à l'organisme} (deux réponses parmi) :
\begin{itemize}
\item un \textbf{apport énergétique} permanent (ATP fourni par la respiration cellulaire), car les synthèses de protéines, lipides et ADN consomment beaucoup d'énergie ;
\item un \textbf{apport alimentaire régulier en matières premières} : acides aminés (protéines), acides gras (lipides), nucléotides (ADN), vitamines et minéraux ;
\item le fonctionnement continu des \textbf{systèmes de dégradation} (protéasome, lysosomes), eux-mêmes consommateurs d'énergie.
\end{itemize}
\end{enumerate}
\end{corrige}

\begin{corrige}[Exercice 9 — Type Probatoire : le renouvellement d'une protéine musculaire]
\begin{enumerate}
\item \textbf{Définition (2 pts).} Le renouvellement moléculaire (turnover) est le processus permanent par lequel les molécules de la cellule sont dégradées puis remplacées par des molécules neuves. Il repose sur deux grands types de mécanismes : les \textbf{mécanismes de synthèse} (traduction des protéines sur les ribosomes, synthèse des lipides, réplication et réparation de l'ADN) et les \textbf{mécanismes de dégradation} (protéasome, lysosomes et autophagie).

\item \textbf{Construction de la courbe (3 pts).} Échelle : 1 cm pour 2 jours ; 1 cm pour 1000 cpm. Points à placer : J0 $\rightarrow$ 9600 ; J3 $\rightarrow$ 6800 ; J6 $\rightarrow$ 4800 ; J9 $\rightarrow$ 3400 ; J12 $\rightarrow$ 2400 ; J15 $\rightarrow$ 1700. La courbe est une décroissance régulière (exponentielle), tracée à main levée en passant par tous les points. \emph{Barème indicatif : 1 pt axes orientés et légendés, 1 pt échelles respectées, 1 pt points correctement placés et courbe régulière.}

\item \textbf{Demi-vie (2 pts).} Valeur initiale : 9600 cpm ; moitié : $9600 / 2 = \num{4800}$ cpm. On trace l'horizontale à 4800 cpm : elle coupe la courbe exactement au point d'abscisse J6 (donnée du tableau). D'où :
\[ t_{1/2} = 6 \text{ jours}. \]
\emph{Point de vigilance : la construction graphique (traits de lecture en pointillés) doit apparaître sur la copie ; une réponse sans construction ne vaut pas la totalité des points.}

\item \textbf{Comparaison avec le cristallin (2 pts).} La radioactivité constante sur 15 jours signifie que les molécules marquées ne sont presque pas remplacées : la protéine du cristallin a une demi-vie très longue et se renouvelle extrêmement lentement, contrairement à la protéine musculaire ($t_{1/2} = 6$ j). \textbf{Interprétation :} la vitesse de renouvellement dépend du tissu et de sa fonction ; le cristallin, structure stable qui doit rester transparente toute la vie, renouvelle très peu ses protéines, alors que le muscle, très actif métaboliquement, les renouvelle vite.

\item \textbf{Conséquence alimentaire (1 pt).} Même si la masse musculaire ne varie pas, ses protéines sont détruites et resynthétisées en permanence (ici, la moitié du stock en 6 jours). L'organisme doit donc recevoir chaque jour des protéines alimentaires pour fournir les acides aminés indispensables à ces synthèses, à tout âge de la vie.
\end{enumerate}
\textbf{Points de vigilance (ensemble de l'exercice) :} définir le turnover par les deux termes « dégradation » et « synthèse » ; citer nommément au moins un organite ou complexe de dégradation ; justifier toute lecture graphique par un calcul de la moitié de la valeur initiale ; ne pas confondre « radioactivité constante » avec « absence de molécules ».
\end{corrige}

\begin{corrige}[Exercice 10 — Type Probatoire : vieillissement et cicatrisation]
\begin{enumerate}
\item \textbf{Analyse du tableau (2 pts).} Deux différences significatives (parmi les quatre possibles) :
\begin{itemize}
\item la demi-vie du collagène passe de 15 à 40 jours : le renouvellement du collagène est presque trois fois plus lent chez la personne âgée ;
\item l'activité du protéasome chute de 100 à 45 unités : le système de dégradation des protéines est plus de deux fois moins actif ;
\item (conséquences) les protéines défectueuses accumulées passent de 2 \% à 12 \% et la durée de cicatrisation de 10 à 28 jours.
\end{itemize}

\item \textbf{Rôle du protéasome (1 pt).} Le protéasome est le complexe enzymatique qui dégrade les protéines cellulaires usées ou défectueuses, préalablement marquées par l'ubiquitine, en les découpant en petits peptides et acides aminés recyclables.

\item \textbf{Chaîne causale (3 pts).}
Ralentissement du renouvellement moléculaire avec l'âge $\rightarrow$ baisse de l'activité du protéasome (et des synthèses) $\rightarrow$ accumulation de protéines défectueuses et remplacement trop lent du collagène $\rightarrow$ dégradation de la qualité et de la cohésion du derme $\rightarrow$ retard de la réparation tissulaire, donc cicatrisation plus lente (28 jours contre 10).
\emph{Barème indicatif : 1 pt pour l'enchaînement logique, 1 pt pour l'utilisation d'au moins deux données chiffrées du tableau, 1 pt pour la conclusion reliant la chaîne au retard de cicatrisation.}

\item \textbf{Accumulation de protéines défectueuses (2 pts).} Les protéines défectueuses (collagène altéré, enzymes inefficaces) ne remplissent plus leur fonction : elles encombrent les cellules, perturbent leur métabolisme et fragilisent la trame du derme. Or la cicatrisation exige la synthèse rapide de collagène neuf et la multiplication des cellules de réparation ; des cellules encombrées de protéines anormales travaillent mal, ce qui compromet et ralentit la reconstruction du derme.

\item \textbf{Conseils (2 pts, 1 pt par conseil justifié).} Par exemple :
\begin{itemize}
\item une \textbf{alimentation suffisamment riche en protéines} (poisson, œufs, haricots niébé, arachides) : elle fournit les acides aminés nécessaires à la synthèse du collagène et des protéines de réparation ;
\item une \textbf{activité physique régulière et modérée} : elle stimule le métabolisme, la vascularisation de la peau et le renouvellement des protéines musculaires et cutanées ;
\item (accepté aussi) une bonne hydratation et une alimentation riche en vitamines (fruits et légumes locaux), cofacteurs des enzymes de synthèse et de réparation.
\end{itemize}
\end{enumerate}
\textbf{Points de vigilance :} toute affirmation doit s'appuyer sur une donnée chiffrée du tableau ; la chaîne causale doit être ordonnée (flèches logiques) et non une simple liste ; les conseils doivent être reliés explicitement au renouvellement moléculaire.
\end{corrige}

\begin{corrige}[Exercice 11 — Type Probatoire : question de synthèse et argumentation]
\textbf{Partie A — Question de synthèse (5 pts).}

\emph{Éléments de réponse attendus (une dizaine de lignes) :}
La composition chimique d'une cellule reste globalement constante grâce à un \cle{équilibre dynamique} : chaque molécule détruite est aussitôt remplacée par une molécule neuve. Les \textbf{mécanismes de synthèse} (traduction des protéines sur les ribosomes, synthèse des lipides, réplication et réparation de l'ADN) compensent exactement les \textbf{mécanismes de dégradation} (protéasome, lysosomes, autophagie). Chaque type de molécule possède sa propre \textbf{demi-vie}, temps au bout duquel la moitié du stock marqué a été remplacée. Par exemple, une protéine hépatique présente en $\num{8e9}$ exemplaires avec une demi-vie de 24 h voit $\num{4e9}$ de ses molécules dégradées chaque jour, mais la cellule en synthétise simultanément $\num{4e9}$ de neuves : le stock reste constant alors que les molécules elles-mêmes changent sans cesse. Ainsi, la cellule est comme un fleuve : l'eau passe, mais le niveau reste le même — c'est l'état stationnaire du turnover moléculaire.

\emph{Barème indicatif : 1 pt notion de synthèse, 1 pt notion de dégradation (avec acteurs nommés), 1 pt définition et usage correct de la demi-vie, 1 pt équilibre dynamique explicité, 1 pt exemple chiffré cohérent et exact.}

\textbf{Partie B — Argumentation guidée (5 pts).}

\begin{enumerate}
\item \textbf{Deux arguments réfutant l'affirmation (2 pts, 1 pt chacun) :}
\begin{itemize}
\item \textbf{Argument sur l'ADN :} l'ADN subit en permanence des altérations (UV, substances toxiques, erreurs de réplication). Les mécanismes de réparation existent mais deviennent moins efficaces avec l'âge : des mutations s'accumulent et peuvent conduire à des cancers ou à la mort cellulaire, malgré le renouvellement permanent.
\item \textbf{Argument sur les systèmes de dégradation :} avec l'âge, l'activité du protéasome et de l'autophagie diminue ; les protéines défectueuses et les organites usés ne sont plus éliminés assez vite et s'accumulent, ce qui altère le fonctionnement cellulaire : le renouvellement ralentit et devient insuffisant.
\end{itemize}

\item \textbf{Paragraphe explicatif (3 pts, 5 à 8 lignes) :}
« Le renouvellement moléculaire est bien permanent, mais il n'est ni parfait ni inusable. Avec l'âge, les systèmes de dégradation (protéasome, autophagie) perdent de leur efficacité : les protéines défectueuses s'accumulent au lieu d'être recyclées. Les mécanismes de réparation de l'ADN, eux aussi, deviennent défaillants, et les mutations s'additionnent au fil des années. Enfin, les synthèses elles-mêmes ralentissent : le collagène de la peau, par exemple, se renouvelle presque trois fois plus lentement chez la personne âgée. Le renouvellement continue donc, mais trop lentement et avec trop d'erreurs pour compenser toutes les altérations : c'est précisément ce déséquilibre progressif qui constitue le vieillissement cellulaire et qui explique l'apparition des maladies liées à l'âge. »
\end{enumerate}
\emph{Barème indicatif Partie B : 2 pts pour les deux arguments précis (ADN + systèmes de dégradation), 2 pts pour la cohérence biologique du paragraphe, 1 pt pour la clarté et l'adaptation au destinataire (expliquer à un élève).}
\textbf{Point de vigilance :} ne pas se contenter d'affirmer que l'élève « a tort » ; chaque argument doit nommer un mécanisme précis (réparation de l'ADN, protéasome, autophagie) et sa conséquence.
\end{corrige}

\newpage

\coursec{Fiche bilan}

\begin{popfigure}
\begin{center}\tcbox[colback=popblue,colframe=popblue,coltext=white,arc=9pt,boxsep=4pt,left=6pt,right=6pt]{\bfseries\small Renouvellement moléculaire cellulaire}\end{center}
\vspace{-2pt}
\begin{tcbraster}[raster columns=3,raster equal height=rows,raster column skip=6pt,raster row skip=6pt,enhanced,blankest]
\begin{tcolorbox}[enhanced,colback=white,colframe=poporange,boxrule=0.9pt,arc=3pt,title={Évidence expérimentale},fonttitle=\bfseries\scriptsize,coltitle=white,colbacktitle=poporange,left=4pt,right=4pt,top=2pt,bottom=2pt,boxsep=1pt,toptitle=1.5pt,bottomtitle=1.5pt,halign=left]
\begin{itemize}[nosep,leftmargin=*,itemsep=1pt,font=\scriptsize]
\item Marquage isotopique
\item $^{15}$N, $^{14}$C, leucine tritiée
\end{itemize}
\end{tcolorbox}
\begin{tcolorbox}[enhanced,colback=white,colframe=popgreen,boxrule=0.9pt,arc=3pt,title={Synthèse},fonttitle=\bfseries\scriptsize,coltitle=white,colbacktitle=popgreen,left=4pt,right=4pt,top=2pt,bottom=2pt,boxsep=1pt,toptitle=1.5pt,bottomtitle=1.5pt,halign=left]
\begin{itemize}[nosep,leftmargin=*,itemsep=1pt,font=\scriptsize]
\item Protéines : transcription + traduction
\item ADN : réplication
\item Lipides : réticulum
\end{itemize}
\end{tcolorbox}
\begin{tcolorbox}[enhanced,colback=white,colframe=poppurple,boxrule=0.9pt,arc=3pt,title={Dégradation},fonttitle=\bfseries\scriptsize,coltitle=white,colbacktitle=poppurple,left=4pt,right=4pt,top=2pt,bottom=2pt,boxsep=1pt,toptitle=1.5pt,bottomtitle=1.5pt,halign=left]
\begin{itemize}[nosep,leftmargin=*,itemsep=1pt,font=\scriptsize]
\item Protéasome
\item Lysosome
\item Enzymes hydrolytiques
\end{itemize}
\end{tcolorbox}
\begin{tcolorbox}[enhanced,colback=white,colframe=poppink,boxrule=0.9pt,arc=3pt,title={Turnover},fonttitle=\bfseries\scriptsize,coltitle=white,colbacktitle=poppink,left=4pt,right=4pt,top=2pt,bottom=2pt,boxsep=1pt,toptitle=1.5pt,bottomtitle=1.5pt,halign=left]
\begin{itemize}[nosep,leftmargin=*,itemsep=1pt,font=\scriptsize]
\item Demi-vie variable
\item Équilibre synthèse/dégradation
\end{itemize}
\end{tcolorbox}
\begin{tcolorbox}[enhanced,colback=white,colframe=popgold,boxrule=0.9pt,arc=3pt,title={Intégrité cellulaire},fonttitle=\bfseries\scriptsize,coltitle=white,colbacktitle=popgold,left=4pt,right=4pt,top=2pt,bottom=2pt,boxsep=1pt,toptitle=1.5pt,bottomtitle=1.5pt,halign=left]
\begin{itemize}[nosep,leftmargin=*,itemsep=1pt,font=\scriptsize]
\item Renouvellement des structures
\item Élimination des molécules altérées
\end{itemize}
\end{tcolorbox}
\begin{tcolorbox}[enhanced,colback=white,colframe=popdark,boxrule=0.9pt,arc=3pt,title={Vieillissement},fonttitle=\bfseries\scriptsize,coltitle=white,colbacktitle=popdark,left=4pt,right=4pt,top=2pt,bottom=2pt,boxsep=1pt,toptitle=1.5pt,bottomtitle=1.5pt,halign=left]
\begin{itemize}[nosep,leftmargin=*,itemsep=1pt,font=\scriptsize]
\item Ralentissement du turnover
\item Accumulation de déchets
\item Perte de fonctions
\end{itemize}
\end{tcolorbox}
\end{tcbraster}

\legende{Carte mentale de la leçon : le renouvellement permanent des molécules (protéines, lipides, ADN) résulte d'un équilibre entre synthèse et dégradation ; son ralentissement conduit au vieillissement cellulaire.}
\end{popfigure}

\begin{bilan}
\textbf{Définitions clés}
\begin{itemize}
  \item \cle{Renouvellement moléculaire} (ou \emph{turnover}) : remplacement permanent des molécules constitutives de la cellule (protéines, lipides, ADN) par des molécules néoformées, sans changement apparent de la structure.
  \item \cle{Synthèse moléculaire} : fabrication de nouvelles molécules à partir de précurseurs (acides aminés, nucléotides, acides gras) ; elle consomme de l'énergie (ATP).
  \item \cle{Dégradation moléculaire} : destruction des molécules usées ou altérées par des enzymes (protéases, lipases, nucléases) au sein du \cle{protéasome} et du \cle{lysosome}.
  \item \cle{Demi-vie} d'une molécule : temps nécessaire pour que la moitié d'une population de molécules soit renouvelée.
  \item \cle{Vieillissement cellulaire} : ensemble des altérations liées au ralentissement du renouvellement moléculaire (accumulation de molécules défectueuses, perte d'activité).
\end{itemize}

\textbf{Repères quantitatifs et mécanismes à connaître}
\begin{center}
\begin{tabularx}{\linewidth}{|l|X|}
\hline
\rowcolor{popblueL} \textbf{Notion} & \textbf{À retenir} \\
\hline
Preuve expérimentale & Marquage isotopique ($^{15}$N, $^{14}$C, acides aminés radioactifs) : les molécules marquées disparaissent et sont remplacées, même dans une cellule qui ne se divise pas. \\
\hline
Synthèse des protéines & ADN $\rightarrow$ ARNm (transcription, noyau) $\rightarrow$ protéine (traduction, ribosomes). \\
\hline
Dégradation & Protéasome (protéines marquées par l'ubiquitine) ; lysosome (hydrolyses en milieu acide). \\
\hline
Bilan énergétique & Synthèse : consomme de l'ATP ; dégradation : recycle les constituants (acides aminés, nucléotides). \\
\hline
Équilibre & Masse cellulaire constante $\Rightarrow$ vitesse de synthèse $=$ vitesse de dégradation. \\
\hline
\end{tabularx}
\end{center}

\textbf{Méthodes express}
\begin{itemize}
  \item \emph{Exploiter un marquage isotopique} : repérer la molécule marquée, suivre la radioactivité au cours du temps, conclure sur le renouvellement (disparition du marqueur $=$ dégradation + remplacement).
  \item \emph{Relier turnover et santé} : un renouvellement efficace élimine les protéines altérées et maintient l'intégrité des membranes et de l'ADN.
  \item \emph{Argumenter sur le vieillissement} : ralentissement du turnover $\rightarrow$ accumulation de molécules défectueuses et de pigments (lipofuscine) $\rightarrow$ dysfonctionnements cellulaires.
\end{itemize}
\end{bilan}

\begin{aretenir}[Checklist avant le Probatoire]
\begin{itemize}
  \item[$\square$] Je sais définir le renouvellement moléculaire (turnover) en une phrase.
  \item[$\square$] Je cite les trois grandes familles de molécules renouvelées : protéines, lipides, ADN.
  \item[$\square$] Je sais décrire le principe du marquage isotopique et ce qu'il démontre.
  \item[$\square$] Je connais les deux étapes de la synthèse des protéines (transcription, traduction) et leur lieu.
  \item[$\square$] Je cite les deux structures de dégradation : protéasome et lysosome.
  \item[$\square$] Je sais que la synthèse consomme de l'ATP et que la dégradation recycle les constituants.
  \item[$\square$] Je sais définir la demi-vie d'une molécule.
  \item[$\square$] Je sais expliquer pourquoi la masse d'une cellule reste stable malgré le turnover.
  \item[$\square$] Je relie le renouvellement moléculaire au maintien de l'intégrité cellulaire.
  \item[$\square$] Je sais expliquer le lien entre ralentissement du turnover et vieillissement cellulaire.
  \item[$\square$] Je sais exploiter une courbe de décroissance de radioactivité pour conclure sur un renouvellement.
  \item[$\square$] Je sais argumenter sur l'importance du renouvellement moléculaire pour la santé (ex. renouvellement des globules rouges, cicatrisation).
\end{itemize}
\end{aretenir}

\findecours

\end{document}
